% Review source — Unpacking Alpöge's Proof of the Riemann Hypothesis, v1.1-review
% The .tex file is the publication source; the PDF is its compiled output.
% All diagrams, including the merged roadmap, are embedded as editable TikZ.
% Build: latexmk -pdf -interaction=nonstopmode -halt-on-error Unpacking_Alpoges_Proof_Of_the_Riemann_Hypothesis_v1.1-review.tex.
% All externally included TikZ figures are in figures/.
\documentclass[11pt,openany,oneside]{scrbook}
\usepackage[paperwidth=6.6in,paperheight=9.3in,top=.87in,bottom=.86in,inner=.85in,outer=.77in,headsep=.28in]{geometry}
\usepackage[T1]{fontenc}
\usepackage[utf8]{inputenc}
\usepackage{amsmath,amsthm,mathtools}
\usepackage{newtxtext,newtxmath}
\usepackage{microtype}
\usepackage{tikz}
\usepackage{booktabs,tabularx,enumitem}
\usepackage{fancyhdr}
\usepackage{float}
\usepackage[backend=biber,style=authoryear,sorting=nyt,maxcitenames=2,maxbibnames=8,doi=true,url=true]{biblatex}
\addbibresource{SCHOLARLY_BIBLIOGRAPHY.bib}
\usepackage{hyperref}
\usepackage{xcolor}
\usepackage{listings}
\usepackage{needspace}
\usepackage{caption}
\usetikzlibrary{arrows.meta,positioning,calc,decorations.pathreplacing}
\definecolor{ink}{RGB}{25,35,52}
\definecolor{accent}{RGB}{32,77,111}
\definecolor{muted}{RGB}{95,102,110}
\definecolor{sourcegray}{RGB}{246,247,249}
\hypersetup{colorlinks=true,linkcolor=accent,urlcolor=accent,citecolor=accent,pdftitle={Unpacking Alpoge's Proof of the Riemann Hypothesis},pdfauthor={Mathematical expository draft},pdfsubject={Version v1.1-review; mathematical review ongoing},pdfkeywords={Manuscript version v1.1-review; proposed proof; mathematical exposition}}
\KOMAoptions{headings=normal,parskip=false}
\setlength{\parindent}{1.13em}
\setlength{\parskip}{.12em}
\linespread{1.105}
\addtolength{\emergencystretch}{2em}
\setlist[itemize]{leftmargin=1.7em,topsep=3pt,itemsep=1.8pt}
\setlist[enumerate]{leftmargin=1.8em,topsep=3pt,itemsep=2pt}
\pagestyle{fancy}
\fancyhf{}
\fancyhead[L]{\footnotesize\scshape Unpacking Alp\"oge's Proof}
\fancyhead[R]{\footnotesize\scshape Draft for Review}
\fancyfoot[C]{\thepage}
\fancyfoot[R]{\footnotesize v\manuscriptversion}
\renewcommand{\headrulewidth}{0.25pt}
\fancypagestyle{plain}{\fancyhf{}\fancyfoot[C]{\thepage}\fancyfoot[R]{\footnotesize v\manuscriptversion}\renewcommand{\headrulewidth}{0pt}}
\newtheorem{theorem}{Theorem}[chapter]
\newtheorem{proposition}[theorem]{Proposition}
\newtheorem{lemma}[theorem]{Lemma}
\newtheorem{corollary}[theorem]{Corollary}
\theoremstyle{definition}
\newtheorem{definition}[theorem]{Definition}
\newtheorem{example}[theorem]{Example}
\theoremstyle{remark}
\newtheorem{remark}[theorem]{Remark}
\newcommand{\CC}{\mathbb C}
\newcommand{\RR}{\mathbb R}
\newcommand{\QQ}{\mathbb Q}
\newcommand{\ZZ}{\mathbb Z}
\newcommand{\NN}{\mathbb N}
\newcommand{\F}{\mathbb F}
\newcommand{\Rea}{\operatorname{Re}}
\newcommand{\ord}{\operatorname{ord}}
\newcommand{\Gal}{\operatorname{Gal}}
\newcommand{\Shasym}{\operatorname{Sha}}
\newcommand{\sig}{\sigma}
\newcommand{\eps}{\varepsilon}
\newcommand{\manuscriptversion}{1.1-review}
\newcommand{\baseurl}{https://github.com/CocoSgt/riemann-hypothesis-lean/blob/7cbf7a0c32fca1badf05931879a37c3f8a2b03b4/}
\newcommand{\sourcefile}[2]{\href{\baseurl#1}{\texttt{#2}}}
\lstset{basicstyle=\ttfamily\footnotesize,backgroundcolor=\color{sourcegray},frame=single,rulecolor=\color{black!10},breaklines=true,columns=fullflexible,keepspaces=true,xleftmargin=5pt,xrightmargin=5pt,aboveskip=7pt,belowskip=8pt}
\captionsetup{font=small}
\title{\vspace{-1.3cm}{\LARGE\bfseries Unpacking Alp\"oge's Proof\\Of the Riemann Hypothesis}\\[.7em]{\Large A Guided Reconstruction of the Lean Argument}\\[1.2em]{\normalsize Draft for mathematical review and feedback}}
\author{\parbox{.78\textwidth}{\centering\small\itshape With boundless curiosity, I (Brian Kramer) worked with ChatGPT to reverse-engineer Alp\"oge\textquotesingle s formalized proof of the Riemann Hypothesis and make its mathematics more accessible.}\\[2.1em]\small Feedback: \texttt{zetacontinuation@gmail.com}}
\date{21 November 2026\\[.5em]\small Version v\manuscriptversion}
\begin{document}
\frontmatter
\maketitle
\thispagestyle{empty}
\vspace*{1.2cm}
\noindent\textbf{DRAFT --- COMMENTS AND CRITICAL FEEDBACK WELCOME.} This manuscript is a working exposition of Alp\"oge's announced proof, not an independent certification of the proof or of every explanation here. Readers are invited to identify inaccuracies, missing hypotheses, and gaps in the reconstruction. Please send corrections and feedback to \href{mailto:zetacontinuation@gmail.com}{\texttt{zetacontinuation@gmail.com}}.

\medskip
\noindent\textbf{About the mathematical status.} In November 2026 Levent Alp\"oge announced a proof of the Riemann Hypothesis and its Dirichlet $L$-function generalization, developed with Claude and formalized in a very large Lean archive. Public independent checking has advanced substantially: Andrew Sutherland's MIT report of 21 November 2026 records compilation of all 175,187 modules, Lean kernel replay, acceptance by independently written kernels, and comparison of the five final statements with an independent challenge. Some long-running export and comparator checks remained pending at that update \parencite{Sutherland2026Verification}. This book did \emph{not} rerun those checks, nor does acceptance by a formal kernel certify the accuracy of every informal explanation given here. We distinguish short arguments proved in these pages, deep results cited from the formal source, and the interpretation that connects them. The remaining mathematical-review tasks are recorded in the appendices.

\medskip
\noindent\textit{Manuscript version:} \href{https://doi.org/10.5281/zenodo.24433694}{\textbf{v\manuscriptversion}}, with an additional example in Section~\ref{sec:imprimitive-euler}.\\
\textit{Lean source revision:} \texttt{7cbf7a0c32fca1badf05931879a37c3f8a2b03b4}.\\
\textit{Announcement:} \url{https://alpo.ge/remark.pdf}.\\
\textit{Source mirror:} \url{https://github.com/CocoSgt/riemann-hypothesis-lean}.

\medskip
\noindent\textit{Related exposition.} The 55-page \emph{Dirichlet L-functions and the Critical Line} informed comparisons of the analytic, Fourier, and lattice arguments and led to fuller treatment of the analytic pole comparison. A shorter exposition under the same title was circulated at \url{https://ulam.ai/research/rh-short.pdf}. These are secondary expositions, not independent verification of the Lean development.
\clearpage

\chapter*{Preface}
\addcontentsline{toc}{chapter}{Preface}
In 1859 Bernhard Riemann published a short paper about the distribution of prime numbers. One of its assertions, concerning the complex zeros of a function now bearing his name, became the Riemann Hypothesis. For more than a century and a half it has remained a reliable source of difficult, consequential mathematics: techniques developed to approach it have reshaped number theory, even though the original question has remained out of reach.

In November 2026 Levent Alp\"oge announced that he and Claude had developed a proof of the Riemann Hypothesis and formalized it in Lean. The announcement came with an archive containing millions of lines of code, a concrete formal development that mathematicians could examine. Its scale creates a different difficulty for the reader: the architecture of the argument and the reasons its constructions belong together are hard to discern from the archive alone. Formal verification and human mathematical review remain distinct tasks; both deserve careful reporting.

The archive allows us to identify the analytic estimate that would exclude zeros away from the critical line and the arithmetic constructions intended to yield it. We can examine those connections without claiming to have independently reconstructed every deep theorem in the formal development.

Alp\"oge's four-page \emph{Remark} offers a glimpse of the work behind the formal archive. Its timeline begins on October 6, 2026, with a Stepanov determinant, and ends 41 days later, on November 16, with the entry ``GRH.'' Between those dates lie changes of direction and tools drawn from widely separated areas of mathematics \parencite{Alpoege2026Remark}. The archive's file count cannot tell us how many distinct ideas were needed to make those tools work together.

We begin with the analytic argument connecting a bound on weighted prime sums to the location of zeros. It tells us precisely what the arithmetic constructions must supply: an upper bound close to the square-root scale, on suitably selected intervals, for a positive combination of character sums. With that target in view, we can examine how the successive constructions produce and preserve the needed estimate, rather than tracing the Lean archive declaration by declaration.

Once that analytic target is fixed, two features of the proposed method become central. First, a chain of theta lifts is used to reproduce an advantage that Stepanov's method obtains from repeated vanishing over finite fields. Second, the resulting arithmetic information must be measured: congruences and divisibility at finite primes must be compared with norms at the infinite places. The comparison is delicate enough to require auxiliary number fields, Galois representations, cohomological bounds, and carefully chosen character averages. It is this chain of mathematical reasons, rather than the sequence of Lean declarations, that we shall follow.

Several elementary deductions can be proved directly here, including the finite Euler-factor reduction, the ramp transform, and a positivity calculation for pole residues. Other steps depend on substantial construction theorems whose statements can be identified in the source but whose proofs have not been independently reconstructed for this exposition. We indicate the relevant hypotheses and source boundaries where those results are used, so that the reader can distinguish the calculations supplied here from the deeper imported assertions.

\bigskip
\noindent\textit{Sources and scope.} References to the Lean archive use a fixed Git revision, so that definitions cannot shift during the exposition. Quotations from Alp\"oge's \emph{Remark} are brief, attributed, and used only where they bear on the mathematical construction under discussion. Diagrams show specific mathematical operations or quantifier restrictions; they are not substitutes for theorems.

\chapter*{Outline of the argument}
\addcontentsline{toc}{chapter}{Outline of the argument}

The analytic goal is nonvanishing of the Riemann zeta function to the right of the critical line, apart from its pole at $s=1$. Once that has been established, the functional equation reflects a hypothetical zero to the left of the line into the forbidden right half-plane. The arithmetic and analytic constructions therefore concentrate on a one-sided zero-free estimate.

\medskip
\begin{center}
\begin{tikzpicture}[>=Latex, every node/.style={font=\footnotesize,align=center}, stage/.style={draw=black!32,rounded corners=2.5pt,fill=blue!3,minimum width=.89\textwidth,text width=.83\textwidth,inner ysep=6pt,inner xsep=8pt}]
\node[stage] (a) {\textbf{Build arithmetic information.}\\Theta lifts, auxiliary fields, and divisibility control many related quantities.};
\node[stage,below=6mm of a] (b) {\textbf{Convert it into a usable estimate.}\\Lattice-size bounds and character averaging isolate a positive family of prime sums.};
\node[stage,below=6mm of b] (c) {\textbf{Apply a localized analytic criterion.}\\A sufficiently strong upper bound on arbitrarily late intervals rules out zeros to the right of $1/2$.};
\node[stage,below=6mm of c,fill=green!6] (d) {\textbf{Finish by symmetry.}\\The functional equation transfers the zero-free conclusion to the left of the critical line.};
\foreach \u/\v in {a/b,b/c,c/d} \draw[->,accent,line width=.85pt] (\u)--(\v);
\end{tikzpicture}
\end{center}

\noindent\textbf{Analytic starting point.} Chapter~\ref{ch:halfplane} explains the zero-free goal and the final reflection. Chapter~\ref{ch:statistic} introduces the weighted prime sum and the analytic criterion that makes an estimate on that sum relevant to zeros. Those two chapters specify what the remaining arithmetic machinery is being asked to deliver.

\medskip
\noindent The diagram records the principal implications. The arithmetic inputs, their quantitative dependencies, and the corresponding cohomological windows are displayed more fully in the unnumbered chapter \emph{Arithmetic dependencies of the analytic criterion} after Chapter~\ref{ch:statistic}. Each arrow depends on hypotheses proved in later chapters or cited from the pinned Lean source; the diagram itself supplies none of those proofs.

\medskip
\noindent\textbf{Reading by subject.} The analytic argument can be read in Chapters~\ref{ch:halfplane} and~\ref{ch:statistic}, then in the final deduction of Chapter~\ref{ch:assembly}. For multiplicity and theta correspondence, begin with the first two chapters of Part~II. The arithmetic material in Parts~III and~IV is more interdependent; the chapter introductions identify the quantities imported from earlier arguments. The appendices provide the exact source boundaries for readers who want to follow a technical step further.

\tableofcontents
\mainmatter
\part{Zeros and the analytic criterion}

\chapter{The zero-free half-plane}\label{ch:halfplane}
\section{The logarithmic derivative and prime powers}
Two familiar descriptions of the Riemann zeta function drive the argument. For a complex number $s$ with $\Rea(s)>1$,
\begin{equation}\label{eq:euler}
\zeta(s)=\sum_{n\ge1}\frac1{n^s}=\prod_p(1-p^{-s})^{-1}.
\end{equation}
The series adds the terms $n^{-s}$ over all positive integers; the product organizes the same arithmetic information prime by prime. Both descriptions converge absolutely when $\Rea(s)>1$ \parencite{Davenport2000,IwaniecKowalski2004}. Their coexistence is the reason a statement about the location of zeros can constrain the distribution of primes.

The Euler product lays out the primes one by one. It also gives a zero-free region, but only where absolute convergence makes the argument valid—well to the right of the critical line. We need its arithmetic information where the product can no longer be used directly.

To see the link in its most economical form, differentiate the logarithm of the Euler product. Absolute convergence permits termwise differentiation, giving
\begin{align}\label{eq:logderiv}
-\frac{\zeta'(s)}{\zeta(s)}
 &=\sum_p\frac{(\log p)p^{-s}}{1-p^{-s}}
 =\sum_p\sum_{k\ge1}\frac{\log p}{p^{ks}}
 =\sum_{n\ge1}\frac{\Lambda(n)}{n^s},
\end{align}
where $\Lambda(n)=\log p$ if $n=p^k$ for some prime $p$ and integer $k\ge1$, and $\Lambda(n)=0$ otherwise. The function $\Lambda$ is called the von Mangoldt function. It does not merely indicate primes; it weights prime powers by the logarithm of their underlying prime.

Equation~\eqref{eq:logderiv} has put zeros and prime powers into the same calculation. On the left, a zero of $\zeta$ appears as a singularity of its logarithmic derivative. On the right, the same analytic function is a generating series for the prime powers. If we can continue the left side far enough while controlling its singularities, we obtain information about the partial sums of $\Lambda$.

Write
\begin{equation}\label{eq:psi}
\psi(x)=\sum_{n\le x}\Lambda(n).
\end{equation}
The prime number theorem is equivalent to $\psi(x)\sim x$. In other words, the logarithmically weighted prime powers accumulate at the expected average rate. The point of the Riemann Hypothesis is that it predicts an essentially square-root-sized scale for the fluctuation of this weighted count, up to logarithmic factors.

The main term $x$ comes from the pole of $\zeta$ at $s=1$; the fluctuations are governed by its zeros. Away from prime-power discontinuities and with the necessary limiting procedure, the classical explicit formula reads
\begin{equation}\label{eq:explicit}
\psi(x)=x-\sum_{\rho}\frac{x^\rho}{\rho}-\log(2\pi)-\frac12\log(1-x^{-2}),\qquad x>1\text{ not a prime power},
\end{equation}
where $\rho$ ranges over nontrivial zeros with multiplicity. The sum is not an absolutely convergent series that may be rearranged at will; it must be understood through the standard symmetric truncation in height. At a prime power $x=p^r$, this truncation instead gives the midpoint value $\psi(x)-\tfrac12\log p$. A zero with real part $\Rea(\rho)=\beta$ contributes a term of size governed by $x^\beta/|\rho|$ before cancellation among the zero terms is considered; the location of the zeros is therefore visible in the explicit formula.

For $\rho=\beta+i\gamma$,
\[
\left|x^\rho\right|=x^\beta.
\]
The real part $\beta$ is therefore the exponent controlling the size of an individual oscillatory contribution. A zero with $\beta=0.7$ carries a power $x^{0.7}$; a zero on the line $\beta=1/2$ carries the square-root power $x^{1/2}$. These are not assertions about the total sum---different zero contributions can cancel---but they explain why moving the allowable real parts of the zeros changes the error-term problem so dramatically.

\begin{remark}
A bound for $\psi(x)$ strong enough to imply RH cannot be read off from individual terms of~\eqref{eq:explicit}. It also requires estimates for the density of zeros, truncation errors, and cancellation among their contributions.
\end{remark}

\section{Absolute convergence of the Euler product}
Within $\Rea(s)>1$, every factor in~\eqref{eq:euler} is nonzero and the infinite product converges in the required manner. Consequently $\zeta$ has no zeros there. That is a start, but it is still far from the conjecture. The critical line lies at $\Rea(s)=1/2$, strictly beyond the region of absolute convergence.

A finite sum can vanish even when none of its terms does. For example,
\[
1+z+z^2+\cdots+z^{m-1}=\frac{1-z^m}{1-z}
\]
vanishes at every nontrivial $m$th root of unity. No individual term needs to vanish. The cancellation is global, and understanding it demands structure beyond a bound on each term. In the zeta problem the relevant variable is $s$, and the competing phases of $n^{-s}=n^{-\sig}e^{-it\log n}$ make cancellation unavoidable.

A direct cancellation argument would require a uniform estimate in a region where the defining Dirichlet series no longer converges. Analytic continuation and the functional equation describe the function in that region, but do not themselves give the required uniform bound for its prime-power coefficients. The arithmetic argument must supply quantitative information that extends beyond the half-plane of absolute convergence.

\section{Functional-equation symmetry about the critical line}
The location $1/2$ is not arbitrary. Introduce the completed zeta function
\begin{equation}\label{eq:xi}
\xi(s)=\frac12s(s-1)\pi^{-s/2}\Gamma\!\left(\frac{s}{2}\right)\zeta(s).
\end{equation}
It extends to an entire function and satisfies the functional equation
\begin{equation}\label{eq:fe}
\xi(s)=\xi(1-s).
\end{equation}
The map $s\mapsto1-s$ is a half-turn about $s=1/2$: it reverses both the horizontal displacement and the imaginary part. Reflection across the vertical line $\Rea(s)=1/2$ is instead $s\mapsto1-\overline{s}$. Since $\xi(\overline{s})=\overline{\xi(s)}$, the functional equation together with conjugation yields this same-height reflection of zeros. The critical line is the axis of that reflection.

Symmetry alone does not place the zeros on the critical line. If $\xi(\rho)=0$, it also gives $\xi(1-\rho)=0$. It follows, using complex conjugation as well, that zeros occur in pairs reflected across the critical line, not that each pair collapses to the line. A point with real part $0.3$ is paired at the same imaginary height with one having real part $0.7$. Symmetry alone permits both.

\begin{figure}[ht]
\centering
\begin{tikzpicture}[x=7.25cm,y=.12cm,>=Latex]
\fill[blue!3] (0,-18) rectangle (1,18);
\fill[blue!9] (.5,-18) rectangle (1,18);
\draw[->,gray!75] (-.08,0)--(1.13,0) node[right,black!75]{$\Rea(s)$};
\draw[->,gray!65] (0,-19)--(0,19.5) node[above,black!75]{$\Im(s)$};
\draw[gray!50] (1,-18)--(1,18);
\draw[accent,line width=1.2pt] (.5,-18)--(.5,18);
\draw[gray!35,dashed] (0,11.5)--(1,11.5) (0,-11.5)--(1,-11.5);
\coordinate (r) at (.30,11.5);
\coordinate (rc) at (.30,-11.5);
\coordinate (fr) at (.70,-11.5);
\coordinate (fcr) at (.70,11.5);
% Three actions on rho: diagonal half-turn, horizontal same-height reflection,
% and vertical conjugation. Their endpoints are the exact symmetry quartet.
\draw[blue!75!black,very thick,-{Latex[length=2mm]}] (.318,10.8)--(.683,-10.8);
\draw[red!75!black,thick,dashed,-{Latex[length=2mm]}] (.325,11.5)--(.675,11.5);
\draw[green!45!black,thick,densely dotted,-{Latex[length=2mm]}] (.30,10.4)--(.30,-10.4);
\foreach \p in {r,rc,fr,fcr} \filldraw[black!85] (\p) circle (.017);
\node[above left,font=\scriptsize] at (r) {$\rho$};
\node[above right,font=\scriptsize] at (fcr) {$1-\overline\rho$};
\node[below left,font=\scriptsize] at (rc) {$\overline\rho$};
\node[below right,font=\scriptsize] at (fr) {$1-\rho$};
\node[below,font=\scriptsize] at (0,-.6) {$0$};
\node[below,font=\scriptsize] at (.5,-.6) {$1/2$};
\node[below,font=\scriptsize] at (1,-.6) {$1$};
\node[above,font=\scriptsize,text=accent] at (.5,18) {critical line};
\end{tikzpicture}

{\footnotesize\textcolor{blue!75!black}{\rule{10pt}{1pt}} $\rho\mapsto1-\rho$ (half-turn)\quad
\textcolor{red!75!black}{---} $\rho\mapsto1-\overline\rho$ (critical-line reflection)\\[3pt]
\textcolor{green!45!black}{$\cdots$} $\rho\mapsto\overline\rho$ (conjugation)}
\caption{The fourfold symmetry of a hypothetical nontrivial zero. The functional equation sends $\rho$ to $1-\rho$ by a half-turn about $1/2$; conjugation sends it to $\overline\rho$; composing these operations gives the same-height reflection $1-\overline\rho$ across the critical line. The shaded right half is the region that the desired nonvanishing theorem excludes. This is schematic, not numerical evidence for a zero.}
\label{fig:strip}
\end{figure}

It is therefore enough to prove nonvanishing on the right of the critical line, with the pole at $s=1$ treated separately. The functional equation then excludes zeros on the left, as the following theorem makes precise.

\section{Nonvanishing for Dirichlet \texorpdfstring{$L$}{L}-functions}
The Lean archive contains the following proposition, defined in the file \sourcefile{src/Development/AutomorphicArithmetic/CharacterBoxes/EveryCharacter.lean}{EveryCharacter.lean}. We write it mathematically, separating the pole exception:
\begin{equation}\label{eq:everychar}
\begin{gathered}
\text{For every Dirichlet character }\chi,\text{ and every }s\in\CC,\\
\Rea(s)>\frac12,\quad (\chi\ne\mathbf{1}\;\text{or}\;s\ne1)
\quad\Longrightarrow\quad L(s,\chi)\ne0.
\end{gathered}
\end{equation}
Here $\mathbf{1}$ means the principal character at the given modulus. The conclusion covers every Dirichlet character, rather than just the zeta function. This breadth is required both by the arithmetic construction, which selects and averages character families, and by the functional equation, which replaces a primitive character by its inverse. The theorem must therefore retain its family-wide quantifier.

For the moment, we treat~\eqref{eq:everychar} as a hypothesis and ask exactly what it gives. In doing so we separate the genuinely difficult burden of the purported proof from a classical reduction that is short enough to verify in full.

\begin{theorem}[One-sided nonvanishing implies RH]\label{thm:rh}
Suppose $\zeta(s)\ne0$ for every $s$ with $\Rea(s)>1/2$ and $s\ne1$. Then every nontrivial zero of $\zeta$ lies on the critical line.
\end{theorem}
\begin{proof}
Let $\rho$ be a nontrivial zero, so $0<\Rea(\rho)<1$. Suppose first that $\Rea(\rho)>1/2$. This contradicts the assumed nonvanishing immediately. Suppose instead that $\Rea(\rho)<1/2$. In the critical strip the factors in~\eqref{eq:xi} that relate $\zeta$ to $\xi$ are finite and nonzero; hence $\xi(\rho)=0$. By~\eqref{eq:fe}, $\xi(1-\rho)=0$. The reflected point remains in the critical strip, so $\zeta(1-\rho)=0$. But
\[
\Rea(1-\rho)=1-\Rea(\rho)>\frac12,
\]
contradicting the nonvanishing hypothesis. Therefore neither strict inequality is possible and $\Rea(\rho)=1/2$.
\end{proof}

Here ``nontrivial zero means a zero in the critical strip, not a zero at a negative even integer. The gamma factor in~\eqref{eq:xi} has poles that cancel the trivial zeros of $\zeta$; one cannot casually infer every zero of $\zeta$ from every zero of $\xi$ outside the critical strip. The full Lean reduction handles these edge cases, including the special point $s=0$ and the pole at $s=1$. The classical proof above isolates the mathematical heart of the reduction without making the exceptional cases disappear.

\section{Dirichlet characters and arithmetic progressions}
For a Dirichlet character $\chi$ modulo $q$, define in the region of absolute convergence
\[
L(s,\chi)=\sum_{n\ge1}\frac{\chi(n)}{n^s}=
\prod_p\bigl(1-\chi(p)p^{-s}\bigr)^{-1}.
\]
Characters are periodic multiplicative functions, with value zero on integers not coprime to their modulus. They allow the primes to be separated according to congruence classes. For example, the nontrivial character modulo $4$ satisfies $\chi_4(n)=0$ when $n$ is even, $\chi_4(n)=1$ when $n\equiv1\pmod4$, and $\chi_4(n)=-1$ when $n\equiv3\pmod4$. The weighted sum
\[
\sum_{n\le x}\chi_4(n)\Lambda(n)
\]
therefore measures, to first order and with prime powers included, an imbalance between primes of the two odd residue classes modulo $4$.

The character $\chi_4$ is a simple example of a general device. Multiplying the prime-power coefficients by characters separates arithmetic progressions, while the functional equations relate $L(s,\chi)$ to the function associated with $\chi^{-1}$. A theorem proved uniformly across this family can consequently be used both for prime-counting questions in progressions and for the reflection argument. This is the reason the source aims for nonvanishing of every Dirichlet $L$-function at once.

For primitive $\chi$ there is a completed function of the form
\begin{equation}\label{eq:completedL}
\Lambda(s,\chi)=\left(\frac{q}{\pi}\right)^{(s+a)/2}
\Gamma\!\left(\frac{s+a}{2}\right)L(s,\chi),
\qquad a\in\{0,1\},
\end{equation}
where $\chi(-1)=(-1)^a$. Its functional equation relates $\Lambda(s,\chi)$ to $\Lambda(1-s,\overline\chi)$ by a nonzero constant of absolute value one. On character values of modulus one, $\overline\chi=\chi^{-1}$, which explains the inverse-character notation in the Lean code.

\begin{proposition}[Primitive-character reflection]\label{prop:primitive}
Assume~\eqref{eq:everychar}. If $\chi$ is primitive and $L(\rho,\chi)=0$ with $0<\Rea(\rho)<1$, then $\Rea(\rho)=1/2$.
\end{proposition}
\begin{proof}
In the critical strip the gamma and exponential factors in~\eqref{eq:completedL} do not vanish and are finite. Thus $L(\rho,\chi)=0$ gives $\Lambda(\rho,\chi)=0$. The functional equation yields $\Lambda(1-\rho,\overline\chi)=0$, whence $L(1-\rho,\overline\chi)=0$. If $\Rea(\rho)<1/2$, the reflected point lies in the region forbidden by~\eqref{eq:everychar}. If $\Rea(\rho)>1/2$, the original zero contradicts that hypothesis. The real part must be $1/2$.
\end{proof}

The proof is conceptually the same as Theorem~\ref{thm:rh}, but the change $\chi\mapsto\chi^{-1}$ is indispensable. The Lean theorem \texttt{Formalism.VW.critLine} in \sourcefile{src/Formalism/GammaFactorNonvanishing/CritLine.lean}{CritLine.lean} implements precisely this argument. It checks the necessary gamma-factor nonvanishing and invokes the completed $L$-function functional equation.

\section{Imprimitive characters and finite Euler factors}\label{sec:imprimitive-euler}
Every Dirichlet character is induced from a primitive character whose conductor divides its modulus. Passing to a larger modulus introduces finitely many extra Euler factors. We check directly that these factors are nonzero for $\Rea(s)>0$, so that this change of modulus cannot create a zero in the critical strip.

Let $\chi$ have conductor $f$ and let $\chi^*$ be the associated primitive character. Away from a finite set of primes dividing the larger modulus but not the conductor, the Euler products agree. Analytic continuation gives an identity of the form
\begin{equation}\label{eq:imprim}
L(s,\chi)=L(s,\chi^*)\prod_{p\in S}
\bigl(1-\chi^*(p)p^{-s}\bigr).
\end{equation}
If $\chi$ is viewed modulo $q$, the finite set is $S=\{p:\ p\mid q,\ p\nmid f\}$. Primes dividing $f$ already have $\chi^*(p)=0$ and contribute no additional factor. The distinction between the modulus and the primitive conductor will matter again for auxiliary characters.

Let $\xi$ be the even quartic character modulo $16$ with $\xi(5)=i$. For $q\in\{547,727\}$ and $0\le j<4$, induce the eight primitive characters underlying $(\cdot/q)\xi^j$ to modulus $16\cdot547\cdot727$. This multiplies the product of their $L$-functions by
\[
(1-2^{-2s})(1-547^{-4s})(1-727^{-2s})^2.
\]
The nonzero values are $\pm1$ at $2$, all four fourth roots at $547$ and two copies of each sign at $727$. Given the conductor degrees $2,4,4$, the respective trace pairs $(0,2)$, $(0,0)$ and $(0,4)$ determine the polynomials by conjugate pairing. Under Dirichlet GRH these first two power sums are recovered from corrected limiting extrema of the $547/727$ race at widths $p$ and $p^2$ \parencite{Crossley2026WidthCusp}.

\begin{lemma}[Finite Euler factors are zero-free]\label{lem:euler}
If $p$ is prime, $\Rea(s)>0$, and $|z|\le1$, then $1-zp^{-s}\ne0$.
\end{lemma}
\begin{proof}
Since $p>1$,
\[
|zp^{-s}|=|z|p^{-\Rea(s)}\le p^{-\Rea(s)}<1.
\]
If $1-zp^{-s}=0$, the left side would equal $1$, a contradiction.
\end{proof}

The conclusion now follows from the elementary fact that a finite product of nonzero complex numbers is nonzero. In the critical strip, the nontrivial zeros of an imprimitive $L$-function therefore coincide with those of its underlying primitive $L$-function. This step appears explicitly in the Lean files \sourcefile{src/Development/AutomorphicArithmetic/CharacterBoxes/ClaimStrip.lean}{ClaimStrip.lean} and \sourcefile{src/Development/AutomorphicArithmetic/CharacterBoxes/Imprimitive.lean}{Imprimitive.lean}.

The finite-factor argument needs the strict inequality $|\chi^*(p)p^{-s}|<1$ in the half-plane $\Rea(s)>0$. Lemma~\ref{lem:euler} provides that inequality, and~\eqref{eq:imprim} then transfers nonvanishing from a primitive character to any character induced from it. The formal reduction records the modulus, the exceptional Euler factors, and the nonvanishing region explicitly; no further arithmetic estimate is needed for this step.

\section{The required arithmetic estimate}
The Lean declaration \texttt{UpperTheory.EveryCharacter} requires nonvanishing in $\Rea(s)>1/2$ for every Dirichlet character. The classical argument of Theorem~\ref{thm:rh} explains why this is sufficient. The mathematical work is concentrated in obtaining the one-sided assertion in the narrow region $1/2<\Rea(s)\le1$, beyond the half-plane of absolute convergence of the Euler product.

For each fixed Dirichlet character $\chi$, a centered smoothed prime sum $G_\chi$ will be associated with $L(s,\chi)$. Its Mellin transform has poles at zeros of $L$. For a finite positive family $F$, the corresponding real statistic $G_F$ has residues that cannot cancel at a common zero. The localized analytic criterion converts an upper bound for this fixed family on arbitrarily late multiplicative windows into a zero-free half-plane. The arithmetic chapters provide precisely this kind of bound.

The source of the final implication is \sourcefile{src/Supplements/Kernels/MainTheorems.lean}{MainTheorems.lean}, where \texttt{allChi} feeds \texttt{rhOfAll}. The upstream arithmetic theorem is called \texttt{everyCharacter} in \sourcefile{src/Development/AutomorphicArithmetic/CharacterBoxes/MainConclusion.lean}{MainConclusion.lean}. The rest of the exposition concerns the mathematical content of the intermediate results.

\section{The logarithmic derivative beyond absolute convergence}
In the region $\Rea(s)>1$, the Dirichlet series for $\zeta(s)$ and its Euler product are available together. Farther to the left, the expression $\sum n^{-s}$ no longer gives an ordinary convergent series. One therefore cannot argue about a hypothetical zero near the critical line by pairing terms in the original series and observing their cancellation.

The logarithmic derivative provides a better route. For $\Rea(s)>1$ it is represented by a series supported on prime powers; elsewhere it has a meromorphic continuation whose poles reveal the zeros. A weighted sum of prime powers can consequently be chosen so that an estimate on that sum excludes those poles. The next chapter introduces this weighted prime sum explicitly. It is the analytic problem against which the elaborate arithmetic constructions should be judged.

\chapter{A weighted prime sum and its Mellin transform}\label{ch:statistic}

\section{The centered ramp sum}

The half-plane argument needs an estimate strong enough to contradict the presence of a zero with real part greater than $1/2$. The von Mangoldt function supplies the arithmetic coefficients of such an estimate. For a Dirichlet character $\chi$, the logarithmic derivative of its $L$-function satisfies
\[
-\frac{L'}{L}(s,\chi)=\sum_{n\ge1}\frac{\Lambda(n)\chi(n)}{n^s},
\qquad \Rea(s)>1,
\]
for any Dirichlet character $\chi$. A sharp cutoff $\sum_{n\le y}\Lambda(n)\chi(n)$ is useful for counting but less convenient when it must later be integrated against a Mellin kernel. The source works with a smoothed sum whose weight has an elementary transform. It also turns out to be the weight naturally produced by one of the arithmetic valuation calculations.

For $y\ge1$ define the \emph{ramp}
\[
 r_y(n)=\left(\frac yn-1\right)_+
 =\begin{cases}y/n-1,&n\le y,\\0,&n>y.\end{cases}
\]
At $n=y$ the weight is zero; for smaller $n$ it equals $y/n-1$. Alp\"oge's timeline dates his choice of this centered, ramp-weighted statistic to 20 October. In the same entry he notes the useful sign in a Guinand--Weil expression; his remarks do not supply the Mellin-transform derivation, which we give below.\footnote{Alp\"oge, \emph{Remark}, p.~3, item~5. The sum in his timeline is the centered one for the trivial character.} Define
\begin{equation}\label{eq:ramp}
R_\chi(y)=\sum_{n\ge1}\Lambda(n)\chi(n)r_y(n).
\end{equation}
This is a finite sum for each $y$. The useful feature of this particular weight is its Mellin transform, which can be calculated exactly. That calculation connects the prime sum to the analytic problem.


\begin{figure}[ht]
\centering
\begin{tikzpicture}[x=6.2cm,y=2cm,>=Latex]
\draw[->] (0.11,0)--(1.38,0) node[right] {$n/y$};
\draw[->] (0.16,-.1)--(.16,1.66) node[above] {$r_y(n)$};
\draw[accent,thick,smooth,domain=.40:1,samples=60] plot (\x,{1/\x-1});
\draw[accent,thick] (1,0)--(1.30,0);
\draw[dashed,gray] (1,0)--(1,1.36);
\filldraw[accent] (1,0) circle (.6pt);
\node[below,font=\small] at (1,0) {$1$};
\node[anchor=west,font=\footnotesize] at (.47,1.05) {$r_y(n)=y/n-1$};
\node[anchor=west,font=\footnotesize] at (1.03,.32) {$r_y(n)=0$};
\end{tikzpicture}
\caption{The ramp as a function of the normalized variable $n/y$. Its cutoff at $n=y$ makes the sum finite; the factor $y/n-1$ yields the elementary Mellin transform in Lemma~\ref{lem:ramp}.}
\label{fig:ramp}
\end{figure}

\begin{lemma}[The ramp transform]\label{lem:ramp}
If $n\ge1$ and $\Rea(s)>1$, then
\[
 \int_1^\infty r_y(n)y^{-s-1}\,dy
 =\frac{n^{-s}}{s(s-1)}.
\]
\end{lemma}
\begin{proof}
The integrand vanishes below $y=n$, and hence
\[
\begin{aligned}
\int_1^\infty r_y(n)y^{-s-1}\,dy
&=\frac1n\int_n^\infty y^{-s}\,dy-\int_n^\infty y^{-s-1}\,dy\\
&=\frac{n^{-s}}{s-1}-\frac{n^{-s}}s
 =\frac{n^{-s}}{s(s-1)}.
\end{aligned}
\]
\end{proof}

The transform is a logarithmic derivative multiplied by $1/(s(s-1))$. Away from $0$ and $1$, its singularities therefore occur exactly where the $L$-function has zeros or poles. This calculation explains the choice of weight: a bound on the prime sum can be brought to bear on the possible location of zeros. The next point is to remove the unwanted singularity at $s=1$.

\section{Centering the prime sum}

For a nonprincipal primitive character $\chi$, the function $L(s,\chi)$ is nonzero at $s=1$, so $L'/L(1,\chi)$ is defined. Let
\begin{equation}\label{eq:gchi}
g_\chi(y)=R_\chi(y)+y\frac{L'}{L}(1,\chi).
\end{equation}
For the trivial character we instead set
\begin{equation}\label{eq:Gtriv}
G(y)=\sum_{n\le y}\Lambda(n)\left(\frac yn-1\right)
-y(\log y-1-\gamma),
\end{equation}
where $\gamma$ is Euler's constant. Both corrections remove the apparent singularity of the Mellin transform at $s=1$. This normalization gives the \emph{centered} prime sum used in the analytic argument.

\begin{proposition}[Mellin transform of the centered character sum]\label{prop:transform}
For nonprincipal primitive $\chi$ and $\Rea(s)>1$,
\begin{equation}\label{eq:mellin-centred}
\int_1^\infty g_\chi(y)y^{-s-1}\,dy
=-\frac1{s(s-1)}\frac{L'}{L}(s,\chi)+\frac1{s-1}\frac{L'}{L}(1,\chi).
\end{equation}
\end{proposition}
\begin{proof}
Lemma~\ref{lem:ramp}, absolute convergence, and the Dirichlet-series expression for $-L'/L$ give
\[
\int_1^\infty R_\chi(y)y^{-s-1}\,dy
=-\frac1{s(s-1)}\frac{L'}{L}(s,\chi).
\]
The additional term in~\eqref{eq:gchi} contributes $(L'/L)(1,\chi)/(s-1)$. This is exactly~\eqref{eq:mellin-centred}. In particular, the correction term removes the pole at $s=1$ because the two singular contributions have equal and opposite residues. The parallel compensation for the trivial character and the unified Mellin formulas appear in Appendix~\ref{app:mellin-normalization}.
\end{proof}

\begin{remark}[Why the normalization matters]
One should retain the full $s$-dependent factor in the correction term. The equally plausible-looking expression
$-[L'/L(s,\chi)-L'/L(1,\chi)]/[s(s-1)]$ is \emph{not} the Mellin transform of~\eqref{eq:gchi}; it would correspond to a different centering. A direct integral check, as above, prevents a small normalization error from propagating into a residue calculation.
\end{remark}

Suppose $\rho$ is a zero of $L(s,\chi)$ of multiplicity $m$. Near $s=\rho$,
\[
\frac{L'}{L}(s,\chi)=\frac{m}{s-\rho}+\text{a holomorphic function}.
\]
Consequently, for $\rho\notin\{0,1\}$, the meromorphic continuation of the Mellin transform has a pole at $\rho$ with residue
\begin{equation}\label{eq:residue}
-\frac{m}{\rho(\rho-1)}.
\end{equation}
The correction term at $s=1$ cannot cancel this pole. A hypothetical zero is therefore recorded in the prime statistic as a definite analytic singularity. The problem becomes whether certain upper bounds for the statistic can be made incompatible with that singularity.

\section{Zero detection from multiplicative windows}

How can bounds confined to arbitrarily late \emph{long windows} exclude a zero when nothing comparable is known in the gaps? The relevant theorem answers that question. In the Lean source it appears under the name
\texttt{Deduction.zeroFree\_of\_rpow\_window}. In schematic notation it states:
\begin{theorem}[Localized one-sided zero-freeness, source-level implication]\label{thm:window}
Let $F$ be a character family in the sense of the Lean development, with associated real statistic $G_F(y)$ and zero-free assertion $\operatorname{ZeroFree}(F,\theta)$. Suppose $0<\theta<1$, $\delta>0$, and there is a constant $C$ such that for every $Y_0$ there exists $Y\ge Y_0$ satisfying
\begin{equation}\label{eq:window}
G_F(y)\le C y^\theta
\qquad (Y\le y\le Y^{1+\delta}).
\end{equation}
Then $\operatorname{ZeroFree}(F,\theta)$ holds.
\end{theorem}
The theorem above reproduces the implication and quantifier order of the pinned declaration in
\sourcefile{src/Development/AutomorphicArithmetic/DedekindLFunctions/LandauTheorems/ZeroFreeOfRpowWindow.lean}{ZeroFreeOfRpowWindow.lean}; it does not independently re-prove its specialized analytic lemmas. The precise definitions of $F$, $G_F$, and the sign conditions underlying the character family are important. One should not read it as saying that \emph{an arbitrary oscillatory function} satisfies such a principle.

The mechanism comes into focus under the substitution $y=e^u$. Since $y^{-s-1}\,dy=e^{-su}\,du$, the Mellin integral becomes a Laplace transform. A multiplicative interval $Y\le y\le Y^{1+\delta}$ is carried to the ordinary interval
\[
\log Y\le u\le(1+\delta)\log Y.
\]
For the fixed real character-family statistic, write $g(y)=G_F(y)$ and $\widehat g(s)=\int_1^\infty g(y)y^{-s-1}\,dy$ in its half-plane of convergence. Repeated differentiation with respect to $s$ amplifies values of the function at large $u$:
\begin{equation}\label{eq:laplace-derivative}
\frac{(-1)^k}{k!}\frac{d^k}{ds^k}\widehat g(s)
=\int_0^\infty g(e^u)\frac{u^k}{k!}e^{-su}\,du.
\end{equation}
The weight $u^ke^{-\beta u}/k!$, for $\beta=\Rea(s)>0$, is positive and sharply concentrated. Ignoring its overall normalization, it is the density of a gamma distribution with mean $(k+1)/\beta$ and standard deviation $\sqrt{k+1}/\beta$. Its relative width is therefore of order $k^{-1/2}$.

\begin{figure}[ht]
\centering
\begin{tikzpicture}[x=1cm,y=1cm]
\draw[->] (0,0)--(10.1,0) node[right] {$u=\log y$};
\draw[->] (.3,0)--(.3,3.35) node[above] {weight};
\fill[blue!6] (3.35,0) rectangle (7.4,3.03);
\draw[blue!45,dashed] (3.35,0)--(3.35,3.07);
\draw[blue!45,dashed] (7.4,0)--(7.4,3.07);
\draw[very thick,accent,smooth,domain=0.5:9.35,samples=110]
 plot (\x,{2.86*exp(-0.47*(\x-5.3)*(\x-5.3))});
\node[below,align=center] at (3.35,0) {\scriptsize $\log Y$};
\node[below,align=center] at (7.4,0) {\scriptsize $(1+\delta)\log Y$};
\draw[->,gray] (5.35,3.08)--(5.35,2.74);
\node[above] at (5.35,3.13) {\small $k/\beta$};
\node[align=center] at (8.15,2.18) {\footnotesize mass outside\\[-2pt]\footnotesize the window is small};
\end{tikzpicture}
\caption{Differentiating the Laplace transform creates a concentrated positive weight. For the right choice of derivative order $k$, most of its mass falls in the late interval where the arithmetic estimate is available. Schematic profile, not an exact gamma density.}
\label{fig:laplace-window}
\end{figure}

A pole at real part $\beta>\theta$ would contribute rapidly growing high derivatives of the Laplace transform. The kernels in~\eqref{eq:laplace-derivative} concentrate the derivative integral near a selected large value of $u$, where the arithmetic upper bound is available. The exact localized theorem also requires the positivity furnished by the chosen character family: an upper bound alone would not constrain an arbitrary complex, sign-changing statistic in the same way. The kernel computation identifies the analytic comparison; the required source-level pole and positivity hypotheses enter through Theorem~\ref{thm:window}.

\section{Quantifiers in the window criterion}

For a fixed positive character family $F$, and after choosing $\theta>1/2$ and a sufficiently small $\eps>0$, a bound of the form
\[
G_F(y)\le y^{1/2+9\eps}
\quad\text{for }Y\le y\le Y^{3/2}
\]
may be enough, even if the construction supplies such a window only along an unbounded sequence of $Y$'s. Suppose for contradiction that an $L$-function in the family has a zero $\rho$ with $\Rea(\rho)>1/2$. Choose
\[
\frac12<\theta<\Rea(\rho).
\]
Then choose $\eps>0$ so small that $1/2+9\eps\le\theta$. If the arithmetic theorems supply arbitrarily late windows with the displayed bound, the localized theorem rules out zeros with real part greater than $\theta$, contradicting the existence of $\rho$.

The choice of $\eps$ \emph{after} the hypothetical zero matters. There is no need to prove a bound at $\eps=0$, nor to find a single window that works simultaneously for every positive $\eps$. A calculation can afford a fixed multiple of $\eps$ in its exponent, because $\eps$ may be made small at the end. A calculation cannot afford a fixed positive exponent such as $2/3$; no choice of $\eps$ makes $2/3$ approach $1/2$.

This gives the arithmetic part of the argument a sharply defined assignment. For each fixed target character and fixed $\eps$, it must select a positive character family and then bound the associated \emph{same} centered ramp sum on arbitrarily late windows, with exponent as close to $1/2$ as desired. The auxiliary fields and primes used to obtain successive windows may change; the analytic family to which Theorem~\ref{thm:window} is applied may not change from window to window. The remaining chapters explain how the theta-lift and arithmetic estimates supply these bounds.

\section{Positive residues for character families}\label{sec:positive-residues}

We have associated a meromorphic function to a single centered prime sum.  The arithmetic part of the argument will give us a sum involving \emph{several} characters, and its coefficients must be chosen with care.  A careless linear combination can cancel a shared pole, erasing the very evidence of a zero that the prime sum was meant to preserve.

For clarity, use the same letter $G_\chi$ for the centered sum of the previous section, including $G_{\mathbf1}=G$.  Take a finite list of primitive characters $\chi_1,\ldots,\chi_r$ and strictly positive real numbers $a_1,\ldots,a_r$.  Consider the real function
\begin{equation}\label{eq:positive-real-family}
H(y)=\sum_{j=1}^r a_j\Rea G_{\chi_j}(y).
\end{equation}
The word \emph{real} is important here: individual characters may be complex-valued.  Complex conjugation on the coefficients gives $\overline{G_\chi(y)}=G_{\bar\chi}(y)$ for real $y$, provided the same normalization is used.  Therefore the Mellin transform of $H$ is the sum of the transforms of $G_{\chi_j}$ and $G_{\bar\chi_j}$, each with half the weight.

\begin{proposition}[Residues of a positive character family]\label{prop:positive-poles}
Write $\Phi(s)=\int_1^\infty H(y)y^{-s-1}\,dy$ in its initial half-plane of convergence, and continue it meromorphically using the logarithmic-derivative formulas.  At a point $\rho$ with $0<\Rea(\rho)<1$, if any member of the family has a zero there, $\Phi$ has a simple pole whose residue is
\begin{equation}\label{eq:weighted-residue}
-\frac{b_\rho}{\rho(\rho-1)},\qquad
b_\rho=\frac12\sum_{j=1}^r a_j\bigl(\ord_\rho L(s,\chi_j)+\ord_\rho L(s,\bar\chi_j)\bigr)>0.
\end{equation}
At a real pole $\rho=\beta\in(0,1)$, the residue is positive.
\end{proposition}
\begin{proof}
A holomorphic function with a zero of order $m$ at $\rho$ has the local form $(s-\rho)^m h(s)$, with $h(\rho)\ne0$, and consequently its logarithmic derivative is $m/(s-\rho)+h'(s)/h(s)$.  The pole at $\rho$ in its centered Mellin transform therefore has residue $-m/[\rho(\rho-1)]$.  The real-part identity gives the displayed sum.  Every summand in $b_\rho$ is nonnegative and at least one is positive, so cancellation is impossible.  Finally $\beta(\beta-1)<0$ for $0<\beta<1$, giving a positive real residue.
\end{proof}

At any common zero, the coefficient $b_\rho$ is a sum of nonnegative zero orders with positive weights, so the pole contributions cannot cancel. This is a statement about residues of the Mellin transform, not about pointwise nonnegativity of $H(y)$: the real statistic may change sign. The arithmetic input provides an \emph{upper} bound on selected intervals, without a corresponding lower bound or control of $|H|$. The localized theorem uses precisely that one-sided information.

There are two distinct uses of positivity. In Proposition~\ref{prop:positive-poles}, positive coefficients make the residue of a common zero nonzero. In the window estimate, the statistic itself may be negative, and the positive-part inequality converts an upper bound into control of the relevant derivatives. Neither statement replaces the other: the residue calculation addresses cancellation among transforms, whereas the derivative bound addresses the contribution of a real-valued function on a selected interval.

\begin{figure}[ht]
\centering
\begin{tikzpicture}[x=.94cm,y=1cm,>=Latex]
\draw[->] (0,0)--(10.5,0) node[right] {$\Rea(s)$};
\draw[->] (.7,-.3)--(.7,3.25) node[above] {residue coefficient};
\fill[blue!6] (3.6,.03) rectangle (8.9,3.05);
\draw[densely dashed,gray] (5.2,0)--(5.2,3.02);
\draw[accent,line width=6pt] (5.2,.03)--(5.2,1.05);
\draw[green!55!black,line width=6pt] (5.2,1.05)--(5.2,2.45);
\node[anchor=west,font=\scriptsize] at (5.55,.58) {$b_{\rho,1}>0$};
\node[anchor=west,font=\scriptsize] at (5.55,1.72) {$b_{\rho,2}>0$};
\node[above,font=\scriptsize,align=center] at (5.2,2.6) {$b_\rho=b_{\rho,1}+b_{\rho,2}>0$};
\node[below,font=\scriptsize] at (5.2,-.07) {$\Rea(\rho)$};
\node[font=\scriptsize,align=center] at (8.05,1.15) {the pole cannot\\cancel};
\end{tikzpicture}
\caption{Two positive contributions to the residue coefficient are stacked at the \emph{same} common zero $\rho$: $b_{\rho,j}=\tfrac{a_j}{2}(\ord_\rho L(s,\chi_j)+\ord_\rho L(s,\bar\chi_j))$. Additional family members add nonnegative segments at that same point. The horizontal coordinate shows $\Rea(\rho)$, not distinct zero locations; the diagram does not assert that the prime sum is positive.}
\label{fig:positive-poles}
\end{figure}

\section{Gamma kernels for localized derivatives}\label{sec:gamma-kernels}

The long-window theorem uses derivatives of a Laplace transform concentrated at a large logarithmic scale.  Put $y=e^u$ and $h(u)=H(e^u)$.  The Mellin transform of $H$ becomes the Laplace transform
\[
\Phi(s)=\int_0^\infty h(u)e^{-su}\,du.
\]
The natural rough bound $|h(u)|\le C_0e^{2u}$ suffices for differentiation when $\Rea(s)>2$.  For every integer $k\ge0$,
\begin{equation}\label{eq:normalized-derivative}
A_k(s):=\frac{(-1)^k}{k!}\Phi^{(k)}(s)
=\int_0^\infty h(u)\frac{u^k}{k!}e^{-su}\,du.
\end{equation}
We shall use these derivatives at points $s=\sigma+it$, with $\sigma>2$.  The modulus of their kernel is
\begin{equation}\label{eq:gamma-weight}
w_{k,\sigma}(u)=\frac{u^k}{k!}e^{-\sigma u},\qquad u\ge0.
\end{equation}
Its integral is $\sigma^{-k-1}$.  After division by that integral we obtain a gamma probability density.  Its mean and variance can be calculated without probability theory:
\[
\frac{\int_0^\infty u^{k+1}e^{-\sigma u}\,du}{\int_0^\infty u^{k}e^{-\sigma u}\,du}=\frac{k+1}{\sigma},\qquad
\operatorname{Var}(u)=\frac{k+1}{\sigma^2}.
\]
Thus choosing $k$ proportional to a large number $U$ puts the mass near $u\asymp U$, while its standard deviation is only of order $\sqrt U$.  The selected logarithmic window $[U,(1+\delta)U]$ has width $\delta U$.  For any fixed $\delta>0$, the kernel eventually occupies only a small portion of that window.

\begin{figure}[ht]
\centering
\begin{tikzpicture}[x=.73cm,y=1.0cm,>=Latex]
\draw[->] (0,0)--(13.4,0) node[right] {$u=\log y$};
\draw[->] (.4,0)--(.4,3.55) node[above] {relative weight};
\fill[blue!8] (4.25,0) rectangle (10.35,3.25);
\draw[dashed,blue!70!black] (4.25,0)--(4.25,3.22) (10.35,0)--(10.35,3.22);
\node[below,font=\small] at (4.25,0) {$U$};
\node[below,font=\small] at (10.35,0) {$(1+\delta)U$};
\draw[gray!60,thick,smooth,domain=1.0:12.5,samples=100] plot (\x,{1.35*exp(-((\x-7.15)^2)/7.5)});
\draw[accent,very thick,smooth,domain=2.2:12.1,samples=130] plot (\x,{2.8*exp(-((\x-7.15)^2)/1.2)});
\node[font=\footnotesize,gray!70!black] at (2.7,1.55) {smaller $k$};
\node[font=\footnotesize,text=accent] at (8.65,2.65) {larger $k$};
\node[font=\footnotesize,align=center] at (7.1,.45) {window where an\\upper bound is available};
\end{tikzpicture}
\caption{Illustration of different gamma-kernel concentrations at one center, using different pairs $(k,\sigma)$. This comparison is schematic: in the actual localization argument $\sigma$ is fixed once $\delta$ is chosen, and $k$ varies with $U$, so the kernel center moves to each successive window. The width is $O(\sqrt U)$ for $k\asymp U$, smaller than the window width $\delta U$.}
\label{fig:gamma-window}
\end{figure}

The arithmetic window in the original variable is $Y\le y\le Y^{1+\delta}$, where $Y=e^U$.  Since taking logarithms turns it into the shaded interval in Figure~\ref{fig:gamma-window}, a derivative of the right order can concentrate on precisely the part of the prime sum that we control.  The method does not require knowing the sum well everywhere.

\section{The one-sided window estimate}\label{sec:one-sided-window}

For a real function, its positive part gives the following one-sided derivative inequality.  Because $h(u)$ is real, its positive part $h_+(u)=\max(h(u),0)$ satisfies
\begin{align}\label{eq:positive-part-inequality}
|A_k(\sigma+it)|+\Rea A_k(\sigma)
&\le\int_0^\infty (|h(u)|+h(u))w_{k,\sigma}(u)\,du\\
&=2\int_0^\infty h_+(u)w_{k,\sigma}(u)\,du.
\end{align}
The first term on the left is bounded by the triangle inequality, while the second is the integral of $h$ against the same \emph{nonnegative} real kernel.  They are combined before the inequality is applied.  In particular, where $h$ is negative, its contribution disappears from the right side.  This is the elementary reason an upper bound may suffice for a statistic that changes sign.

\begin{figure}[ht]
\centering
% A schematic real function and its positive part; no mathematical data implied.
\begin{tikzpicture}[x=.93cm,y=.76cm,>=Latex,every node/.style={font=\footnotesize}]
\fill[blue!4] (3.35,.52) rectangle (9.35,5.8);
\draw[dashed,blue!45] (3.35,.52)--(3.35,5.8) (9.35,.52)--(9.35,5.8);
\node[anchor=west,font=\scriptsize] at (3.35,6.03) {$U$};
\node[anchor=east,font=\scriptsize] at (9.35,6.03) {$(1+\delta)U$};
\draw[->,black!70] (.5,4.35)--(11.8,4.35) node[right] {$u$};
\node[anchor=east] at (.4,5.42) {$h(u)$};
% Fill exactly the positive part of a smooth schematic h using top-half clip.
\begin{scope}
  \clip (.6,4.36) rectangle (11.6,5.82);
  \fill[accent!19] (.6,4.35) -- plot[domain=.6:11.5,samples=220,smooth] (\x,{4.35+.85*sin(72*\x)+.18*cos(165*\x)}) -- (11.5,4.35)--cycle;
\end{scope}
\draw[very thick,black!64,smooth] plot[domain=.6:11.5,samples=250] (\x,{4.35+.85*sin(72*\x)+.18*cos(165*\x)});
\node[font=\scriptsize,text=accent] at (8.3,5.62) {shaded: $h_+(u)$};
% Below: separate positive kernel concentrated inside the chosen interval.
\draw[->,black!65] (.5,1.02)--(11.8,1.02) node[right] {$u$};
\node[anchor=east,align=right,font=\scriptsize] at (.4,2.27) {positive\\kernel};
\fill[accent!11] (1.0,1.02) -- plot[domain=1.0:11.1,samples=160] (\x,{1.02+1.65*exp(-.63*(\x-6.45)*(\x-6.45))}) --(11.1,1.02)--cycle;
\draw[very thick,accent,smooth] plot[domain=1.0:11.1,samples=160] (\x,{1.02+1.65*exp(-.63*(\x-6.45)*(\x-6.45))});
\node[font=\scriptsize,text=accent] at (8.5,2.5) {$w_{k,\sigma}(u)$};
\end{tikzpicture}

\caption{The positive part of a real statistic on a selected logarithmic window. The top curve is schematic: shaded portions show $h_+(u)$, while negative excursions do not contribute to the right side of~\eqref{eq:positive-part-inequality}. Below, the positive kernel $w_{k,\sigma}$ concentrates inside the same window for an appropriate selected derivative order. Neither curve represents computed prime-sum data.}
\label{fig:positive-part-window}
\end{figure}

\needspace{9\baselineskip}
For the illustrative case of a unit coefficient, suppose that on a chosen window $[U,(1+\delta)U]$ we know
\begin{equation}\label{eq:local-upper-h}
h(u)\le e^{\theta u},\qquad 0<\theta<1.
\end{equation}
A fixed positive coefficient $C$ in front of $e^{\theta u}$ gives the same conclusion after changing the constants in the derivative bound; the source-level Theorem~\ref{thm:window} allows that coefficient.
There $h_+(u)\le e^{\theta u}$, so the window contribution to~\eqref{eq:positive-part-inequality} is at most
\[
2\int_U^{(1+\delta)U}\frac{u^k}{k!}e^{-(\sigma-\theta)u}\,du
\le\frac{2}{(\sigma-\theta)^{k+1}}.
\]
Outside the window we have only the rough bound $|h(u)|\le C_0e^{2u}$.  A quantitative concentration estimate, using exponential rather than merely quadratic tails, gives a fixed $\sigma>2$ and a choice $k\asymp U$ such that the contribution outside is $O(\sigma^{-k})$.  The fixed $\sigma$ depends on $\delta$ but not on the location $U$ of the window; the derivative order $k$ is instead selected for each window, with $k=\lceil(1+\delta/2)\sigma U\rceil$. Appendix~\ref{app:kernel} derives the required tail comparison and this dependence on $U$.

We therefore reach an inequality of the useful form
\begin{equation}\label{eq:two-derivative-bound}
|A_k(\sigma+it)|+\Rea A_k(\sigma)
\le C(\sigma-\theta)^{-k},\qquad t\in\RR,
\end{equation}
for arbitrarily large selected derivative orders $k$.  The harmless extra factor $1/(\sigma-\theta)$ has been absorbed in $C$.  The key point is that the right side depends on the line $\Rea(s)=\theta$, whereas a pole to the right of that line would force derivatives to grow faster.

To see the last sentence, suppose $\rho$ is an isolated simple pole with residue $c_\rho\ne0$, and $z$ is a point a distance $d=|z-\rho|$ from it.  After subtracting its principal part, Cauchy's estimate bounds the $k$-th normalized derivative of the holomorphic remainder by $Mr^{-k}$ for any circle of radius $r>d$ containing no other pole.  The pole itself contributes
\begin{equation}\label{eq:pole-derivatives}
\frac{c_\rho}{(z-\rho)^{k+1}}.
\end{equation}
Its magnitude is $|c_\rho|d^{-k-1}$.  If $d<\sigma-\theta$, this grows relative to the bound in~\eqref{eq:two-derivative-bound}.  For a nonreal nearest pole, one can move the point $\sigma+it$ slightly to arrange a unique nearest pole; the real-axis derivative supplies the required comparison term.  For a real pole, Proposition~\ref{prop:positive-poles} supplies a \emph{positive} residue, allowing the same contradiction to be read from the real derivative.  The two cases are handled carefully in the source's localized Landau argument, and Appendix~\ref{app:kernel} records the precise extra hypotheses.

\begin{figure}[ht]
\centering
\begin{tikzpicture}[x=1cm,y=1cm,>=Latex]
\draw[->] (.5,.25)--(10.2,.25) node[right] {$\Rea(s)$};
\draw[->] (2.0,.0)--(2.0,3.6) node[above] {$\Im(s)$};
\draw[accent,thick] (6.1,.25)--(6.1,3.4);
\draw[dashed,gray!70] (3.25,.25)--(3.25,3.4);
\filldraw[red!70!black] (4.55,2.0) circle (2pt);
\filldraw[accent] (6.1,2.0) circle (2pt);
\draw[<->,thick] (4.55,2.28)--(6.1,2.28) node[midway,above,font=\scriptsize] {$d$};
\draw[dashed,accent!70] (6.1,2.0) circle (1.95cm);
\node[font=\scriptsize,below] at (4.55,1.95) {pole $\rho$};
\node[font=\scriptsize,right] at (6.14,1.87) {$z=\sigma+it$};
\node[font=\scriptsize,above,align=center] at (3.25,3.48) {$\Rea(s)=\theta$};
\node[font=\scriptsize,above,align=center] at (6.1,3.5) {$\Rea(s)=\sigma$};
\node[font=\scriptsize,align=center] at (8.5,2.65) {distance to pole\\controls derivatives};
\end{tikzpicture}
\caption{Why a pole to the right of $\Rea(s)=\theta$ is dangerous: the pole contributes derivatives growing like $d^{-k}$, while the localized prime-sum estimate controls them on the larger scale $(\sigma-\theta)^{-k}$. The circle indicates local isolation and is schematic.}
\label{fig:pole-distance}
\end{figure}

This is a localized Landau-type comparison between a transform's singularities and one-sided control of its underlying real function \parencite{Landau1905}. Concentration in the logarithmic window replaces a global sign assumption. The window width fixes a real parameter $\sigma>2$, while its location $U$ determines the derivative order $k\asymp U$; the chosen windows may begin arbitrarily far out. Appendix~\ref{app:kernel} gives the quantitative tail estimates for these choices.

\section*{Source of the localized analytic theorem}
The exact character-family predicate and the positive statistic appearing in Theorem~\ref{thm:window} are defined in the Lean formalism; the statement used here is \texttt{Deduction.zeroFree\_of\_rpow\_window}. Its proof passes through \texttt{lemmaL} and the conversion from power windows to the analytic window hypothesis. The theorem cannot be applied to an arbitrary complex-valued prime sum without checking those definitions and the sign assumptions.


% --- BEGIN ROADMAP MERGE (source: standalone v1.0.3 text master) ---
% Scope: mathematical dependencies only. The source diagram is fully editable TikZ.
% Included as an unnumbered chapter so all existing counters and cross-references remain stable.
\definecolor{rmdeepblue}{RGB}{36,73,103}
\definecolor{rmmuted}{RGB}{87,98,110}
\definecolor{rmink}{RGB}{43,53,64}
\definecolor{rmframegray}{RGB}{186,203,214}
\definecolor{rmbluepanel}{RGB}{236,243,248}
\definecolor{rmgreenpanel}{RGB}{237,245,240}
\definecolor{rmsandpanel}{RGB}{248,243,233}
\definecolor{rmlavenderpanel}{RGB}{241,239,247}
\definecolor{rmgraypanel}{RGB}{245,247,248}
\definecolor{rmskywhite}{RGB}{226,239,245}
\chapter*{Arithmetic dependencies of the analytic criterion}
\phantomsection
\addcontentsline{toc}{chapter}{Arithmetic dependencies of the analytic criterion}
The localized analytic criterion in Theorem~\ref{thm:window} reduces zero-freeness to an estimate for a positively weighted family of centered prime sums on suitably late windows. The arithmetic constructions supply that estimate through three parallel sets of results, which must eventually be applied to the same family and interval.

For a primitive character $\psi$ of order $o$, let $C_\psi$ denote the set of distinct primitive characters inducing the powers $\psi^j$, $0\le j<o$. Here $\psi$ denotes a Dirichlet character, whereas $\psi(x)$ in Chapter~\ref{ch:halfplane} denoted Chebyshev's prime-counting function. In terms of the centered prime sums of Chapter~\ref{ch:statistic}, the target real statistic is
\begin{equation}\label{eq:Hpsi}
H_\psi(y)=G_1(y)+\frac{1}{3o+1}\sum_{\chi\in C_\psi\setminus\{1\}}\Rea G_\chi(y).
\end{equation}
Here $1$ denotes the trivial primitive character; a power $\psi^j$ may have smaller conductor, so its primitive inducer must be used. These are the desired positive-family components, not the auxiliary quadratic twists appearing in the larger host field.

\section*{The zero-free criterion}

The final step of the argument is a nonvanishing theorem for Dirichlet $L$-functions. In the critical strip it excludes zeros in the half-plane $\Re(s)>\tfrac12$. For the zeta function, the functional equation and complex conjugation pair a nontrivial zero $\rho$ with $1-\overline\rho$. Excluding zeros on the right therefore excludes them on the left as well. The critical line is all that remains.

The analytic theorem used for this deduction starts with a centered, weighted prime sum. Its Mellin transform involves the logarithmic derivative $-L'/L$, so a zero of the $L$-function appears as a pole. With a suitable positive family of characters, the pole contributions cannot cancel. A localized form of Landau's argument then turns an upper bound for that family on selected large intervals into a zero-free region.

This identifies what the arithmetic work must establish: a bound of sufficient strength on the \emph{same} weighted family over sufficiently large windows. The precise normalization and the localization theorem were developed in Chapter~\ref{ch:statistic}.

\section*{The three arithmetic inputs}

The first group supplies divisibility and archimedean norm comparisons for the theta-lift lattices; the second supplies admissible auxiliary fields and cohomological estimates on selected windows; the third performs the character projection and controls the residual twists. The arrows indicate where these conclusions are used in forming the same positive prime-sum estimate. The compatibility of the choices is part of the source-level arithmetic argument, rather than a consequence of the diagram.

% ----------------- Portrait proof map ----------------------------
% Everything below is editable text. No pre-rendered graphic is used.
% Original 612x792 pt coordinate layout, scaled to the 6.6x9.3 inch book trim.
% Diagram is generated in the main manuscript from editable TikZ commands.
% Every label and mathematical symbol is defined here.
\clearpage
\thispagestyle{empty}
\begin{tikzpicture}[remember picture,overlay,x=1pt,y=1pt,
  navy/.style={draw=rmdeepblue,line width=0.85pt},
  route/.style={draw=rmdeepblue,line width=1.3pt,-{Latex[length=6pt,width=5pt]}},
  card/.style={rounded corners=7pt,line width=0.9pt,draw=rmframegray},
  sec/.style={font=\sffamily\bfseries\fontsize{9}{11}\selectfont,text=rmdeepblue},
  title/.style={font=\bfseries\fontsize{14}{16}\selectfont,text=rmdeepblue},
  body/.style={font=\fontsize{11}{13.4}\selectfont,text=rmink},
  smallbody/.style={font=\fontsize{11.2}{13.6}\selectfont,text=rmink}]
\begin{scope}[shift={(current page.south west)}]
\begin{scope}[shift={(-18pt,8pt)},scale=0.81,transform shape]
  % Colors: named rather than embedded in a PDF so all artwork is editable.
  % The monograph page background is already white; do not overpaint its trim.
  \node[anchor=west,font=\bfseries\fontsize{24}{27}\selectfont,text=rmdeepblue]
    at (49,746) {How the argument is organized};
  \node[anchor=west,font=\fontsize{11.7}{14}\selectfont,text=rmmuted]
    at (50,722) {The analytic criterion is proved separately; three arithmetic arguments supply its input.};

  % Analytic theorem, logically separate from the construction below.
  \draw[card,fill=rmbluepanel] (50,624) rectangle (562,701);
  \node[sec,anchor=west] at (64,680) {ANALYTIC CRITERION};
  \node[body,anchor=north west,inner sep=0pt,text width=480pt,align=left]
    at (64,667) {A positively weighted, centered prime-sum family has a Mellin transform
    with poles at zeros. A suitable one-sided bound on sufficiently large windows
    excludes zeros with real part greater than $\tfrac12$.};

  \node[sec,anchor=west] at (50,603) {ARITHMETIC INPUTS};
  % Row 1: the theta chain and its size control.
  \draw[card,fill=rmsandpanel] (50,488) rectangle (284,581);
  \node[title,anchor=west] at (63,558) {Theta lifts and periods};
  \node[smallbody,anchor=north west,inner sep=0pt,text width=206pt,align=left]
    at (63,544) {Iterated lifts supply multiplicity-like congruences. Period identities
    control their archimedean cost.};
  \draw[card,fill=rmgreenpanel] (330,488) rectangle (562,581);
  \node[title,anchor=west] at (343,558) {Integral lattices and slopes};
  \node[smallbody,anchor=north west,inner sep=0pt,text width=204pt,align=left]
    at (343,544) {Divisibility gains are compared with size through lattice, frame and
    slope inequalities.};
  \draw[route] (289,535) -- (323,535);

  % Row 2: a different enlargement argument, followed by the window theorem.
  \draw[card,fill=rmlavenderpanel] (50,377) rectangle (284,470);
  \node[title,anchor=west] at (63,447) {Galois representations};
  \node[smallbody,anchor=north west,inner sep=0pt,text width=206pt,align=left]
    at (63,433) {A separate enlargement argument reduces suitable odd-rank information
    to smaller rank.};
  \draw[card,fill=rmbluepanel] (330,377) rectangle (562,470);
  \node[title,anchor=west] at (343,447) {Cohomological windows};
  \node[smallbody,anchor=north west,inner sep=0pt,text width=204pt,align=left]
    at (343,433) {Local--global obstructions are bounded on an interval with a fixed
    auxiliary-prime pair.};
  \draw[route] (289,424) -- (323,424);

  % Row 3: the host field and the Fourier projection.
  \draw[card,fill=rmgraypanel] (50,266) rectangle (284,359);
  \node[title,anchor=west] at (63,336) {Auxiliary character families};
  \node[smallbody,anchor=north west,inner sep=0pt,text width=206pt,align=left]
    at (63,322) {Suitable fields and finite character groups provide the averaging data.};
  \draw[card,fill=rmgreenpanel] (330,266) rectangle (562,359);
  \node[title,anchor=west] at (343,336) {Finite Fourier extraction};
  \node[smallbody,anchor=north west,inner sep=0pt,text width=204pt,align=left]
    at (343,322) {Orthogonality isolates the required Dirichlet characters with positive
    family weights.};
  \draw[route] (289,313) -- (323,313);

  % The shared rail means ALL three results are needed for the estimate.
  \draw[draw=rmmuted,line width=1.1pt] (575,535)--(575,238);
  \foreach \yy in {535,424,313}{\draw[draw=rmmuted,line width=1.1pt] (562,\yy)--(575,\yy);}
  \draw[route] (575,238)--(522,238);

  % Quantitative output: real vector mathematical text, no image object.
  \draw[rounded corners=7pt,fill=rmdeepblue,draw=rmdeepblue] (72,161) rectangle (521,244);
  \node[anchor=west,font=\sffamily\bfseries\fontsize{9}{11}\selectfont,text=rmskywhite]
    at (87,222) {COMMON ARITHMETIC ESTIMATE};
  \node[anchor=west,font=\fontsize{11.3}{13.5}\selectfont,text=white]
    at (87,203) {Compatible choices give a bound on arbitrarily late windows:};
  \node[anchor=center,font=\fontsize{15}{18}\selectfont,text=white]
    at (296.5,178) {$\displaystyle H_{\psi}(y)\le y^{1/2+9\varepsilon}
      \qquad (Y\le y\le Y^{3/2})$};

  \draw[route] (296,161)--(296,145);
  \draw[card,fill=rmgraypanel] (106,90) rectangle (486,143);
  \node[title,anchor=west] at (121,121) {Apply the analytic criterion};
  \node[body,anchor=west] at (121,103)
    {No zeros of the required $L$-functions for $\Re(s)>\tfrac12$};
  \draw[route] (296,90)--(296,76);
  \draw[card,fill=rmbluepanel] (106,24) rectangle (486,74);
  \node[title,anchor=west] at (121,52) {Functional-equation symmetry};
  \node[body,anchor=west] at (121,35)
    {Every nontrivial zeta zero lies on the critical line};
\end{scope}
\end{scope}
\end{tikzpicture}
\null
\clearpage
% ---------------- End of editable proof map ---------------------

\begingroup\linespread{1.035}\selectfont
\section*{Multiplicity, size, and local arithmetic}

Stepanov's method supplies the guiding idea for the theta-lift construction: extra multiplicity makes a counting argument stronger. Over the integers, ordinary differentiation does not provide the multiplicity used in the classical function-field setting. Alpöge instead uses long chains of theta lifts to create related arithmetic conditions. Seesaw period identities help control the cost of iteration.\footnote{For Alpöge's discussion of the multiplicity idea and chains of theta lifts, see \href{https://alpo.ge/remark.pdf}{\emph{Remark}} (2026).}

The resulting finite-prime divisibility must be compared with the archimedean sizes of lattice maps. Frame inequalities control the combined images of many subspaces; arithmetic slopes quantify the loss. The comparison also requires exact cancellations among ramification, level, and period terms. Without them, the final exponent would be too weak.

\section*{Galois representations and cohomological windows}

The Galois argument contains a separate use of enlargement. Under suitable hypotheses, two enlarged representations with disjoint auxiliary characters allow information at odd rank to descend to a smaller rank. This is distinct from the long theta-lift chains used to supply multiplicity.

The arithmetic input must hold across an interval. Cohomological lengths measure certain global classes that vanish on restriction to selected local places. The window theorem controls these obstructions while the field may vary with the scale. The same auxiliary-prime pair must work throughout the interval.

\section*{Character averaging and the final estimate}

Finite Fourier analysis decomposes the arithmetic data into character components. Orthogonality selects the required Dirichlet characters with positive coefficients. That positivity belongs to the family weights and the relevant pole residues, not to the value of the prime sum at every scale.

The geometric, cohomological, and Fourier estimates must apply to a common family and interval. Their intended output is the centered prime-sum bound illustrated in the diagram, with exponent $\tfrac12+9\varepsilon$ on selected windows $[Y,Y^{3/2}]$ extending arbitrarily far. The exact numerical specialization is conditional on the source parameters discussed in Chapter~\ref{ch:fourier}; a fixed coefficient of $\varepsilon$ is the essential analytic requirement. The estimate is not asserted at every large argument.
\endgroup
% --- END ROADMAP MERGE ---

\part{Multiplicity and automorphic arithmetic}

\chapter{Multiplicity in finite-field arguments}\label{ch:multiplicity}
\section{The multiplicity bound}
Suppose a nonzero polynomial $P(X)$ over a field has degree at most $D$. If it has $M$ distinct roots, then $M\le D$. The bound is familiar, but it deserves a closer look: the degree limits how many independent vanishing conditions a nonzero polynomial can satisfy.

We gain more by counting the order of each zero, rather than simply counting distinct roots. If $a$ is a root of multiplicity $m$, then $(X-a)^m$ divides $P$. Distinct roots contribute independent factors, and their degrees add.

\begin{lemma}[Multiplicity budget]\label{lem:mult}
Let $P\ne0$ be a polynomial of degree at most $D$ over a field $K$. Suppose $a_1,\ldots,a_M\in K$ are distinct and $P$ vanishes to order at least $m_j$ at $a_j$. Then
\[
\sum_{j=1}^M m_j\le D.
\]
In particular, if every multiplicity is at least $m$, then $M\le D/m$.
\end{lemma}
\begin{proof}
For each $j$, $(X-a_j)^{m_j}$ divides $P$. As the linear factors are pairwise coprime, their product divides $P$. Taking degrees gives the asserted inequality.
\end{proof}

The lemma reveals why multiplicity can act as a mathematical amplifier. One point at which $P$ vanishes simply costs one unit of degree. One point at which $P$ vanishes to order ten costs ten units. If a counting argument can force order-ten vanishing at each solution, the same bound on $D$ permits at most one tenth as many solutions.

The inequality itself is elementary. The difficulty in an application is to construct a nonzero auxiliary polynomial which vanishes repeatedly at many relevant points while keeping its degree under control. Stepanov's method makes this possible by exploiting additional algebraic relations over finite fields \parencite{Stepanov1969,Bombieri1974}.

\section{A finite-field example motivating Stepanov's method}
Let $p$ be prime. Over $\F_p$, Fermat's identity gives $a^p=a$ for every $a\in\F_p$. Hence
\[
F(X)=X^p-X=\prod_{a\in\F_p}(X-a)
\]
vanishes at all $p$ field elements. Its formal derivative is $F'(X)=pX^{p-1}-1=-1$, so these roots are simple. For an integer $m\ge1$, now form
\[
P_m(X)=(X^p-X)^m.
\]
The degree of $P_m$ is $mp$, and each of the $p$ field elements is a root of multiplicity exactly $m$. The multiplicity budget becomes an equality: $mp=mp$.

\begin{figure}[ht]
\centering
% A concrete naive Stepanov budget: F_5, m=3, degree 15.
\begin{tikzpicture}[x=.96cm,y=.84cm,>=Latex,every node/.style={font=\footnotesize}]
\node[anchor=west,font=\footnotesize\bfseries] at (.25,4.65) {The naive construction over $\F_5$};
\node[anchor=west,text=black!65] at (.25,4.10) {$P_3(X)=(X^5-X)^3$};
\draw[gray!65] (.75,1.28)--(10.2,1.28);
\foreach \x/\a in {1.5/0,3.15/1,4.8/2,6.45/3,8.1/4} {
  \foreach \j in {0,1,2} {
    \filldraw[fill=blue!13,draw=accent!70,line width=.5pt] (\x,{1.48+.54*\j}) rectangle ++(.77,.54);
  }
  \node[below] at (\x+.385,1.26) {$\a$};
}
\node[anchor=west,align=left,text width=3.0cm] at (9.08,2.75) {Five roots, each of\\vanishing order three};
\draw[->,accent] (9.10,2.28)--(8.9,2.28);
\draw[decorate,decoration={brace,amplitude=4pt},black!65] (1.33,1.46)--(1.33,3.10);
\node[anchor=east,font=\scriptsize,align=right] at (1.2,2.25) {$m=3$};
\draw[gray!45] (.55,.50)--(12.25,.50);
\node[anchor=west] at (.75,-.05) {Total vanishing order: $5\cdot3=15$};
\node[anchor=west] at (6.25,-.05) {Polynomial degree: $\deg P_3=15$};
\end{tikzpicture}

\caption{The naive auxiliary polynomial for $p=5$ and $m=3$. Every element of $\F_5$ is a root of order three, giving total multiplicity $15$. The polynomial $(X^5-X)^3$ also has degree $15$, so the degree bound yields no improvement. Stepanov's method needs a more efficient auxiliary construction; the figure does not model that deeper construction.}
\label{fig:stepanov-budget}
\end{figure}

For all that extra vanishing, the bound remains $p$. Multiplicity and degree have risen in exactly the same proportion. Stepanov's method needs repeated vanishing at a smaller degree cost. Over a curve, the relevant measure of complexity is often the degree of a divisor of poles; Frobenius relations and the geometry of the curve help arrange repeated vanishing at rational points. If the total multiplicity that would be required exceeds this degree, the assumed number of rational points is impossible. The arithmetic content lies in a special construction that makes repeated vanishing possible while keeping the pole degree small enough to yield a useful restriction on the number of rational points.  On a smooth projective curve, however, a nonzero function with poles of total degree at most $d$ cannot vanish to order at least $m$ at $N$ distinct rational points unless $mN\le d$.  Thus the improvement cannot come from choosing a generic auxiliary space with more than $mN$ free coefficients and pole degree below $mN$; one needs the special Frobenius equations and a separate nonvanishing argument.

The two quantities to track are thus distinct:
\begin{equation}\label{eq:tradeoff}
\boxed{\text{number of constrained points}\times\text{vanishing order}
\ \le\ \text{algebraic complexity}.}
\end{equation}
The method succeeds when the auxiliary construction makes the number of imposed zero conditions large relative to its degree. Hasse derivatives, or an equivalent characteristic-$p$ jet formalism, express those conditions without the factorial obstruction of ordinary repeated differentiation; the degree comparison turns them into a bound.

\section{The limits of polynomial multiplicity over the integers}
For a curve over $\F_p$, multiplicity is the order of vanishing of a polynomial or section at an algebraic point, and its repeated vanishing gives literal linear conditions. A corresponding polynomial with a geometric coordinate and a useful algebraic order of vanishing is not available for the prime powers $n$ weighted by $\Lambda(n)$. The latter are discrete multiplicative data, organized instead by local factors, conductors, and characters. Analytic differentiation in a complex variable does not supply the same algebraic multiplicity invariant.

Alp\"oge puts the obstruction plainly: ``multiplicity is hard'' over the integers.\footnote{Alp\"oge, \emph{Remark}, p.~1; the short quotation omits the rest of his sentence.} Over finite fields, an auxiliary function can vanish to high order at each rational point. A prime sum has no comparable geometric list of points equipped with a natural derivative. In the announced argument, consecutive theta lifts are intended to produce additional arithmetic congruences whose cumulative effect plays the part of multiplicity. The analogy concerns the \emph{strength of the constraints}, not an identification of theta lifting with differentiation.

The burden is twofold: make the constraints accumulate with the lifts, then show that satisfying them has a global arithmetic cost. The second part leads to a simple arithmetic comparison between divisibility and absolute size, which remains recognizable even after the objects become large spaces of automorphic forms.

\section{Theta series and Poisson summation}
Before discussing theta lifts between representation spaces, it helps to see a classical theta transformation. Let $z$ lie in the upper half-plane and define
\begin{equation}\label{eq:theta}
\theta(z)=\sum_{n\in\ZZ}e^{\pi i n^2z}.
\end{equation}
The sum converges absolutely because $\operatorname{Im}z>0$. Poisson summation gives the transformation law
\begin{equation}\label{eq:thetamod}
\theta(-1/z)=(-iz)^{1/2}\theta(z),
\end{equation}
with the standard holomorphic choice of square root. It relates two apparently different regimes: $z$ and its reciprocal with a minus sign.

For intuition, take $z=it$ with $t>0$. Then
\[
\theta(it)=\sum_{n\in\ZZ}e^{-\pi n^2t},
\qquad
\theta(i/t)=\sum_{n\in\ZZ}e^{-\pi n^2/t},
\]
and the transformation law reads
\begin{equation}\label{eq:poisson}
\sum_{n\in\ZZ}e^{-\pi n^2t}
=t^{-1/2}\sum_{n\in\ZZ}e^{-\pi n^2/t}.
\end{equation}
Poisson summation relates this theta series to a dual lattice sum because the Fourier transform of a Gaussian is again Gaussian. The displayed identity provides a concrete example of a function expressed in two coordinate descriptions with different scaling behavior. It is a prototype for the use of theta kernels, without establishing the integral or nonvanishing properties required of the automorphic lifts below.

The classical transformation~\eqref{eq:thetamod} should be distinguished from the theta correspondence of automorphic representation theory. The latter relates representations of two groups in a dual pair through a representation known as the Weil representation. Nonetheless, the classical theta transformation gives an honest reason for the terminology and a first glimpse of the mechanism: a Fourier-analytic kernel mediates between two mathematical descriptions.

\section{Two enlargement constructions}
The source uses enlarged representations for two purposes. In the main quantitative construction, one repeats theta lifts through a sequence of definite unitary groups. Each new stage is supposed to contribute another layer of arithmetic divisibility. The author's proposed chain has the schematic form
\[
U(X_0)\longrightarrow U(X_1)\longrightarrow\cdots\longrightarrow U(X_m).
\]
The arrows denote theta correspondences between automorphic data attached to the indicated groups; they are not homomorphisms of the unitary groups themselves.
A separate construction, formalized under the name \texttt{TwoUp}, starts from an object in even rank, under additional parity, field-degree, and local-unramified hypotheses, and forms \emph{two} enlargements in odd rank. Its purpose is to recover a compatible Galois representation, not directly to obtain a stronger prime-counting estimate. We examine that parity argument in Chapter~\ref{ch:galois}. Keeping these operations distinct prevents an attractive verbal analogy from hiding their quite different hypotheses.

\section{Divisibility and norm costs under iteration}
High multiplicity is useful only when it can be obtained without paying too much for the auxiliary construction. For a polynomial, an increase in the prescribed order of vanishing introduces more linear conditions and often forces a larger degree. Stepanov's method succeeds by finding an auxiliary space whose dimension and effective degree do not grow at the same rate.

A repeated theta lift has a different cost. Even if each stage supplies an additional congruence, the representation changes, its local parameters acquire extra constituents, and its norm may grow. To compare the resulting divisibility with a global size estimate, the argument needs integral structures, local congruence formulas, and identities for periods and norms. Alp\"oge explicitly names a Rallis seesaw, Bost's slope inequality, and detailed mass and covolume computations in this part of his account. Their role becomes intelligible once we write down the elementary inequality relating valuations and absolute size in Chapter~\ref{ch:slopes}.

The two enlargement constructions have different jobs. The pair of odd-rank enlargements described in Chapter~\ref{ch:galois} allows one to recover Galois information in even rank. The long sequence of theta lifts discussed below is meant to strengthen an arithmetic estimate by repetition. The operations are related through representation theory, but they serve different purposes.

\chapter{From multiplicity to repeated theta lifts}\label{ch:stepanov}

\section{Hasse derivatives and linear conditions}

For a polynomial over a finite field, the conditions $P(a)=P^{[1]}(a)=\cdots=P^{[m-1]}(a)=0$ are linear in its coefficients. A point can therefore impose up to $m$ linear restrictions on one auxiliary function; these restrictions need not be independent on a chosen auxiliary space. In the prime-counting problem there is no obvious function space in which the prime powers impose an analogous family of restrictions. Analytic differentiation with respect to $s$ does not by itself give the required arithmetic divisibility. The challenge is to find a construction that supplies repeated constraints by a different route.

Let $P(T)$ be a nonzero polynomial of degree at most $D$ over a field. If $P$ vanishes at $N$ distinct points and vanishes to order at least $m$ at each of them, then
\begin{equation}\label{eq:stepanov-count}
 mN\le D.
\end{equation}
The proof is almost too simple to notice. Each zero $a$ contributes a factor $(T-a)^m$, and those factors are relatively prime for distinct $a$. Their product divides $P$, and that product has degree $mN$.

But the usefulness of~\eqref{eq:stepanov-count} depends on finding the polynomial $P$. A natural strategy is to construct it from a space of auxiliary functions with many adjustable coefficients. One imposes linear conditions at the points where high-order vanishing is wanted. If the number of coefficients is larger than the number of conditions, a nonzero solution exists. At the same time, the degrees of the auxiliary functions must be kept sufficiently small. This is the \emph{budget} of Stepanov's method: the number of adjustable coefficients pays for the imposed vanishing, while degree pays for the final root count.

\begin{proposition}[A local linear-algebra model]\label{prop:jet}
Let $V$ be a vector space of polynomials over a field with dimension $M$. Let $a_1,\ldots,a_N$ be distinct points, and let $m\ge1$. If $M>mN$, there exists a nonzero $P\in V$ such that
\[
P^{[j]}(a_i)=0\qquad(1\le i\le N,\ 0\le j<m),
\]
where $P^{[j]}(a)$ denotes the coefficient of $(T-a)^j$ in the expansion of $P$ at $a$.
\end{proposition}
\begin{proof}
For each $i,j$, extracting the coefficient $P^{[j]}(a_i)$ is a linear functional on $V$. There are $mN$ such functionals. The combined linear map $V\to K^{mN}$ has a nontrivial kernel whenever $\dim V=M>mN$. Any nonzero element of this kernel has the required vanishing. The use of Hasse derivatives $P^{[j]}$ avoids difficulties with factorials in positive characteristic.
\end{proof}

For the whole polynomial space $V=K[T]_{\le D}$ the dimension is $D+1$. The condition $D+1>mN$ then guarantees an auxiliary polynomial only when the elementary multiplicity estimate already permits one. The elementary dimension count leads back to the same barrier. The useful auxiliary spaces exploit equations of the curve and powers of Frobenius. Such constructions use special equations and Frobenius relations to enforce vanishing conditions, together with a separate proof that the resulting auxiliary function is nonzero. The generic rank--nullity model by itself cannot provide more than $mN$ independent coefficients in a space of functions with pole degree below $mN$ on a complete curve; the geometry and additional algebraic identities, not a dimension count alone, are essential.

On a finite field $\F_q$, the identity $x^q=x$ for $x\in\F_q$ makes certain high-degree expressions behave like lower-complexity functions when evaluated at rational points. More generally, the geometry of a curve supplies relations that can reduce the effective number of independent high-order conditions. Agreement of values under Frobenius is not agreement of Hasse jets: for instance $T^q$ and $T$ agree on $\F_q$ but their first Hasse derivatives differ.  Higher-order vanishing requires additional identities. The resulting bound is not a generic application of linear algebra to arbitrary polynomials.

\section{The absence of an algebraic derivative for prime sums}

For the prime-power statistic $\sum_{n\le X}\Lambda(n)\chi(n)$, no polynomial in a geometric coordinate has the required prime powers as its zeros with a useful common algebraic multiplicity. One can differentiate the associated analytic functions, but this operation is distinct from imposing vanishing conditions at points of a curve and comparing their total multiplicity with a polynomial degree. The theta construction must replace that missing algebraic mechanism.

Alp\"oge describes beginning with a Fourier-analytic toy problem in search of an analogue of repeated vanishing. He later found ``multiplicity 2'' in a determinant construction on $O(2,n)$, an observation that suggested further theta lifts might supply additional arithmetic constraints.\footnote{Alp\"oge, \emph{Remark}, p.~1 (the toy setting) and p.~3 (the determinant). The remark reports the line of thought; it does not give a derivation of the determinant construction.}

The automorphic construction must play the role of multiplicity without literally differentiating integers. The analogy demands two ingredients. First, there must be a construction that gains strength every time it is repeated. Second, there must be a global accounting identity that bounds how many times the construction can succeed. Repetition without the second ingredient merely produces a larger and more elaborate object.

The announced construction uses chains of theta correspondences between spaces of automorphic data associated with groups $U(X_j)$.  In the following schematic, the arrows are operators on those data, not group homomorphisms:
\begin{equation}\label{eq:theta-chain}
U(X_0)\xrightarrow{\Theta_0}U(X_1)
\xrightarrow{\Theta_1}\cdots
\xrightarrow{\Theta_{m-1}}U(X_m),
\end{equation}
where the $X_j$ are suitable Hermitian spaces and the groups are totally definite at infinity. Each lift is intended to supply another unit of arithmetic multiplicity. The mechanism is representation-theoretic, not differential: the lift transports automorphic information to a new group while changing its local and global arithmetic data in a controlled way.

The chain is different from the separate \texttt{TwoUp} construction, whose field-degree, rank-parity, and local hypotheses are described in Chapter~\ref{ch:galois}. That construction serves a parity reduction in the construction of Galois representations: two different odd-rank enlargements are used to recover even-rank data. The long chain in~\eqref{eq:theta-chain} serves a quantitative purpose. They share a language and some tools, but solve different problems.

\needspace{9\baselineskip}
\section{The scale of the theta-lift chain}
Alp\"oge reports taking the length of the lift chain to be of order $\sqrt X$ for a von Mangoldt sum of length $X$.\footnote{Alp\"oge, \emph{Remark}, p.~4, item~8. The length is reported in his account, not deduced from the general properties of theta lifts.} This choice is explicitly part of his account. The number cannot be justified from the announcement alone: it depends on how the extra divisibility gained in a lift compares with the norm and level costs imposed by that same lift.

The honest way to understand the choice is to identify the two quantities which repetition affects. Each additional stage is intended to increase an arithmetic order of divisibility; at the same time it changes periods, masses, dimensions, and the size of the integral structure supporting the lift. The final estimate depends on a precise comparison of these contributions. At present we can explain the \emph{need} for such a comparison and the form of the resulting valuations, but we cannot deduce a formula for the optimal number of stages from the few declarations reviewed here.

A simple hypothetical cost model illustrates why square-root scales often appear in amplification arguments. Suppose an error took the form $X/m+m$ for a continuous parameter $m>0$. Then $X/m+m\ge2\sqrt X$, with equality at $m=\sqrt X$; an integer parameter would also require rounding. This model is not an asserted cost estimate for Alp\"oge's theta lifts: the actual chain has rank, parity, integrality and period constraints, and its source-level cost must be established by the geometric and arithmetic theorems. The square-root scale in the announced argument therefore cannot be deduced from this elementary minimization.

The model describes only the competition between a term that decreases with the number of iterations and a term that increases with it. The actual chain requires more information: each stage changes a representation space, and the comparison must transport integral data and period normalizations between those spaces. A specific length or error exponent for the arithmetic chain must come from these source-dependent comparisons, not from assigning an interpretation to the two terms of the model after the fact.

\section{Orders of vanishing and nonarchimedean valuations}

Both kinds of multiplicity can be expressed through orders in local rings. If a polynomial $P(T)$ vanishes to order $m$ at $a$, then $(T-a)^m$ divides $P(T)$. At a nonarchimedean place, repeated divisibility is measured by the power of a uniformizer $\varpi$. Their common feature is an order, although the objects carrying that order and the estimates used to bound it are different.

For an integer $N$, its $p$-adic valuation $v_p(N)$ counts how many copies of $p$ divide it. The statement $p^m\mid N$ means $v_p(N)\ge m$. The analogy is exact at the level of valuations: $\operatorname{ord}_{T=a}P$ and $v_p(N)$ are both orders of vanishing in a discrete valuation ring.

What differs is the global comparison. For a polynomial, the sum of the multiplicities of its zeros is bounded by its degree. For a nonzero integer, a product of many prime powers dividing it is bounded by its absolute size. In its most elementary form,
\[
\sum_p v_p(N)\log p=\log|N|.
\]
The left side counts nonarchimedean multiplicity; the right side measures archimedean size. An algebraic number or a section of an arithmetic vector bundle has a more complicated version of the same balance, governed by the product formula and slope inequalities.

The announced strategy aims to exploit that comparison in a sophisticated setting. A chain of theta lifts should force unusually strong divisibility at many primes. Independent estimates of compatible integral maps, lattice structures, and periods at infinite places should prevent the divisibility from being arbitrarily large.  Iterating a theta operator does not by itself give those bounds: the source imposes separate local-index, nonvanishing and rank, period, frame, and slope hypotheses. The incompatibility becomes an upper bound for the sum of primes that supplied the divisibility in the first place.

A controlled archimedean upper bound places an upper bound on total prime-weighted divisibility. If a nonzero arithmetic quantity is divisible by many prime powers, its finite-place valuations contribute to its arithmetic size; the product formula requires those contributions to be balanced by its other absolute values. The theta-lift argument arranges an upper bound on the archimedean terms, making this balance quantitative.

\section{Additional hypotheses for the arithmetic comparison}
The local-order interpretation of multiplicity is exact, but it does not yet provide an estimate for prime sums. In the arithmetic setting one needs a nonzero integral object, a lower bound for its divisibility at appropriate primes, and an upper bound for its archimedean size. We next describe the theta operators to which these comparisons are applied. The later period and slope calculations, together with independent local-divisibility and compatibility inputs, are responsible for turning their repeated application into a quantitative statement.  The required formal interfaces include \texttt{LinkDIV} and \texttt{LinkMu} from \sourcefile{src/Formalism/SharpRestriction/GlobalTheoremInputs.lean}{GlobalTheoremInputs.lean} and the local-to-global comparison in \sourcefile{src/Formalism/LocalToGlobal/Divisibility/RecBridge.lean}{RecBridge.lean}.

\chapter{Theta correspondence on definite unitary groups}\label{ch:theta}

\section{Totally definite unitary groups}

Theta lifting moves automorphic forms between different groups \parencite{Howe1979}. To see the construction in a setting where its arithmetic properties can be measured, let $E$ be a CM field with complex conjugation $c$, and let $E^+$ be its totally real fixed field. A Hermitian form on $E^n$ is a function
\[
 h: E^n\times E^n\longrightarrow E
\]
which is linear in one argument, conjugate-linear in the other, and satisfies $h(y,x)=c(h(x,y))$. Its unitary group is
\[
U(h)=\{g\in\operatorname{GL}_n(E):h(gx,gy)=h(x,y)\}.
\]
For $E=\CC$ and the standard positive-definite form $h(x,y)=\sum x_i\overline{y_i}$, this is the familiar compact unitary group $U(n)$. Over a CM number field one has a unitary group at each embedding of $E^+$ into $\RR$. When the Hermitian form is positive definite at every real place, the group is called \emph{totally definite}.

Definiteness removes a formidable analytic burden. In many automorphic settings an automorphic form lives on a noncompact quotient and must satisfy intricate growth and cuspidality conditions. A definite unitary group is compact at infinity, so that, for fixed finite-level conditions, its automorphic forms can be modeled on a \emph{finite set of arithmetic classes}. A familiar precursor is the finite set of ideal classes of a number field. In the present setting the classes may be represented by suitable Hermitian lattices.

Let $\mathcal C$ be such a finite class set. For the simplest scalar-valued model, a form is a function $f:\mathcal C\to\CC$. An automorphism group $\operatorname{Aut}(M)$ is attached to each class $M$. The natural mass-weighted inner product is
\begin{equation}\label{eq:massip}
\langle f,g\rangle_{\mathcal C}
=\sum_{M\in\mathcal C}\frac{f(M)\overline{g(M)}}{|\operatorname{Aut}(M)|}.
\end{equation}
The denominators matter because they record stabilizers when the class set is viewed as a quotient. They are also a reason the integral structure is subtler than a matrix of ordinary integers; denominators appear through the automorphism orders.

\section{Theta lifting as a kernel operator}

Suppose $\mathcal C_0$ and $\mathcal C_1$ are class sets belonging to two unitary groups. A theta kernel is a function of two variables, $\vartheta(M,M')$, constructed from a theta series and a chosen Hecke character. Formally, the lifting operator takes the shape
\begin{equation}\label{eq:kerneltheta}
(\Theta f)(M')
=\sum_{M\in\mathcal C_0}
\frac{\vartheta(M,M')}{|\operatorname{Aut}(M)|}f(M).
\end{equation}
This is a finite-dimensional integral transform in the definite setting. The original theta correspondence involves global representations and local compatibility conditions, but~\eqref{eq:kerneltheta} explains the operational idea with no specialized terminology: one mixes the values of the original function using an arithmetically constructed kernel, thereby producing a function on another class set. The formula is schematic: the actual stage kernel also includes determinant-character factors and prescribed archimedean exponents, and the stabilizer denominators must be tracked through the integral model. Appendix~\ref{app:theta-technical} sketches where those factors enter and explains the stabilizer-denominator issue; it does not display the exact archimedean exponents, which remain in the cited source construction. Suppressing them in this elementary model does not remove their contributions to the arithmetic cost.
\begin{example}[A small matrix model]
Let $\mathcal C_0=\{M_1,M_2\}$ and $\mathcal C_1=\{N_1,N_2,N_3\}$, and temporarily suppose all automorphism orders are $1$. Then an operator of the form~\eqref{eq:kerneltheta} is simply
\[
\begin{pmatrix}f(M_1)\\f(M_2)\end{pmatrix}
\longmapsto
\begin{pmatrix}
\vartheta(M_1,N_1)&\vartheta(M_2,N_1)\\
\vartheta(M_1,N_2)&\vartheta(M_2,N_2)\\
\vartheta(M_1,N_3)&\vartheta(M_2,N_3)
\end{pmatrix}
\begin{pmatrix}f(M_1)\\f(M_2)\end{pmatrix}.
\]
The matrix formula is useful only after we remember where its entries come from. A theta kernel is built from arithmetic theta series and carries transformation laws under the groups involved. Those laws relate the resulting operator to local representation theory and to global period identities. They are what permits more to be said about its norm and integral values than about an arbitrary matrix of complex numbers.
\end{example}

The finite matrix model illustrates the linear operation defined by a kernel. An arbitrary matrix need not produce a nonzero lift on the required vector, nor does it carry the integral lattices and height bounds needed for the arithmetic comparison. For the theta correspondence, nonvanishing, integrality and norm control are separate mathematical inputs. Their source-level statements are needed before the kernel picture can be used in a divisibility argument.

The chain in~\eqref{eq:theta-chain} is a product of such operators,
\begin{equation}\label{eq:composite}
\Theta^{(m)}=\Theta_{m-1}\Theta_{m-2}\cdots\Theta_0.
\end{equation}
A long chain of lifts threatens a large norm at the end. Period identities and arithmetic normalizations control that cost; merely constructing each map $\Theta_j$ gives no useful bound on the product $\Theta^{(m)}$.

\section{Local spectral data under theta lifting}

A representation of a group records how that group acts on a vector space. An automorphic representation packages a family of global and local symmetries into a structured spectral object. Associated to an unramified local component are parameters which, in the case of $\operatorname{GL}_n$, may be encoded by an unordered multiset of $n$ complex numbers, the Satake parameters. Such parameters determine local Euler factors, and matching them is a basic method for comparing arithmetic representations.

The author's announcement summarizes the effect of a typical step by the symbolic relation
\begin{equation}\label{eq:isobaric}
\pi\longmapsto\pi\boxplus\mu_j[t].
\end{equation}
The notation means, roughly, that the lifted representation retains the original automorphic component while adjoining a controlled block determined by a Hecke character $\mu_j$ and a length parameter $t$. The symbol $\boxplus$ is not the elementary direct sum of vector spaces; it denotes an isobaric operation in the automorphic representation theory of general linear groups, as it appears after the relevant transfer. The notation $\mu_j[t]$ likewise abbreviates a structured representation-theoretic block. The details of this transfer are mathematically substantial and must not be suppressed when a later theorem requires their precise local hypotheses.

The added block $\mu_j[t]$ makes the lift useful because its contribution to the local spectral data can be described. It is possible to compare the enlarged representation with the initial one after accounting for these known constituents. This is similar to multiplying a characteristic polynomial by a known factor, but a genuine automorphic statement also records ramification, weights, and normalization. Later we shall see a precise version of the ``remove a known constituent'' operation in the Galois-representation chapter.

\section{Nonvanishing, integrality, and norms}
An iterated lift has little quantitative use if it vanishes identically. The first requirement is therefore to know that the relevant automorphic component survives each stage. Local theta correspondence determines which representations can occur, and global theta-lifting theorems supply the conditions under which a nonzero lift exists. These are substantive hypotheses: the matrix model of the preceding section certainly does not prevent an arbitrary vector from belonging to the kernel.

There are two further requirements. The values of the lifted form must lie in a suitable integral lattice, so that congruence statements at finite primes have a precise meaning. And the norm of the lifted form must be controlled at the archimedean places. A vector divisible by a high power of a prime but allowed to be arbitrarily large carries no contradiction; a norm bound without integrality similarly says nothing about the prime valuations.

Rallis's inner-product formula compares the norms of theta lifts with special $L$-values and local factors \parencite{GanQiuTakeda2014}. A seesaw identity transfers a period pairing to another representation space where it can be evaluated or bounded. Mass formulas and integral models control the lattices and metrics in which the resulting periods lie. These are distinct norm, period and integrality inputs; the arithmetic estimate requires them to remain compatible at successive lifts.

\section{Adjointness and seesaw period identities}

The basic reason a seesaw identity can exist is an ordinary identity about integrating a kernel in two orders. Suppose $K(x,y)$ is a sufficiently integrable function on $X\times Y$ and $f,g$ are test functions. Fubini's theorem yields
\[
\int_Y\left(\int_X K(x,y)f(x)\,dx\right)\overline{g(y)}\,dy
=
\int_X f(x)\overline{\left(\int_Y\overline{K(x,y)}g(y)\,dy\right)}\,dx.
\]
In operator language, this is $\langle Tf,g\rangle=\langle f,T^*g\rangle$. In theta correspondence, pairs of reductive groups embedded in a larger symplectic group give much richer identities of this kind. The groups, characters, measures, and convergence arguments require serious care. Nevertheless the central idea is familiar: a period of a theta lift may be evaluated by unfolding the theta kernel and interpreting the result as a different period.

\begin{figure}[ht]
\centering
\begin{tikzpicture}[>=Latex,node distance=14mm,box/.style={draw=black!30,rounded corners=3pt,fill=blue!3,minimum width=28mm,minimum height=11mm,align=center}]
\node[box] (a) {$U(X_0)$\\\scriptsize initial form};
\node[box,right=9mm of a] (b) {$U(X_1)$\\\scriptsize first theta lift};
\node[box,right=9mm of b] (c) {$U(X_m)$\\\scriptsize iterated lift};
\draw[->,thick] (a)--node[above,align=center]{\scriptsize theta kernel\\[-2pt]\scriptsize $\Theta_0$}(b);
\draw[->,thick] (b)--node[above]{\scriptsize $\Theta_1\cdots\Theta_{m-1}$}(c);
\node[below=13mm of b,align=center] (p) {\footnotesize period or norm identity\\[-2pt]\footnotesize changes the form of the comparison};
\draw[->,gray] (c.south) to[bend left=15] (p.east);
\draw[->,gray] (p.west) to[bend left=15] (a.south);
\end{tikzpicture}
\caption{What must accompany a theta-lift chain. The arrows representing lifts are not sufficient by themselves; a second family of identities compares norms and periods across the chain.}
\end{figure}

Alp\"oge says that a Rallis seesaw expresses the relevant period as a Bessel period ``independent of the Hecke characters.''\footnote{Alp\"oge, \emph{Remark}, p.~4, item~12.} The independence has a practical consequence, but the normalization is essential. Suppose the arithmetic calculation produces $\sum_\mu c_\mu P_\mu$, and the period comparison gives $P_\mu=a_\mu P$ for a common Bessel period $P$ and nonzero comparison scalars $a_\mu$. Then the identity available without further assumptions is
\[
\sum_\mu c_\mu P_\mu=P\sum_\mu c_\mu a_\mu.
\]
It reduces to $P\sum_\mu c_\mu$ only when the relevant normalizations make all $a_\mu=1$; a common scalar may instead be factored out, or the varying scalars must be included in the coefficients with appropriate bounds. Character orthogonality applied to the original $c_\mu$ cannot simply ignore varying $a_\mu$. Thus the seesaw comparison explains why a character-independent period is useful, while its precise hypotheses and quantitative normalizations remain indispensable to the later averaging.

\section{Integral data along a theta-lift chain}\label{sec:theta-chain}

We have described a theta lift as an operator between spaces of automorphic forms.  Iterating operators of this kind is formally easy: if the output of one is an acceptable input to the next, compose them.  The difficulty in this argument is that the final inequality must remember information that an ordinary statement about the existence of lifts would forget.  It must remember where an evaluation is divisible by a prime, how large an evaluation is at the infinite places, and how the integral structures at intermediate stages are related.

The obstruction noted in Alp\"oge's remark is the lack of a comparable arithmetic multiplicity over the integers.\footnote{Alp\"oge, \emph{Remark}, p.~1.}  Differentiating a polynomial over a finite field can be arranged to produce repeated vanishing.  Over the integers the proof cannot simply replace a derivative by an operator bearing the same name.  What is needed is a construction whose arithmetic consequences \emph{play the role} of increasing multiplicity.  In the announced argument, successive theta lifts supply the repeated stages, while the valuation and period comparisons give the numerical meaning to their accumulation.

At each step there are three pieces of data.  First is an automorphic space $\mathcal A_j$ of a prescribed type.  Second is a kernel operator
\[
T_j:\mathcal A_j\longrightarrow\mathcal A_{j+1},
\]
obtained by integrating an appropriate theta kernel.  Third is a lattice $M_j$ of integral or rationally normalized sections whose images can be compared with the next stage.  The arithmetic statements concern not merely the composite $T_{m-1}\cdots T_0$, but the corresponding integral maps and the denominators needed to make them well defined over a common number field.

Here is a small model of the accounting.  Suppose $v_p(f)$ denotes the $p$-adic valuation of the evaluation of a section $f$.  If an operator carries a section with $v_p(f)\ge a$ to one with $v_p(Tf)\ge a+b$, the repeated gain after $m$ steps would be $mb$ \emph{provided} the same inequality applies at every intermediate step.  The italicized condition is substantial.  In the true construction the form changes its representation, its level and its Hermitian space as the chain proceeds.  There is no single scalar operator $T$ whose powers one can estimate once and for all.

The source uses separate families of Hermitian spaces for the changing rank and for the changing signature.  In the description of the compatible construction, parameters $k,m$ are even and $t$ is odd, with inequalities ensuring that all the intermediate spaces exist in the required range.  The spaces $X_j$ have rank $k+j(t+1)$, and $Y_j$ have rank $r+jt$, where $r=k+m-1$.  For $j<m$ the construction uses isometric embeddings $X_j\hookrightarrow Y_j$; at the terminal stage $Y_m$ has codimension one in $X_m$.  A third family $W_j$, of fixed rank $r$, has signature $(k+j-1,m-j)$ for $1\le j\le m$, at each archimedean place.  The $W$-family becomes definite at its terminal stage. A fourth family $V_j$ consists of auxiliary definite Hermitian spaces, with rank $r+t$ asserted in the source for $1\le j\le m$. These families have different roles.  The changing ranks and signatures organize the theta construction; the definite spaces allow lattice and section estimates. Figure~\ref{fig:v108-theta-spaces} records the stagewise rank relations, while Figure~\ref{fig:theta-integral} distinguishes the analytic operators from the integral comparisons. Appendix~\ref{app:theta-technical} lists the accompanying structural inequalities.


\begin{figure}[H]
\centering
% Extracted from DESIGN_PROTOTYPES/FIG03_theta_space_families_book.tex
% Requires TikZ and arrows.meta.
\begin{tikzpicture}[>=Latex,font=\small,
 bx/.style={draw=black!65,fill=black!2,rounded corners=1pt,minimum width=41mm,minimum height=7mm,align=center,inner sep=1.7pt},
 sub/.style={draw=black!40,fill=black!4,rounded corners=1pt,minimum width=41mm,minimum height=7mm,align=center,inner sep=1.7pt},
 ar/.style={-{Latex[length=1.8mm]},line width=.7pt},
 lab/.style={font=\scriptsize}]
% stage cards: start at x=0 and x=5.15 cm; rows y=0 and -5.8
\foreach \cx/\cy/\tag in {0/0/0,5.2/0/1,0/-5.75/{m-1},5.2/-5.75/m}{
\draw[black!20,rounded corners=2pt] (\cx-2.42,\cy+1.02) rectangle (\cx+2.42,\cy-4.43);
\node[font=\small\bfseries] at (\cx,\cy+.57) {Stage $j=\tag$};
}
\node[bx] (x0) at (0,0) {$X_0:\ k$};
\node[bx] (y0) at (0,-1.54) {$Y_0:\ r$};
\draw[ar] (x0.south)--node[right,lab] {isometry}(y0.north);
\node[sub] at (0,-2.65) {$W_0$: no signature asserted};
\node[sub] at (0,-3.68) {$V_0$: no rank asserted};
\node[bx] (x1) at (5.2,0) {$X_1:\ k+t+1$};
\node[bx] (y1) at (5.2,-1.54) {$Y_1:\ r+t$};
\draw[ar] (x1.south)--node[right,lab] {isometry}(y1.north);
\node[sub] at (5.2,-2.65) {$W_1:\ (k,m-1)$};
\node[sub] at (5.2,-3.68) {$V_j$: rank $r+t$; definite};
\node[bx] (xp) at (0,-5.75) {$X_{m-1}:\ k+(m-1)(t+1)$};
\node[bx] (yp) at (0,-7.29) {$Y_{m-1}:\ r+(m-1)t$};
\draw[ar] (xp.south)--node[right,lab] {isometry}(yp.north);
\node[sub] at (0,-8.4) {$W_{m-1}:\ (k+m-2,1)$};
\node[sub] at (0,-9.43) {$V_j$: rank $r+t$; definite};
\node[bx] (xm) at (5.2,-5.75) {$X_m:\ k+m(t+1)$};
\node[bx] (ym) at (5.2,-7.29) {$Y_m:\ r+mt$};
\draw[ar] (ym.north)--node[right,lab] {codim. 1}(xm.south);
\node[sub] at (5.2,-8.4) {$W_m:\ (r,0)$ definite};
\node[sub] at (5.2,-9.43) {$V_j$: rank $r+t$; definite};
\node[font=\scriptsize,anchor=west] at (-2.4,-10.58) {$r=k+m-1$. Ranks are over $E$; signatures shown only for $W_j$.};
\node[font=\scriptsize,anchor=west] at (-2.4,-10.97) {Solid arrows are within-stage isometric embeddings, not theta operators.};
\end{tikzpicture}

\caption{Hermitian-space ranks and selected signatures required by the theta construction. The arrow at the terminal stage reverses because $\dim_EY_m=\dim_EX_m-1$. The displayed $V_j$ ranks and $W_j$ signatures apply only for $1\le j\le m$. All arrows shown are same-stage isometric embeddings, not theta operators; the diagram records structural source requirements rather than a separate existence proof.}
\label{fig:v108-theta-spaces}
\end{figure}

\begin{figure}[ht]
\centering
\begin{tikzpicture}[>=Latex,node distance=6mm, every node/.style={font=\small,align=center},bx/.style={draw=black!40,rounded corners=2pt,fill=blue!3,minimum width=19mm,minimum height=10mm}]
\node[bx] (a) {$\mathcal A_0$};
\node[bx,right=10mm of a] (b) {$\mathcal A_1$};
\node[bx,right=10mm of b] (c) {$\cdots$};
\node[bx,right=10mm of c] (d) {$\mathcal A_m$};
\draw[->,thick] (a)--node[above,font=\scriptsize] {$T_0$}(b);
\draw[->,thick] (b)--node[above,font=\scriptsize] {$T_1$}(c);
\draw[->,thick] (c)--node[above,font=\scriptsize] {$T_{m-1}$}(d);
\node[bx,below=12mm of a,fill=green!5] (e) {$M_0$};
\node[bx,below=12mm of b,fill=green!5] (f) {$M_1$};
\node[bx,below=12mm of c,fill=green!5] (g) {$\cdots$};
\node[bx,below=12mm of d,fill=green!5] (h) {$M_m$};
\draw[->,dashed] (e)--(f);\draw[->,dashed] (f)--(g);\draw[->,dashed] (g)--(h);
\draw[->,gray!60] (a)--(e);\draw[->,gray!60] (b)--(f);\draw[->,gray!60] (c)--(g);\draw[->,gray!60] (d)--(h);
\node[font=\footnotesize,align=left] at (4.9,-3.15) {The lower arrows require integral comparison,\\not merely the existence of upper arrows.};
\end{tikzpicture}
\caption{Two levels of a theta chain. The upper arrows are analytic theta operators between automorphic spaces; the lower arrows represent the arithmetic comparison of selected lattices. The proof needs both rows and their compatibility.}
\label{fig:theta-integral}
\end{figure}

A useful consequence of this distinction is that the \emph{number} of lifts and the \emph{cost} of making them integral are separate questions.  The suggested square-root length, reported in Alp\"oge's remark, is part of the design of the chain.  But a length of order $\sqrt X$ by itself says nothing about the final prime-sum exponent.  Denominators, discarded parts of lattices, and archimedean norms all grow during the construction; they determine whether repeated lifting actually improves the estimate.

\section{Rational period pairings and integral comparisons}\label{sec:theta-period-pairing}

An integral theta chain is useful only if its final output can be tested by a period or evaluation pairing that has a controlled algebraic realization.  For each column of the construction, the terminal pairing produces a bilinear map between a column space and one common row space.  After choosing a coefficient number field $K$, the source realizes these as rational maps
\begin{equation}\label{eq:rational-pairings}
B_u:N_u\otimes K\longrightarrow R^\vee\otimes K,
\end{equation}
where $N_u$ and $R$ are full lattices.  The \emph{same} maps must recover the analytic pairings at a chosen complex embedding.  Otherwise a complex norm estimate and a finite-prime divisibility estimate would be statements about different objects, with no reason to combine.

The comparison with a Bessel pairing is an example of this compatibility.  In the exact \texttt{BesselId} source interface there is \emph{one} unitary change of coordinates $V$ on the common row space, while for each column $u$ there may be a separate unitary map $U_u$ and nonzero complex scalar $a_u$, such that
\[
B_u(\phi,F)=a_u\,\operatorname{besS}_u(U_u\phi,VF).
\]
This identifies the pairings up to the stated transformations; it neither makes the scalars independent of $u$ nor bounds $|a_u|$.  Separately, \texttt{BuRational} asserts rationality of the column pairings over the coefficient field, and the quantitative \texttt{LinkDIV\_at} interface requires operator-height, period and local-error bounds for the \emph{same} arithmetic maps.  Rationality alone does not give uniform denominator or height estimates.  Thus the seesaw identification, algebraic realization, and quantitative estimates are distinct inputs that must be made compatible; none follows merely from the existence of a theta lift.\footnote{See \sourcefile{src/Formalism/SharpRestriction/GlobalTheoremInputs.lean}{GlobalTheoremInputs.lean}, definitions \texttt{BesselId} and \texttt{LinkDIV\_at}, and \sourcefile{src/Formalism/LocalToGlobal/Divisibility/DivisibilityRemainder/RationalColumns/BuRationalOf.lean}{BuRationalOf.lean}.}

The seesaw identity of the previous section supplies the principal simplification of the periods. The underlying Bessel period is independent of the Hecke characters introduced in the chain. Extracting a factor common to the columns nevertheless requires a \emph{separate matching of test vectors, period normalizations and column-comparison scalars}; it does not follow from \texttt{BesselId} alone. Any common factor may be extracted only \emph{after} the column-dependent comparison scalars have been accounted for. In particular, a nonzero scalar with uncontrolled magnitude cannot be dropped from a numerical estimate, and a character-dependent scalar cannot be omitted from a Fourier average. The separately stated height and period bounds are the quantitative inputs. We have not reconstructed the full joint matching theorem from its producer chain.

The next chapter's finite-place divisibility estimate and its archimedean norm bound can now be understood as two measurements of the maps~\eqref{eq:rational-pairings}.  At a finite prime, integrality asks how large an ideal divides their matrix coefficients.  At an infinite place, Hermitian norms ask how large those coefficients can be.  The arithmetic product formula makes these measurements comparable, but the comparison depends on using the identical rational maps and compatible metrics. The general arithmetic intersection framework is discussed in \textcite{GilletSoule1990}; the construction-specific normalization here is supplied by the cited formal inputs.  This is the exact point where the correspondence becomes a number-theoretic estimate rather than a representation-theoretic construction alone.

\section{Geometric and integral input hypotheses}\label{sec:theta-input-boundary}

It is possible to give a correct account of what these maps must accomplish without claiming to reconstruct their full construction.  The formal source must prove that there is a common coefficient field, suitable full positive-rank lattices, integral comparisons along the chain, uniform operator-height bounds, and near-orthogonality of the resulting family of evaluation maps.  It must also control the part discarded when passing to a preferred automorphic type.  The last condition matters because a small vector-space complement may have an unexpectedly large arithmetic degree.

The source's geometric input includes arithmetic sections, covering levels, comparison of metrics, and arithmetic Hilbert--Samuel estimates. These enter with specified coefficient fields, source and target ranks, integral lattices, period normalizations, and allowed frame defects. Appendix~\ref{app:theta-technical} identifies the named theorems and the restrictions under which their conclusions are used. The main argument applies those conclusions only after the relevant norm, rank and integrality hypotheses have been stated; a general existence theorem for theta correspondence alone would not provide these bounds.

\section{Finite class sets for definite unitary groups}

For a totally definite unitary group, the automorphic quotient at a fixed level can be represented by a finite arithmetic class set. In the useful model, automorphic forms are functions on that set, so finite-dimensional linear algebra and integral lattices become available. Their norms and covolumes can then be compared using the product formula. This is the particular role of definiteness in a construction that ultimately estimates an analytic prime-power statistic.

Prime numbers reappear through Frobenius and Hecke operators and, in the final arithmetic estimate, through valuations of the theta kernels. The definite setting makes the forms finite-dimensional objects equipped with norms and integral lattices. We can therefore compare their divisibility at finite places with their size at infinite places. The elementary version of that comparison is already present in unique factorization.

\chapter{Divisibility and archimedean size}\label{ch:slopes}

\section{Prime valuations and archimedean bounds}

Consider a nonzero integer $N$ divisible by a large product $\prod_{p\in S}p^{a_p}$. Then
\begin{equation}\label{eq:integer-product}
\sum_{p\in S}a_p\log p\le\log|N|.
\end{equation}
Unique factorization is all that goes into this inequality. Even so, it already has the form of the arithmetic comparison we need. The left records divisibility gained at finite primes; the right measures its cost at the real place. If independent reasoning bounds $|N|$ from above, then~\eqref{eq:integer-product} bounds the total prime-weighted divisibility from above.

For a nonzero algebraic integer $\alpha$ in a number field, the analogous observation uses the field norm $N_{E/\QQ}(\alpha)$. A lower bound on its nonarchimedean valuations forces a lower bound on $|N_{E/\QQ}(\alpha)|$, while
\begin{equation}\label{eq:fieldnorm}
\log|N_{E/\QQ}(\alpha)|
=\sum_{\sigma:E\hookrightarrow\CC}\log|\sigma(\alpha)|
\end{equation}
measures its conjugates at infinity. Thus \emph{a number cannot be extremely divisible at many finite places and simultaneously very small in every complex embedding}, unless it is zero. The exceptional possibility $\alpha=0$ is exactly why a proof must control nonvanishing of the constructed object.

Theta lifts do not normally produce a single algebraic integer. They produce maps between spaces of automorphic forms, and the relevant arithmetic data live in lattices and determinants. The proper generalization of~\eqref{eq:integer-product} is therefore a statement about covolumes and slopes.


\begin{figure}[ht]
\centering
\begin{tikzpicture}[x=1.05cm,y=1cm]
\def\alog{2.07944}\def\blog{2.19722}
\draw[->] (0,0)--(7.7,0) node[right] {$\log |N|$};
\fill[blue!30] (0,0.24) rectangle (3.12,0.86);
\fill[orange!32] (3.12,0.24) rectangle (6.42,0.86);
\draw (0,0.24) rectangle (6.42,0.86);
\node at (1.56,0.55) {$3\log2$};
\node at (4.77,0.55) {$2\log3$};
\draw[dashed,black!45] (6.42,0)--(6.42,1.12);
\node[above] at (3.22,1.04) {$72=2^3\cdot3^2$};
\node[below] at (0,0) {$0$};
\node[below] at (6.42,0) {$\log72$};
\end{tikzpicture}
\caption{For $N=72$, the logarithm of its absolute value is exactly the sum of its prime valuations weighted by $\log p$. An upper bound for the total size of $N$ therefore limits how much divisibility can be packed into it.}
\label{fig:valuation-budget}
\end{figure}

\section{Lattice covolumes and slopes}

A Euclidean lattice $\mathcal L$ is a free abelian group of finite rank $r$ embedded discretely in a Euclidean vector space. If $v_1,\ldots,v_r$ is a lattice basis, its covolume is
\[
\operatorname{covol}(\mathcal L)
=\sqrt{\det(\langle v_i,v_j\rangle)_{i,j}}.
\]
A small covolume means that the lattice points are densely packed in the span. The logarithmic quantity
\begin{equation}\label{eq:slope}
\mu(\mathcal L)=-\frac1r\log\operatorname{covol}(\mathcal L)
\end{equation}
for $r>0$ is called its slope. In geometry, degree and slope control when a vector bundle has sections. In arithmetic geometry, similar ideas are applied to lattices equipped with norms at infinite places. The slope $\mu_{\max}(\mathcal L)$ is the maximal slope of its nonzero sublattices, with the appropriate saturation conventions.

\begin{proposition}[An elementary form of the slope inequality]\label{prop:slope}
Let $\mathcal L$ be a \emph{nonzero} Euclidean lattice, and let $\varphi:\mathcal L\hookrightarrow\mathcal M$ be an injective homomorphism of Euclidean lattices, extended to a linear map of their real spans. Then
\begin{equation}\label{eq:basic-slope}
\mu(\mathcal L)\le\mu_{\max}(\mathcal M)+\log\|\varphi\|,
\end{equation}
where $\|\varphi\|$ is its operator norm.
\end{proposition}
\begin{proof}
Let $r=\operatorname{rk}\mathcal L$. For any lattice basis $v_1,\ldots,v_r$, the Euclidean volume of the parallelepiped spanned by $\varphi(v_1),\ldots,\varphi(v_r)$ is at most $\|\varphi\|^r$ times the original volume. Consequently
\[
\operatorname{covol}(\varphi(\mathcal L))
\le\|\varphi\|^r\operatorname{covol}(\mathcal L).
\]
After taking logarithms and dividing by $r$,
\[
\mu(\mathcal L)
\le\mu(\varphi(\mathcal L))+\log\|\varphi\|.
\]
Since $\varphi(\mathcal L)$ is a nonzero sublattice of $\mathcal M$, its slope is at most $\mu_{\max}(\mathcal M)$. This yields~\eqref{eq:basic-slope}.
\end{proof}

The proposition expresses the basic idea behind Bost's slope inequality in a transparent finite-dimensional setting \parencite{Bost2001}. In an arithmetic application the map and its norms must be compatible with the integral structures at all finite primes; additional local terms may be necessary when rational rather than integral maps are used. The Lean archive has a formalized slope-theory development, including a theorem named \texttt{slopeInequality\_of}. One should not replace that formalized theorem by Proposition~\ref{prop:slope}; the latter merely shows why an inequality of this shape is natural.

\section{A matrix example with visible divisibility}

Let
\[
A=\begin{pmatrix}p^a&0\\0&p^b\end{pmatrix}
\]
be an integral map on $\ZZ^2$, where $a,b\ge0$. Its determinant has valuation $a+b$, so the image lattice has index $p^{a+b}$. If the map is viewed as a real operator, however, its norm is $p^{\max(a,b)}$. The determinant measures a \emph{sum} of divisibilities, while the operator norm measures a \emph{maximum} archimedean expansion. This is already the seed of the arithmetic tension.

A more interesting matrix will not be diagonal. Smith normal form then recovers the elementary divisors of its cokernel: if $A$ is a full-rank integral matrix, there are unimodular $U,V$ such that
\[
UAV=\operatorname{diag}(d_1,\ldots,d_r),
\qquad d_1\mid d_2\mid\cdots\mid d_r.
\]
The valuations of minors control the sums of the valuations of the $d_i$. Thus divisibility of minors is a natural intermediate statement when the arithmetic construction is an operator rather than a scalar.

The announcement and its formal source family single out a \emph{two-by-two minor} arising in the local data of a theta lift. The arithmetic meaning is that some local rows become congruent modulo a uniformizer. For a toy model, if
\[
a_2\equiv a_1\pmod{\varpi},\qquad
b_2\equiv b_1\pmod{\varpi},
\]
then the determinant $a_1b_2-a_2b_1$ is divisible by $\varpi$. If each successive lift produces a controlled instance of such congruence, one may hope to accumulate higher valuations in a determinant built from the whole chain.

This is a \emph{model} of the source's local divisibility step, not a claim that the full local theorem reduces to the two congruences written here. The genuine theorem has specific hypotheses involving local unitary and theta data. Its importance lies in connecting representation-theoretic iteration to the elementary valuation mechanism.

\section{Frames of evaluation maps}\label{sec:frame}

The inequality for one lattice map, proved earlier in the chapter, compares its finite divisibility with the size of its values at the infinite places.  The construction eventually asks for something larger: many evaluation maps must be combined into a comparison for one common target lattice.  A single source may have much smaller rank than that target.  There is no contradiction here.  A short list of vectors, taken one at a time, spans only a line; several lines in sufficiently different directions can fill a plane.  What must be measured is how evenly the collection covers the target.

Start with a concrete example.  Let $H=\RR^2$, and take four maps $B_i:\RR\to H$ whose values at $1$ are the unit vectors
\[
(1,0),\qquad(0,1),\qquad(-1,0),\qquad(0,-1).
\]
Each $B_i$ has one-dimensional image, so no one of them can fill $H$.  Yet the four together span $H$, and
\[
\sum_{i=1}^4 B_iB_i^*=2I_H.
\]
Here $B_i^*$ denotes the adjoint map.  The operator on the left counts the directions provided by the source family.  It is a scalar multiple of the identity because the four sources distribute their contribution equally between the two coordinate directions.  Discard any one map; the other three still span $H$.

\begin{figure}[ht]
\centering
\begin{tikzpicture}[x=.72cm,y=.72cm,>=Latex]
\coordinate (O) at (0,0);
\draw[->,gray] (-3.0,0)--(3.0,0) node[right] {$H$};
\draw[->,gray] (0,-2.6)--(0,2.6);
\draw[accent,very thick,->] (O)--(2.0,0) node[below,font=\scriptsize] {$B_1(1)$};
\draw[accent,very thick,->] (O)--(0,1.95) node[right,font=\scriptsize] {$B_2(1)$};
\draw[accent,very thick,->] (O)--(-2.0,0) node[below,font=\scriptsize] {$B_3(1)$};
\draw[accent,very thick,->] (O)--(0,-1.95) node[right,font=\scriptsize] {$B_4(1)$};
\node[font=\small,align=left] at (7.0,1.0) {Each source has\\rank one.};
\node[font=\small,align=left] at (7.0,-.5) {Together the images\\fill the rank-two target.};
\draw[->,gray!60] (4.55,.78)--(1.65,.25);
\end{tikzpicture}
\caption{A perfectly balanced finite frame. The images are individually one-dimensional, while their Gram sum is $2I$.  The high-dimensional arithmetic construction replaces exact balance by a quantitative approximation.}
\label{fig:four-vectors}
\end{figure}

For general finite-dimensional complex Hermitian spaces $V_i$ and $H$, let $B_i:V_i\to H$ and choose positive weights $c_i$.  The operator
\begin{equation}\label{eq:frame-Q}
Q=\sum_i c_i B_iB_i^*:H\longrightarrow H
\end{equation}
is positive semidefinite.  If $Q$ were a positive multiple of the identity, the family would illuminate every direction in $H$ equally.  In the real construction one cannot reasonably expect exact equality.  Instead, the source proves that the trace of $Q^2$ is close to the minimum compatible with the trace of $Q$.

If $\lambda_1,\ldots,\lambda_D$ are the eigenvalues of $Q$, Cauchy--Schwarz gives
\[
\sum_{j=1}^D\lambda_j^2\ge\frac1D\left(\sum_{j=1}^D\lambda_j\right)^2,
\]
with equality precisely when all eigenvalues are equal.  Thus a bound
\begin{equation}\label{eq:frame-flatness}
\operatorname{Tr}(Q^2)\le\frac{1+\eta}{D}\bigl(\operatorname{Tr}Q\bigr)^2
\end{equation}
asserts that the frame is nearly flat.  The trace-flatness estimate limits how much of the target image can lie in a small collection of directions.  Consequently, removing a low-dimensional portion of the sources cannot concentrate nearly all the remaining mass in a small target subspace.

\section{Frame bounds for restricted source spaces}\label{sec:frame-lemma}

The source spaces come with distinguished subspaces $N_i\subset V_i$.  Later they will be the portions of arithmetic lattices whose slopes are sufficiently large.  We need to know how much of $H$ is spanned by their images.  The next proposition is a readable form of the finite-dimensional estimate isolated in the Lean file
\sourcefile{src/Formalism/HarderNarasimhan/HarderNarasimhanProofs/FrameLinearAlgebraFirst.lean}{FrameLinearAlgebraFirst.lean}.

\begin{proposition}[Almost-spanning subspaces]\label{prop:frame-cover}
For finite-dimensional Hermitian spaces $V_i,H$, maps $B_i:V_i\to H$ and positive weights $c_i$, define $Q$ by~\eqref{eq:frame-Q}.  Let $R=\sum_i\dim V_i>0$ and $D=\dim H>0$.  Suppose $0\le e\le1/4$, $\eta\ge0$, $\varepsilon\ge0$, and
\begin{equation}\label{eq:frame-assumptions}
\sum_i\|c_iB_i^*B_i-I_{V_i}\|_{\mathrm{HS}}^2\le e^2R,
\qquad
\operatorname{Tr}(Q^2)\le\frac{1+\eta}{D}(\operatorname{Tr}Q)^2.
\end{equation}
If $N_i\subset V_i$ have total codimension at most $\varepsilon R$, then
\begin{equation}\label{eq:frame-codim}
\operatorname{codim}_H\left(\sum_iB_iN_i\right)
\le(\eta+4\varepsilon+4e)D.
\end{equation}
\end{proposition}
\begin{proof}
Let $S=\sum_i B_iN_i$, let $P$ be the orthogonal projection onto $S^\perp$, and write $P_i$ for the projection onto $N_i^\perp$.  Since $PB_i$ vanishes on $N_i$, we have $PB_i=PB_iP_i$.  Consequently
\[
\operatorname{Tr}(PQ)
=\sum_i\operatorname{Tr}(PB_ic_iB_i^*)
\le\sum_i\operatorname{Tr}(P_ic_iB_i^*B_i).
\]
Write $c_iB_i^*B_i=I+\Delta_i$.  The first sum of traces is at most $\varepsilon R$.  Hilbert--Schmidt Cauchy--Schwarz bounds the errors by
\[
\sum_i |\operatorname{Tr}(P_i\Delta_i)|
\le\Bigl(\sum_i\|P_i\|_{\mathrm{HS}}^2\Bigr)^{1/2}
\Bigl(\sum_i\|\Delta_i\|_{\mathrm{HS}}^2\Bigr)^{1/2}\le eR.
\]
Applying the same estimate with every $P_i=I$ shows $\operatorname{Tr}Q\ge(1-e)R$.  Therefore, if $r=\operatorname{Tr}(PQ)/\operatorname{Tr}Q$, then $r\le 2(\varepsilon+e)$.  In particular $r\ge0$.

On the other hand, $\operatorname{Tr}((I-P)Q)=(1-r)\operatorname{Tr}Q$.  Cauchy--Schwarz and~\eqref{eq:frame-flatness} give
\[
(1-r)^2(\operatorname{Tr}Q)^2
\le(\dim S)\operatorname{Tr}(Q^2)
\le\frac{1+\eta}{D}(\dim S)(\operatorname{Tr}Q)^2.
\]
Cancelling the nonzero trace yields $\dim S/D\ge(1-r)^2/(1+\eta)\ge1-2r-\eta$.  With the bound on $r$, this is~\eqref{eq:frame-codim}.
\end{proof}

If the original family is nearly balanced and only a small fraction of its source directions is discarded, the resulting target images retain a quantitatively large combined rank. The rank-defect bound controls this \emph{pooled} loss; it does not assert that any individual $B_i$ is surjective, or that every $N_i$ retains the same proportion of its original rank. The distinction is important when the later lattice estimates are summed over the family.

\section{Integrating the slope-distribution bound}\label{sec:slope-integration}

The frame estimate applies to the arithmetic slopes introduced earlier.  An Hermitian lattice has a finite list of Harder--Narasimhan slopes, counted with multiplicity \parencite{HarderNarasimhan1975}.  Thinking of them as a sorted list is sufficient for the calculation that follows.  Integral maps of controlled height carry the high-slope part of the source into a slightly lower high-slope part of the target.  This is the place where the maps $B_i$ of the previous section become useful.

For each source lattice $E_i$, subtract the height bound $h_i$ from its slopes and put all the shifted slopes into one list, weighted automatically by the source ranks.  Let their mean be $s$.  Assume no shifted source slope exceeds $s+u$, with $u>0$.  For a threshold parameter $\tau\ge0$, what fraction of these slopes can be smaller than $s-\tau$?  If a proportion $p$ were below that threshold, while everything else were at most $s+u$, the mean would be bounded above by
\[
(1-p)(s+u)+p(s-\tau)=s+u-p(u+\tau).
\]
Since the mean is $s$, we must have
\begin{equation}\label{eq:slope-tail}
p\le\frac{u}{u+\tau}.
\end{equation}
The inequality is elementary; its force comes from applying Proposition~\ref{prop:frame-cover} to the high-slope portions of the sources together.

The frame estimate says that, apart from a controlled exceptional fraction, the target also lies above the slope threshold $s-\tau$.  Writing $\eta,e$ for the frame errors, let $D$ be the lifted frame-space rank in Proposition~\ref{prop:frame-cover}, and let $d_F$ be the arithmetic target rank whose slopes are being averaged. In the source, the comparison is between $\dim(\kappa\oplus\kappa_1)=D$ and $\dim\kappa=d_F$; it is not a measurement of a physical discard operation on the original target.  If $d_F\ge(1-\delta_F)D$ with $0\le\delta_F<1$, the relevant fraction is normalized by $d_F$, not by $D$. The schematic but quantitatively faithful bound is
\begin{equation}\label{eq:target-tail}
\frac{\#\{\text{surviving target slopes}<s-\tau\}}{d_F}
\le\frac{\eta+4e+4u/(u+\tau)}{1-\delta_F}.
\end{equation}
The numerator comes from the frame argument.  The denominator records the separate loss incurred when passing from ambient rank to surviving rank; the factor $(1-\delta_F)^{-1}$ must not be inserted a second time. These are different errors and should not be combined before the estimate is applied.

\begin{figure}[ht]
\centering
\begin{tikzpicture}[x=1cm,y=1cm,>=Latex]
\draw[->] (.6,.5)--(10.6,.5) node[right] {$\tau$};
\draw[->] (.6,.5)--(.6,4.05) node[above] {fraction below $s-\tau$};
\fill[blue!8] (1,.5) -- plot[domain=1:9.6,samples=60] (\x,{.5+3.0/(.72*\x+.35)}) --(9.6,.5)--cycle;
\draw[accent,thick,smooth,domain=1:9.6,samples=80] plot (\x,{.5+3.0/(.72*\x+.35)});
\draw[dashed] (9.65,.5)--(9.65,2.15);
\node[below,font=\scriptsize] at (9.65,.5) {$H_0$};
\node[font=\small,align=center] at (6.0,2.9) {area under the tail bound\\becomes a mean-slope loss};
\node[font=\scriptsize,rotate=-14,text=accent] at (6.5,1.2) {$u/(u+\tau)$};
\end{tikzpicture}
\caption{Schematic plot of only the $\tau$-dependent component $u/(u+\tau)$, not the full target-tail fraction. The latter also includes frame and rank-loss terms in~\eqref{eq:target-tail} (and is of course at most $1$). Integrating its variable component gives $u\log(1+H_0/u)$.}
\label{fig:slope-tail}
\end{figure}

Assume $H_0\ge0$ and that the surviving target slopes are all at least $s-H_0$. For each such slope $\lambda$,
\[
\max(s-\lambda,0)=\int_0^{H_0}\mathbf1_{\{\lambda<s-\tau\}}\,d\tau.
\]
Averaging over the $d_F$ surviving target slopes and inserting~\eqref{eq:target-tail} shows that their mean is at least
\begin{equation}\label{eq:mean-slope-main}
s-\frac{4u\log(1+H_0/u)+(\eta+4e)H_0}{1-\delta_F}.
\end{equation}
The logarithm has an elementary origin:
\[
\int_0^{H_0}\frac{u}{u+\tau}\,d\tau=u\log\left(1+\frac{H_0}{u}\right).
\]
At $u=0$ the expression is interpreted by its limit, namely zero.  The logarithmic term arises by integrating the permissible loss of source rank across the slope thresholds, with the normalization shown in the preceding formula.  The endpoint value follows from the same integral.

The actual arithmetic theorem, \texttt{theoremFSharp}, bounds the normalized target value \texttt{muDef b' X' cc}, not an unshifted geometric slope. It uses a \emph{doubled} Hermitian/Arakelov slope-and-height normalization (including a slope-normalizing constant), rather than literally the Euclidean convention $\mu=-\log(\operatorname{covol})/r$ used in the introductory calculation. Factors of $2$ and the height shifts must therefore be translated consistently, as detailed in Appendix~\ref{app:slope-technical}. The source theorem also tracks the coefficient field, complex embeddings, integrality of the maps, the shift in slope caused by operator height, and positivity of all participating ranks.  Those matters are essential when the formula is used, and are summarized in Appendix~\ref{app:slope-technical}.  The calculation above explains the finite-dimensional and distributional heart of the theorem, which is what the main argument needs from it.

\section{Valuations and centered prime sums}

The field used in the announced construction contains a quadratic imaginary field $K=\QQ(\sqrt{-d})$. Let $\chi_{-d}$ be its quadratic Dirichlet character. For primes $p$ away from the discriminant,
\[
\chi_{-d}(p)=\begin{cases}
+1,&p\text{ splits in }K,\\
-1,&p\text{ is inert in }K.
\end{cases}
\]
The factor $1-\chi_{-d}(p)$ therefore distinguishes inert primes from split primes. At an inert prime it equals $2$; at a split prime it vanishes. For prime powers, the value depends on the parity of the exponent, because $\chi_{-d}(p^k)=\chi_{-d}(p)^k$.

A provisional mathematical guide to the archive describes a divisibility statistic, denoted $\operatorname{ASK}(a)$, of the form
\begin{equation}\label{eq:ask}
\operatorname{ASK}(a)
=\sum_{n\ge1}\Lambda(n)(1-\chi_{-d}(n))
\left(\frac an-w_n\right)_+,
\end{equation}
where the positive weights $w_n$ are associated with Frobenius classes in a finite abelian Galois group. This is an illustrative unramified shorthand, not the exact source definition: the source functional $\operatorname{ASK}(D,N,a,t,w)$ also depends on a chosen translate $t$. Its ramified-prime contributions use the Frobenius cosets $rF_p+I_p$, an index-two coset restriction, and an inertia normalization $2/|I_p|$, as made explicit in~\eqref{eq:finite-gain-A}. In particular, the prime $d$ must not be treated as if its quadratic character always took a value $\pm1$. The expression is directly comparable with the ramp in Chapter~\ref{ch:statistic}: for any $w>0$, putting $y=a/w$ gives
\begin{equation}\label{eq:ramp-rescale}
\left(\frac an-w\right)_+
=w\left(\frac yn-1\right)_+.
\end{equation}
The valuation identity obtained from the theta-lift divisibility estimate can therefore be rewritten as a weighted combination of centered prime sums. Its normalization is the arithmetic input needed for the subsequent Fourier character projection and for the analytic bound on $G_F$.



\section{Periods and the archimedean estimate}

The slope inequality bounds the total divisibility in terms of norms. But an abstract operator-norm bound of the form $\|\Theta^{(m)}\|\le C^m$ would almost certainly be too crude; it may cost more than the divisibility gained. The argument needs the normalization of periods, masses, and local factors to high precision.

The author's announcement names the Chowla--Selberg formula for complex-multiplication periods, Siegel's mass formula, Weyl's formula, and the Rallis inner-product formula. Each has a distinct role. The period formula relates transcendental factors to arithmetic $L$-values. The mass formula computes the aggregate weights of definite lattice classes. Weyl's formula controls representation-theoretic dimensions. The Rallis formula relates norms of theta lifts to special values of $L$-functions. The seesaw identity helps express the period of the iterated lift in a form suitable for averaging.

The same guide gives an archimedean comparison term of the form
\begin{equation}\label{eq:main-cost}
\operatorname{MAIN}(a)
=a\left(\log a-1-\gamma+\frac{L'}{L}(1,\chi_{-d})
-\overline{\log w}\right).
\end{equation}
In this formula $\overline{\log w}:=\frac1f\sum_{\rho\in\Phi}\log w(\rho)$, where $\Phi$ is the index-two subgroup of the finite Galois group used in the later Fourier construction, $f=|\Phi|$, and $w:\Phi\to\RR_{>0}$ is the positive weight function; see Chapter~\ref{ch:fourier}, especially~\eqref{eq:arith-main-M}. The terms are intelligible individually. The $a\log a-a$ behavior is reminiscent of the logarithm of a factorial, as one would expect from a construction playing the role of repeated differentiation. The logarithmic derivative $L'/L(1,\chi_{-d})$ arises from arithmetic normalization of CM periods. The average of $\log w$ records the chosen family of Galois weights. That these contributions combine in precisely this way is a deep identity, not something we infer from their superficial resemblance.

Alp\"oge writes that the initial cost calculation would ``stop at 2/3'' without additional arithmetic improvements.\footnote{Alp\"oge, \emph{Remark}, p.~4, item~14; his qualification is \emph{a priori}.} The difference from $1/2$ is decisive. An exponent of the form $1/2+C\eps$ can be brought arbitrarily close to $1/2$ by choosing $\eps$ small, whereas $2/3$ leaves a fixed gap. His timeline records several proposed ways past that obstruction---isolating a main term, transferring a constant to a function-field comparison, or using twists---before turning to a common integral structure across Hecke characters, from which a gcd-like averaging operation could gain additional divisibility.\footnote{Alp\"oge, \emph{Remark}, p.~4, items~14--15. These are proposed remedies described in his timeline, not a derivation of the final exponent.} This gives a concrete reason to care about accurate period normalizations and integral structures: at this point a small-looking arithmetic refinement must eliminate a genuine power of $x$. The calculation leading to $2/3$ is reported from his remark, not independently reconstructed in this chapter.

The desired comparison between~\eqref{eq:ask} and~\eqref{eq:main-cost} is intended to give an upper bound for a centered prime statistic with exponent arbitrarily close to $1/2$. The elementary identities in this chapter explain why the inequality has a plausible arithmetic shape: local congruences supply valuations, while norm calculations restrict their total size. The exact numerical comparison depends on the period identities, the integral-structure constructions, and the local congruence estimates. Those deep steps are not independently rederived here.

\part{Arithmetic fields and local-global estimates}

\chapter{Auxiliary number fields and splitting conditions}\label{ch:hosts}

\section{Required properties of the auxiliary field}

For each selected scale and target Dirichlet character, the construction introduces an auxiliary number field $E$ with controlled splitting and a character quotient suited to Fourier projection. The field may vary with the scale, while the chosen analytic character family must remain fixed across the windows to which the zero-free criterion is applied. The auxiliary field makes the divisibility estimate selective enough to isolate that family.

Choosing an auxiliary field determines how rational primes split, which Galois characters are available, and how congruences can be organized. Several requirements have to be met at once: certain small primes should split, certain characters should appear as Fourier modes of its Galois group, and its cyclotomic part should supply enough room for an Iwasawa-theoretic window.

The source defines a predicate \texttt{IsChiField} governing eligible host fields. The public architectural guide gives an instructive representative expression,
\begin{equation}\label{eq:host}
E\simeq\QQ(\zeta_{2^{e+2}})^+\,K_\psi\,
\QQ(\sqrt{-q},\sqrt{-d}),
\end{equation}
where $K_\psi$ is the extension cut out by the chosen character $\psi$, the superscript $+$ denotes a maximal real cyclotomic subfield, and $q,d$ are auxiliary primes. The precise compositum, embeddings and fixed-field conditions are part of \texttt{IsChiField};~\eqref{eq:host} is a mathematical description of the intended structure rather than a replacement for that predicate.

\section{Cyclotomic levels as functions of scale}

Let $x$ be the scale of the prime sum and let $\eps>0$ be small. The file
\sourcefile{src/Development/AutomorphicArithmetic/CharacterBoxes/ExistsPageN.lean}{ExistsPageN.lean}
introduces a simple but revealing choice:
\begin{equation}\label{eq:page}
x^{2+8\eps}\le2^e\le2x^{2+8\eps}.
\end{equation}
It is always possible to find such an integer $e$ when $x\ge1$. One simply takes the next sufficiently large power of $2$. The Lean theorem \texttt{exists\_pageN} proves that this can be done uniformly as a choice function of $x$.

Taking logarithms in~\eqref{eq:page}, one finds
\[
e=(2+8\eps)\log_2 x+O(1).
\]
Thus $e$ is only logarithmic in $x$, even though $2^e$ is a large power of $x$. This is why a $2$-power cyclotomic tower is a flexible choice: each increase of $e$ by one doubles the modulus, permitting a controlled adjustment of field size while retaining a nested arithmetic structure.

\begin{figure}[ht]
\centering
\begin{tikzpicture}[>=Latex,node distance=8mm, every node/.style={font=\small,align=center},bx/.style={draw=black!30,rounded corners=2pt,minimum width=48mm,minimum height=9mm}]
\node[bx] (a) {$\QQ$};
\node[bx,above=of a] (b) {$\QQ(\zeta_8)^+$};
\node[bx,above=of b] (c) {$\QQ(\zeta_{16})^+$};
\node[bx,above=of c] (d) {$\QQ(\zeta_{32})^+$};
\node[above=6mm of d] (more) {$\vdots$};
\draw[->] (a)--(b);\draw[->] (b)--(c);\draw[->] (c)--(d);\draw[->] (d)--(more);
\node[right=13mm of c,align=left] {\footnotesize increasing $e$ changes\\[-2pt]\footnotesize cyclotomic conductor\\[-2pt]\footnotesize and available characters};
\end{tikzpicture}
\caption{The real cyclotomic tower. The exact field used at a given scale also incorporates the character field $K_\psi$ and two quadratic fields.}
\end{figure}

Alp\"oge connects his first Iwasawa-tower idea to a ``pigeonhole'' analogy over $\F_{q^n}/\F_q$, and later singles out Behrend's construction in the more precise tower choice.\footnote{Alp\"oge, \emph{Remark}, p.~4, items~11 and~19. The analogy records his mathematical motivation; the remark does not establish the source's field-existence or window hypotheses.} The tower is chosen to align the available arithmetic with the range and the character being tested.

A cyclotomic tower provides a nested sequence of fields whose Galois groups and ramification can be described explicitly. That organization is useful when the same cohomological calculation must be compared at many levels. It does not confer the elementary Frobenius identities available over finite fields; those are replaced by the much more elaborate Galois and automorphic comparisons developed elsewhere in the argument.

\section{Auxiliary-prime congruences and splitting}

The host also depends on auxiliary primes $q$ and $d$. The condition \texttt{WinPair} in \texttt{ExistsPageN.lean} includes
\begin{equation}\label{eq:auxprimes}
q\equiv3\pmod8,\qquad d\equiv7\pmod{16},\qquad q\mid d+1,
\end{equation}
as well as asymmetric quantitative ranges for $\log q$ and $\log d$ and a condition that every odd prime up to a specified cutoff divides $d+1$. The logarithmic size scales differ by a factor of $(\log T)^2$; the exact inequalities at scale $T$ appear in~\eqref{eq:admissible-prime-sizes}. These congruences have an immediate arithmetic interpretation.

Set $K=\QQ(\sqrt{-d})$, with $d\equiv7\pmod{16}$ prime. Then $-d\equiv1\pmod8$. The quadratic discriminant is $-d$; hence the prime $2$ splits in $K$. If $q$ is a distinct odd prime dividing $d+1$, then $-d\equiv1\pmod q$, so the Kronecker symbol satisfies
\[
\left(\frac{-d}{q}\right)=1,
\]
and $q$ splits in $K$. The same reasoning applies to every odd prime $p$ dividing $d+1$ and distinct from $d$:
\[
-d\equiv1\pmod p\quad\Longrightarrow\quad
\left(\frac{-d}{p}\right)=1.
\]
Thus the congruence conditions ensure that a prescribed collection of small primes splits in $K$.

Why arrange for small primes to split? Recall the divisibility statistic~\eqref{eq:ask}. It contains the factor $1-\chi_{-d}(n)$. At split primes this factor vanishes for every positive power. Such primes therefore contribute nothing to the principal inert-prime sum. By forcing the troublesome small primes to split, the construction suppresses terms that are costly to treat uniformly in the later estimates.

\begin{figure}[ht]
\centering
\resizebox{\linewidth}{!}{%
\begin{tikzpicture}[>=Latex, every node/.style={font=\small,align=center}, bx/.style={draw=black!30,rounded corners=2pt,minimum width=29mm,minimum height=12mm}]
\node[bx] (p) {Prime $p$\\\scriptsize $p\mid d+1$};
\node[bx,right=8mm of p] (s) {$-d\equiv1\pmod p$};
\node[bx,right=8mm of s,fill=blue!5] (k) {$p$ splits in $K$\\\scriptsize $\chi_{-d}(p)=1$};
\draw[->] (p)--(s);\draw[->] (s)--node[above,font=\scriptsize]{odd $p$}(k);
\node[below=9mm of s] (z) {$1-\chi_{-d}(p^j)=0$ for every $j\ge1$};
\draw[->,gray] (k.south) to[bend left=14] (z.east);
\end{tikzpicture}%
}
\caption{For odd $p\mid d+1$ with $p\ne d$, the Legendre-symbol criterion gives splitting in $K$. The prime $2$ is a separate case: it splits because the standing hypothesis $d\equiv7\pmod{16}$ implies $-d\equiv1\pmod8$. In either case the inert-prime factor vanishes.}
\end{figure}

The quantitative existence of primes $q,d$ satisfying all the simultaneous size, congruence and divisibility conditions is not an elementary consequence of Dirichlet's theorem in the particular ranges used here. In the source it is packaged into further arithmetic lemmas. We have verified the easy \emph{consequence} of these congruences; the deep point is obtaining them at the right scale without losing the desired square-root estimate.

\section{Embedding the target character in the host field}

A Dirichlet character $\psi$ can be viewed as a one-dimensional representation of a finite abelian Galois group. The associated field $K_\psi$ makes this correspondence concrete: it is the fixed field of the kernel of the character. If $\psi$ has order $o$, then $\operatorname{Gal}(K_\psi/\QQ)$ is cyclic of order $o$. Its characters are exactly the powers $1,\psi,\ldots,\psi^{o-1}$.

That fact will become decisive in the Fourier-averaging step. The theta-lift inequality is naturally expressed in terms of Frobenius classes of the large host field $E$, so it initially contains more character information than is wanted. By ensuring that the desired $\psi$ occurs as a character of the host's Galois group, one can average away unwanted frequencies while retaining the powers of $\psi$.

The subgroup $K_\psi$ is chosen so that $\psi$ descends to a character of the quotient. Averaging over $K_\psi$ then projects the host-character expansion onto the characters compatible with that quotient, including the desired powers of $\psi$. The cyclotomic and cohomological bounds control the arithmetic cost of making this projection inside the auxiliary field.

\section{Long windows and variable auxiliary fields}

The archive does not insist that a suitable host exist uniformly at every sufficiently large $x$ with the same auxiliary data. The predicate \texttt{VWinAt} permits the scale to move into a carefully chosen interval and the host to be rebuilt there. This flexibility is essential because the analytic theorem in Chapter~\ref{ch:statistic} requires only arbitrarily late windows.

There is an important alignment between the quantifiers. After the character $\psi$, small parameter $\eps$ and admissible level selector $e(x)$ are fixed, the cohomological window interface selects global parameters $J,X_0$. For each sufficiently large $X$ and each admissible pair $(q,d)$, \emph{one} interval $[Y,Y^2]$ is selected. The pair $q,d$ stays fixed on that entire interval; the level $e(x)$ and eligible host field may vary with $x$. Formally the estimate is conditional on the host predicate, and the actual host is supplied by separate construction results. In the pinned development, \texttt{PairEx} supplies admissible prime pairs satisfying \texttt{PairGood}; the bridge \texttt{supplyChi} uses \texttt{VWinAt $\psi$ 401} together with the admissible host construction to reach \texttt{HostExW $(3/2)$}. These are distinct producers: \texttt{VWinAt} alone neither constructs $(q,d)$ nor asserts existence of an eligible host field.\footnote{See \sourcefile{src/Development/AutomorphicArithmetic/CharacterBoxes/PairGood.lean}{PairGood.lean}, \sourcefile{src/Development/AutomorphicArithmetic/CharacterBoxes/SupplyChi.lean}{SupplyChi.lean}, and \sourcefile{src/Development/AutomorphicArithmetic/CharacterBoxes/ExistsPageN.lean}{ExistsPageN.lean}.} The estimate is \emph{uniform at every scale $x$ within the selected interval} and for every eligible host at that scale. What is \emph{not} asserted is a corresponding estimate at every sufficiently large scale \emph{outside} these selected windows. The analytic theorem needs only arbitrarily late good windows, so this precise distinction is essential; see Theorem~\ref{thm:quantified-window}.

\chapter{Automorphic and Galois representations}\label{ch:galois}

\section{Automorphic and Galois local parameters}

An automorphic representation has local Hecke eigenvalues at almost every prime. An $\ell$-adic Galois representation has Frobenius eigenvalues at almost every prime. In the cases relevant here, the Langlands correspondence relates these two sets of data by an equality of characteristic polynomials, after a specified normalization. For automorphic Galois constructions in suitable settings, see \textcite{ChenevierHarris2013}; the local hypotheses for the present construction are recorded separately below. This permits information produced by automorphic constructions to be read in the Galois language needed for local-global arithmetic.

At an unramified prime the local compatibility has a concrete description. An automorphic representation of $\operatorname{GL}_n$ has local parameters $\alpha_1,\ldots,\alpha_n$ that determine its local Euler factor. A Galois representation
\begin{equation}\label{eq:rhopi}
r_\pi:\Gal(\overline E/E)\longrightarrow
\operatorname{GL}_n(\overline\QQ_\ell)
\end{equation}
associates to a Frobenius element at an unramified prime $v$ a matrix whose eigenvalues provide a second description of local arithmetic behavior. The comparison is not merely a coincidence of dimensions; its content is a precise identity of characteristic polynomials, after specified normalizations.

At the level of a schematic unramified local correspondence, one encounters
\begin{equation}\label{eq:satake-frob}
\iota\,\operatorname{charpoly}\bigl(r_\pi(\operatorname{Frob}_v^{-1})\bigr)
=\prod_{j=1}^n\left(T-q_v^{(n-1)/2}\alpha_j\right),
\end{equation}
where $q_v$ is the residue-field cardinality, the $\alpha_j\in\CC^\times$ are nonzero Satake parameters, and $\iota:\overline\QQ_\ell\simeq\CC$ identifies the coefficient fields. Here $\operatorname{Frob}_v^{-1}$ denotes the inverse Weil/Frobenius element used in the source comparison: its local diagonal parameters are first written as $\alpha_j^{-1}q_v^{(1-n)/2}$ and then inverted, giving $\alpha_jq_v^{(n-1)/2}$. The exact form of the normalization depends on the chosen Frobenius and reciprocity conventions;~\eqref{eq:satake-frob} records the inverse-element convention traced in \texttt{RecAlgNormDiag.lean}. It is important not to erase that normalization when comparing local factors or estimating their sizes.

The familiar case $n=1$ gives the right intuition. A Dirichlet character is simultaneously a multiplicative function on residue classes and, via abelian class field theory, a one-dimensional character of a Galois group. Its value $\chi(p)$ at an unramified prime is a Frobenius eigenvalue. Higher-dimensional Langlands correspondence can be seen as the vastly more sophisticated generalization of this equivalence of viewpoints.

Local compatibility between the automorphic and Galois descriptions supplies the Galois representation data used to define the coefficient modules and restriction maps of the cohomological obstruction groups in the next chapter. Matching local characteristic polynomials contributes to this comparison under the stated representation hypotheses; the one-dimensional example illustrates that relation without constructing the general Galois representation. The local-global bounds require the additional source-level inputs described there.

\section{Regular algebraicity and polarization}

Not every automorphic representation is expected to come from an $\ell$-adic Galois representation with the desired properties. The source relies on a range of hypotheses familiar from the theory of regular algebraic, essentially conjugate self-dual automorphic representations over CM fields.

The representation is \emph{regular algebraic} when its behavior at the infinite places is compatible with an algebraic representation whose weights are distinct. The distinctness helps prevent different cohomological degrees and filtrations from becoming entangled. It is \emph{polarized}, roughly, when conjugation and duality are related by a controlled character twist. One schematic pattern (not the literal \texttt{RowAt} contract) is
\begin{equation}\label{eq:polar}
\pi^c\simeq\pi^\vee\otimes(\chi\circ N_{E/E^+}).
\end{equation}
Here $c$ denotes complex conjugation on $E$, $\pi^\vee$ is the contragredient representation, $\chi$ is an appropriate character on $E^+$ (under the conventions in which the display is used), and $N_{E/E^+}$ is the norm map. The source states its actual conditions through automorphic, weight, and local-compatibility predicates; the exact parity and twisting conventions cannot be inferred from this schematic equation alone.

Under suitable hypotheses, the Galois representation~\eqref{eq:rhopi} is continuous and semisimple. At primes away from $\ell$ it has local Langlands compatibility; at primes above $\ell$ it satisfies conditions from $p$-adic Hodge theory, such as being de Rham or crystalline in the appropriate situations. These properties matter in the arithmetic comparison with the cohomological estimates in Chapter~\ref{ch:iwasawa}. However, the obstruction groups introduced there initially use cyclotomic Tate-twist coefficients $\QQ_\ell/\ZZ_\ell(-e)$, \emph{not} the $n$-dimensional representation $r_\pi$ as their coefficient module. The automorphic/Galois representations enter the difficult estimates and compatibility constructions rather than directly defining that \texttt{Sha}. Merely knowing an equality of characteristic polynomials at almost every prime would not suffice for all of the local deformation and integral arguments.

The Lean code's \texttt{RowAt} condition in the file
\sourcefile{src/Development/AutomorphicArithmetic/Foundations/GaloisRepresentations/ThetaParseval/SplitLineSubtraction/RowDefOfOdd.lean}{RowDefOfOdd.lean}
combines an existence assertion with additional compatibility at places. More precisely, it first asserts that some representation satisfies \texttt{WeakU}, and then requires the prescribed local compatibility for \emph{every} representation satisfying \texttt{WeakU} and every eligible split place and local parameter data. This dual requirement is significant: constructing $r_\pi$ is only the beginning; to use it in the arithmetic argument, one needs the local refinements.

\section{Subtraction of auxiliary constituents}
A pair of enlarged representations can sometimes reveal an original constituent that neither enlargement would identify by itself. For a first calculation, take polynomials over a field and suppose we are given two products
\[
A(T)=Q(T)R(T),\qquad B(T)=Q(T)S(T),
\]
where $R$ and $S$ describe additional, deliberately introduced information. If $R$ and $S$ are relatively prime, then
\begin{equation}\label{eq:gcd}
\gcd(A,B)=Q
\end{equation}
up to multiplication by a unit, provided the greatest common divisors are chosen monic and $Q$ is monic. The extra factors are not an error; they are a device for making $Q$ accessible through comparison.

For instance, suppose
\[
Q(T)=(T-2)(T-5),\quad
R(T)=T-7,\quad S(T)=T-11.
\]
Then $A=Q(T)(T-7)$ and $B=Q(T)(T-11)$. We see $Q$ as the common factor of two enlargements, even if neither enlargement was constructed by starting with explicit knowledge of $Q$.

\begin{proposition}[Two-enlargement recovery]\label{prop:gcd}
Let $K$ be a field and let $Q,R,S\in K[T]$ be nonzero. If $\gcd(R,S)=1$, then $\gcd(QR,QS)$ is associated to $Q$.
\end{proposition}
\begin{proof}
Write a B\'ezout identity $UR+VS=1$ with $U,V\in K[T]$. Multiplying by $Q$ gives
\[
U(QR)+V(QS)=Q.
\]
Thus every common divisor of $QR$ and $QS$ divides $Q$, while $Q$ plainly divides both. Their greatest common divisor is therefore associated to $Q$.
\end{proof}

The analogy to representation theory should be used with care. In a representation, the extra data are characters, eigenvalues and local parameters, not just factors in one fixed polynomial. Distinct characters can have the same value at a particular group element, so one cannot naively take the greatest common divisor of two characteristic polynomials at a single point and expect the original representation to emerge. One needs compatible identities across enough group elements or places, together with the appropriate semisimplicity and local--global comparison results. In other words, the algebraic picture gives the correct \emph{shape} of the operation but not its full theorem.

\section{Odd-rank enlargement and parity reduction}
A source file entitled \sourcefile{src/Development/AutomorphicArithmetic/Foundations/GaloisRepresentations/ThetaParseval/SplitLineSubtraction/RowDefOfOdd.lean}{RowDefOfOdd.lean} contains definitions named \texttt{TwoUp}, \texttt{TwoUpAll}, and \texttt{Subtract}, followed by a theorem \texttt{rowDef\_of\_odd}. These names are unusually informative when read together.

The starting object is an automorphic representation attached to a unitary group of rank $n$, supplemented by suitable weight data. The property \texttt{TwoUp} asserts the availability of two enlarged automorphic objects of rank $n+t$, both with a controlled description of their added characters, and with the extra condition
\[
\psi_{1,i}\ne\psi_{2,j}\qquad\text{for every pair }(i,j).
\]
The two collections of auxiliary characters are therefore disjoint. In the precise \texttt{TwoUp} quantifiers, these two enlarged objects may be chosen afresh \emph{for each specified split place}; a single pair is not required to be unramified at every place. The local unramifiedness condition singles out the common summand and the first auxiliary character family. The later subtraction and uniqueness/transport argument, not a globally fixed enlargement pair, produces a representation with the required all-place properties. The target rank $n+t$ is required to be odd. That last condition is not incidental: the next stage assumes a Galois-representation statement available at odd ranks, and subtraction is designed to carry the information back down.

More formally, the source defines a proposition we may abbreviate by $\operatorname{RowAt}(n)$, saying roughly that the appropriate automorphic representation in rank $n$ has a compatible $\ell$-adic Galois representation with specified local properties. The file's \texttt{Subtract} proposition has the schematic logical form
\begin{equation}\label{eq:subtract}
\left(\forall d\text{ odd},\ \operatorname{RowAt}(d)\right)
\land \operatorname{TwoUp}(n)
\quad\Longrightarrow\quad\operatorname{RowAt}(n),
\end{equation}
with the field, splitting, weight, and prime hypotheses restored in the actual definition. In particular, the \texttt{Scope} predicate restricts the CM extension to the required finite unramifiedness and quadratic description, requires the maximal totally real subfield to have even degree, imposes splitting at the places over $\ell$, and requires at least three real-subfield places over $\ell$. These are genuine arithmetic hypotheses, not requirements of the elementary parity argument. The theorem \texttt{rowDef\_of\_odd} then splits into two cases, according as $n$ is odd or even.

For odd $n$, the desired statement is already among the hypotheses. For even $n$, one invokes \texttt{TwoUpAll} to obtain the two enlargements at odd rank, and \texttt{Subtract} to descend. The formal proof of this final parity step is very short; its mathematical force rests entirely in the auxiliary theorems that justify the two constructions and their comparison.

\begin{theorem}[Conditional parity reduction, as organized in the archive]\label{thm:parity}
Assume the three source-level hypotheses: the subtraction principle \texttt{Subtract}, the availability of suitable two-lift constructions \texttt{TwoUpAll}, and the required Galois-representation statement \texttt{RowDef} in every odd positive rank. Then \texttt{RowDef} holds in every positive rank.
\end{theorem}
\begin{proof}
Let $n\ge1$. If $n$ is odd, apply the assumed odd-rank result. Otherwise $n$ is even. The \texttt{TwoUpAll} hypothesis supplies a pair of suitable odd-rank enlargements. Apply the assumed odd-rank result to those enlargements and then invoke \texttt{Subtract} to recover the rank-$n$ assertion. This is exactly the case distinction performed by \texttt{rowDef\_of\_odd} in the pinned Lean source.
\end{proof}

\begin{figure}[ht]
\centering
\begin{tikzpicture}[node distance=5.5mm,>=Latex, every node/.style={align=center,font=\small}, bx/.style={draw=black!30,rounded corners=2pt,inner sep=5pt,text width=5.3cm},ar/.style={-{Latex[length=2mm]},line width=.6pt}]
\node[bx,fill=blue!4] (n) {Even-rank object $U(n)$\\\footnotesize target Galois data unavailable directly};
\node[bx,fill=blue!4,below=8mm of n] (ab) {Two extensions at odd rank $n+t$\\\footnotesize auxiliary character families are disjoint};
\node[bx,fill=blue!4,below=8mm of ab] (odd) {Apply available odd-rank theorem\\\footnotesize obtain larger Galois representations};
\node[bx,fill=blue!8,below=8mm of odd] (res) {Subtract auxiliary constituents\\\footnotesize recover desired data in rank $n$};
\draw[ar] (n)--node[right,font=\scriptsize]{\texttt{TwoUp}}(ab);
\draw[ar] (ab)--(odd);\draw[ar] (odd)--node[right,font=\scriptsize]{\texttt{Subtract}}(res);
\end{tikzpicture}
\caption{The parity reduction of \texttt{RowDefOfOdd.lean}. The two enlargements may depend on the split place; the subsequent subtraction and local-transport argument globalizes the output. This is a conditional transfer, not by itself an existence proof for the lifts.}
\end{figure}

The even-rank case is treated through two odd-rank enlargements whose auxiliary constituents can be distinguished. The odd-rank input applies to each enlargement under its hypotheses; comparing the results and subtracting the known constituents then recovers the desired even-rank information. The separate character families and local choices are essential to this subtraction, and are preserved in the cited source propositions.

\section{Extension of identities from dense subsets}

The transition from Frobenius data to all group elements deserves a short mathematical explanation. Let $\Gamma$ be a topological group and let $f,g:\Gamma\to V$ be continuous maps to a Hausdorff space. If $f=g$ on a dense subset of $\Gamma$, then $f=g$ everywhere. Indeed, the equalizer $\{x:f(x)=g(x)\}$ is closed and contains a dense set.

Galois theory supplies many dense subsets formed from Frobenius elements; the Chebotarev density theorem is one route to such density statements in the appropriate profinite group. Thus a polynomial identity verified at almost all Frobenius elements may, when continuity and the required density hypotheses are in place, extend to every group element. This is the elementary topological heart of the source's \texttt{DenseId} statement.

There is a second issue: even equality of characteristic polynomials at every group element need not identify arbitrary nonsemisimple representations, because different extensions can have the same semisimplification. Semisimplicity is therefore important. This is why the source repeatedly records \texttt{IsSSRep}; it is not a cosmetic condition attached to the theorem name.

\section{Local compatibility in the arithmetic estimate}
The comparison between automorphic and Galois representations is required at the same finite places where the theta lifts are studied and at the places above $\ell$ where the cohomological groups are defined. A global correspondence without precise local compatibility would not give the needed information about Frobenius parameters, de Rham conditions, or congruence lattices.

The source's local-compatibility predicate \texttt{Cl} has two parts: a Frobenius-eigenvalue comparison away from $\ell$ and an \texttt{Above} clause for places above $\ell$. The latter incorporates de Rham/crystalline, weight-jump, and Weil-characteristic-polynomial information \emph{conditional on the specified local bridge data}; it is stronger than one unramified characteristic-polynomial identity. Accordingly, \texttt{RowAt} combines existence of a Galois representation with these conditional local properties. The parity-reduction step in the preceding sections is consequently stronger than a dimension argument: subtracting auxiliary characters must leave the correct local information intact. This is one reason a theorem that is almost obvious for semisimple representations at the level of character identities can become difficult in the arithmetic setting.

\chapter{Galois cohomology and arithmetic windows}\label{ch:iwasawa}

\section{Local restrictions and global cohomology}

Calculations made after completing a number field at each relevant prime are often easier than the corresponding global problem. But a collection of locally solvable conditions need not assemble into a global solution. The discrepancy is measured by various local-to-global obstruction groups. Arithmetic duality provides the wider setting \parencite{Milne2006}; the bounds needed in this chapter are more specialized. In the present development the obstruction groups are expressed through Galois cohomology.

For a number field $E$ and a prime $\ell$, one considers a Galois group $G_{E,\ell}$ governing extensions unramified outside the primes above $\ell$ (and the infinite places, where applicable), together with an $\ell$-adic Galois module $M$. Its first cohomology group $H^1(G_{E,\ell},M)$ parametrizes cocycle classes subject to the specified global conditions. For each place $v\mid\ell$, a restriction map sends a global class to a local class in $H^1(E_v,M)$.

The first cohomology group has a concrete description. A $1$-cocycle is a map $c:G\to M$ satisfying
\[
c(gh)=c(g)+g\cdot c(h).
\]
Two cocycles represent the same cohomology class when their difference has the form $c(g)=g\cdot m-m$ for some $m\in M$. One can restrict such a class from the global Galois group to the decomposition group at each chosen place. It may become a coboundary under every such restriction without being a coboundary globally. That possibility is measured by the kernel
\begin{equation}\label{eq:sha}
\Shasym^1(E,M)
=\ker\!\left(H^1(G_{E,\ell},M)
\longrightarrow\prod_{v\mid\ell}H^1(E_v,M)\right).
\end{equation}
This kernel contains global classes invisible to these local tests. This notation is related in spirit to Tate--Shafarevich groups, although the coefficient modules and local conditions here are specialized and should not be conflated with the Tate--Shafarevich group of an elliptic curve.

The source defines the coefficient module
\[
M_{\ell,e}=\QQ_\ell/\ZZ_\ell(-e),
\]
where the parenthesized twist indicates a power of the cyclotomic character. It denotes the group~\eqref{eq:sha} by \texttt{Sha} and its $\ZZ_\ell$-module length by \texttt{shaLen}. In the mathematical notation of this chapter let
\[
a_{\ell,e}(E)=\operatorname{length}_{\ZZ_\ell}\Shasym^1(E,M_{\ell,e}).
\]
The length is an extended nonnegative integer before finiteness is established. Thus it is important that the source's window property explicitly states finiteness for all indices in the prescribed range. A finite sum of nonnegative numbers is not well defined as a useful ordinary bound until those values are known to be finite.

Alp\"oge reports that twisting did not eliminate all the cohomology classes and that the argument instead sought a substantial weakening of a Vandiver-type vanishing assertion.\footnote{Alp\"oge, \emph{Remark}, p.~4, item~17. His note describes a strategic change, not an independently verified Vandiver-type result.} The chapter's quantitative estimates should be read in that light: a very strong assertion that every obstruction vanishes may be unnecessary if the proof can make their cumulative contribution small enough.

\section{Bounds for sums of cohomological lengths}

Vanishing of every obstruction group in~\eqref{eq:sha} would be welcome, but the argument can afford less. Nonzero groups are tolerable if their combined lengths remain small beside the divisibility gained from theta lifts.

The source predicate \texttt{VVle} captures this quantitatively. Abstracting the selection of local data, its essential content is the bound
\begin{equation}\label{eq:vvle}
\sum_{\ell\in S}(\log\ell)
\sum_{1\le e<t}a_{\ell,e}(E)\le B
\qquad\text{for every finite set of primes }S,
\end{equation}
together with the finiteness of each length in the indicated range. The quantifier over all finite sets $S$ is strong: since every summand is nonnegative, it controls arbitrary partial sums of the total contribution of all primes.

\begin{lemma}[The monotonicity of an obstruction budget]\label{lem:vvmono}
Suppose the nonnegative lengths $a_{\ell,e}$ are finite for $1\le e<t$, and~\eqref{eq:vvle} holds for a bound $B$. If $1\le t'\le t$ and $B'\ge B$, then the same bound holds with $t'$ and $B'$ in place of $t$ and $B$.
\end{lemma}
\begin{proof}
Every summand in~\eqref{eq:vvle} is nonnegative. Reducing the index range can only decrease the sum, and enlarging the allowed upper bound cannot invalidate the inequality. The finiteness assertion remains true on the smaller range.
\end{proof}

The formal lemma \texttt{VVle.mono} relaxes the cutoff and upper budget for a cohomological length bound over the \emph{same} number field. It does not transport such an estimate to another field. Transport across an isomorphism uses \texttt{vvle\_congr}, while use of scale-dependent auxiliary hosts is expressed by the universal conditional quantifiers in \texttt{VAt} and \texttt{VWinAt}. These distinctions allow the cutoffs to be adjusted without changing the host assumptions under which the bound was proved.

\section{The cohomological window condition}

The main predicate from \texttt{ExistsPageN.lean} is \texttt{VWinAt}. For a fixed primitive character $\psi$, it asserts that a certain class of host fields satisfies \texttt{VVle} on strategically located scale intervals. In the displayed version relevant to the final argument, the cutoff and bound are
\begin{equation}\label{eq:vwin-cutoff}
t=\lfloor x^{401}\rfloor,
\qquad B=\sqrt{\exp(x^\eps)}=\exp(x^\eps/2),
\qquad 0<\eps<\frac1{52}.
\end{equation}
The exponent $401$ is a parameter of the formalized quantitative architecture. It is not a number of theta lifts, a cohomological degree, or a statement about the dimension of a representation. It specifies how far in the integer twist index the combined obstruction lengths must be controlled. The small parameter $\eps$ governs the acceptable error in the analytic application.

This separation of scales is instructive. Let
\[
b=\exp(x^\eps).
\]
Then~\eqref{eq:vwin-cutoff} permits an obstruction budget of size $b^{1/2}$, while a normalizing factor occurring in the theta-lift construction is of size $b$. If an estimate contains the ratio of the two, that ratio is
\[
\frac{b^{1/2}}b=\exp(-x^\eps/2),
\]
which tends to zero rapidly as $x$ increases. The quotient here is only a hypothetical normalization of a $b^{1/2}$-size error by a $b$-size main term; it is not a derivation of the actual cohomological loss. To make the comparison in the source one needs the stated normalization of the arithmetic estimate, including its host and window hypotheses. The calculation indicates what relative scale could be sufficient once those hypotheses hold.

The window predicate includes a further quantifier pattern. Given sufficiently large $X$, after choosing certain auxiliary primes $q,d$, it permits a number $Y$ satisfying
\begin{equation}\label{eq:chooseY}
X^{1/J}\le Y,
\qquad Y^2\le X,
\end{equation}
for some $J\ge1$, and it requires the cohomological estimate throughout $x\in[Y,Y^2]$. This is the source's way of locating a long arithmetic interval free of scales where the obstruction bound is inadequate.

\section{Selection of admissible scale windows}
The definition \texttt{VWinAt} permits the arithmetic estimate to hold on selected intervals $[Y,Y^2]$ rather than at every large scale. The analytic theorem has exactly the same flexibility: a sequence of sufficiently late long windows is enough. The delicate arithmetic statement is that suitable windows can in fact be found, despite the possible failures of the cohomological estimate in a cyclotomic tower.

It is useful to isolate a simple geometric fact about this kind of selection. In logarithmic coordinates, a multiplicative interval $[Y,Y^2]$ becomes $[\log Y,2\log Y]$. Large gaps in the logarithmic parameter therefore yield the intervals required by the arithmetic predicate. But the mere existence of finitely many bad intervals does not guarantee a long gap: those intervals might cover everything under consideration. One needs a bound on their total size or on their distribution.

\begin{lemma}[A gap left by short excluded intervals]\label{lem:gap}
Let $I=[0,L]$, and suppose a union of at most $R$ closed subintervals of $I$, each of length at most $h$, is removed. If $L>Rh$, the complement contains an open interval of length at least
\[
\frac{L-Rh}{R+1}.
\]
\end{lemma}
\begin{proof}
The removed set has total length at most $Rh$, so its complement has length at least $L-Rh$. The complement has at most $R+1$ connected components. At least one therefore has length at least $(L-Rh)/(R+1)$.
\end{proof}

\begin{figure}[ht]
\centering
\begin{tikzpicture}[x=.83cm,y=.9cm]
\draw[->] (0,0)--(11,0) node[right] {$u=\log x$};
\fill[red!13] (1,0.04) rectangle (2.2,.53);
\fill[red!13] (3.15,0.04) rectangle (4.15,.53);
\fill[red!13] (7.8,0.04) rectangle (8.9,.53);
\fill[blue!14] (4.4,0.04) rectangle (7.55,.53);
\draw[blue!65!black,very thick] (4.4,.04)--(7.55,.04);
% Keep each caption inside the horizontal span of the region it names.
% The two-line central label prevents collision with the right excluded interval.
\node[anchor=south,align=center,font=\footnotesize,inner sep=1pt] at (5.95,.58) {uncovered\\interval};
\node[anchor=south,font=\footnotesize,inner sep=1pt] at (1.60,.58) {excluded};
\node[anchor=south,font=\footnotesize,inner sep=1pt] at (8.35,.58) {excluded};
\end{tikzpicture}
\caption{A model of interval selection in logarithmic coordinates. To infer a gap of specified length one must bound the size of the excluded set; the picture alone proves nothing about the Iwasawa-theoretic case.}
\end{figure}

The lemma explains one elementary feature of gap selection; it is \emph{not} the arithmetic theorem establishing the windows in Alp\"oge's construction. There the excluded scales arise from Galois and Iwasawa modules, and the nontrivial task is to obtain an estimate strong enough for this kind of selection. The Lean predicate \texttt{VWinAt} records the resulting good intervals, with carefully ordered quantifiers. The detailed source theorem proving their existence remains an item for independent audit.

\section{Uniform estimates on selected windows}
A cohomological bound can be useful without asserting that the obstruction groups vanish. The actual window condition asks for their total weighted lengths to remain below a quantitative bound on a selected interval. That information can be inserted into the divisibility estimate as a controlled loss. Because the analytic criterion needs only arbitrarily late windows, the arithmetic statement need not be strengthened to an estimate at every scale.

The quantifiers matter. The field, the auxiliary primes, and the window endpoints are chosen in a prescribed order. Requiring the same field or the same primes to work at every large scale would be a different theorem. It remains to arrange that the resulting bound refers to the particular Dirichlet character under investigation, rather than to an unwieldy family of auxiliary twists.

\section{The quantified cohomological window statement}\label{sec:exact-window}

The window condition introduced earlier is more demanding than the existence of a favorable number field at some scale.  Here it helps to write the choices in their exact order.  Fix a primitive character $\psi$ and a sufficiently small $\eps>0$.  Choose a function $e(x)$ of the real scale $x\ge1$ with
\begin{equation}\label{eq:window-selector}
x^{2+8\eps}\le2^{e(x)}\le2x^{2+8\eps}.
\end{equation}
For auxiliary primes $q,d$, write $E_{\psi,e(x),q,d}$ for the cyclotomic host field described in Chapter~\ref{ch:hosts}.  In the source this is specified by a fixed-field condition, so no accidental choice of embeddings changes the assertion.  The discrete level $e(x)$ is allowed to increase as $x$ moves through a window.

At a scale $T>1$, the pair is called admissible only if it meets both the splitting congruences from Chapter~\ref{ch:hosts} and the prescribed, unequal sizes
\begin{equation}\label{eq:admissible-prime-sizes}
\tfrac12T^{\eps/5}<\log q\le T^{\eps/5},\qquad
\tfrac12T^{\eps/5}(\log T)^2<\log d\le T^{\eps/5}(\log T)^2.
\end{equation}
In addition to $q\mid d+1$, every odd prime up to $T^{\eps/5}$ is required to divide $d+1$.  The exact simultaneous existence of pairs with these properties and the additional analytic twist conditions is a separate source theorem; the present window statement begins only once such a pair is available.

\begin{theorem}[Cohomological window, cited source input]\label{thm:quantified-window}
For the fixed $\psi$, $0<\eps<1/52$, and every selector satisfying~\eqref{eq:window-selector}, the source theorem \texttt{vwin\_of\_row} supplies numbers $J\ge1$ and $X_0>1$ such that the following holds.  For every $X\ge X_0$ and every admissible auxiliary prime pair $(q,d)$ at scale $T=X^{1/J}$, there is a $Y$ with
\[
X^{1/J}\le Y,\qquad Y^2\le X,
\]
for which, \emph{simultaneously for every} $x\in[Y,Y^2]$,
\begin{equation}\label{eq:window-vv-explicit}
\mathrm{VVle}\!\left(E_{\psi,e(x),q,d},\lfloor x^{401}\rfloor,
\exp\left(\frac{x^\eps}{2}\right)\right)
\end{equation}
holds, with the meaning of \texttt{VVle} given in~\eqref{eq:vvle}.
\end{theorem}

For fixed $\psi,\eps$ and one selected window, the source-level statement keeps the same pair $q,d$ throughout that window. Its predicate \texttt{VAt} is a \emph{universal conditional}: for every number field $E$ satisfying \texttt{IsChiField} at the specified parameters, it asserts the associated \texttt{VVle} bound. It does not construct such a field. Existence and identification of eligible hosts are supplied by separate construction and supply modules, and \texttt{vvle\_congr} handles isomorphic fields. Thus the displayed $E_{\psi,e(x),q,d}$ denotes an eligible host, which may change when $x$ crosses a power-of-two level even though $q,d$ and the window endpoints remain fixed. The Galois-representation hypothesis needed to produce this window enters through the \texttt{Row0b} local-global theorem; \texttt{vwin\_of\_row} consumes that input, and \texttt{everyCharacter\_of\_row} uses the resulting window. The relevant source files are \texttt{GaloisRepresentationsOfThetaLift}, \texttt{RowSrc}, \texttt{VwinOfRow} and \texttt{MainConclusion}. The elementary interval lemma in the preceding section does not prove this deep window statement.

The order of selection can be read directly from Theorem~\ref{thm:quantified-window}. The target character, exponent, and scale selector are fixed first. The theorem gives common threshold data; for each sufficiently large $X$ it then applies to every admissible pair $(q,d)$ at the specified auxiliary scale. A suitable $Y$ may depend on that pair, but the selected pair stays fixed as $x$ ranges from $Y$ to $Y^2$. The level $e(x)$, and hence the eligible host, may change inside that interval. This is the compatibility needed when the separate twist estimate chooses its auxiliary primes.

\begin{figure}[ht]
\centering
\begin{tikzpicture}[x=.97cm,y=.9cm,>=Latex]
\draw[->] (.2,.3)--(11.6,.3) node[right] {$\log x$};
\fill[blue!8] (2.0,.32) rectangle (10.2,3.05);
\draw[blue!70,dashed] (2,.3)--(2,3.05) (10.2,.3)--(10.2,3.05);
\node[below,font=\scriptsize] at (2,.3) {$\log Y$};
\node[below,font=\scriptsize] at (10.2,.3) {$2\log Y$};
\draw[thick] (2.15,1.0)--(3.8,1.0)--(3.8,1.45)--(5.9,1.45)--(5.9,1.9)--(8.0,1.9)--(8.0,2.35)--(10.05,2.35);
\node[font=\scriptsize,align=left] at (6.0,2.78) {The level $e(x)$ changes in steps};
\node[font=\scriptsize,align=left] at (6.05,.75) {one selected pair $q,d$ throughout};
\node[font=\scriptsize,rotate=90] at (1.0,1.65) {$e(x)$};
\end{tikzpicture}
\caption{The statement of a cohomological window. The auxiliary pair is held fixed on $[Y,Y^2]$ while the level of the cyclotomic host field changes with $x$. The weighted obstruction bound must hold at every point of the interval.}
\label{fig:coho-window}
\end{figure}

The quantifier ``for every admissible pair'' prevents a serious mismatch. The analytic theorem selects $q,d$ to control unwanted Dirichlet characters; the cohomological theorem must apply \emph{after} that choice. Separate existence claims could produce different pairs, with no guarantee that either satisfies both estimates.

\section{Rescaling the twist cutoff from 401 to 400}\label{sec:400-401}

Why $401$ in~\eqref{eq:window-vv-explicit} when the geometric argument needs twists only below $x^{400}$? The extra unit pays for the rescaling that matches the field's cyclotomic level to the frame size.

Let $c_\psi\ge1$ be the character-dependent factor that occurs in the frame order, and put
\[
\kappa=c_\psi^{1/(2+8\eps)}.
\]
If $x=\kappa x'$, then
\[
c_\psi(x')^{2+8\eps}=x^{2+8\eps}.
\]
Thus a power-of-two level suitable for the cohomological scale $x'$ becomes suitable for the geometric frame at $x$ after the factor $c_\psi$ is included.  But the twist cutoff has changed from $x^{400}$ to $(x')^{401}$.  We must check that the new cutoff really is larger.

\begin{lemma}[Rescaling a window]\label{lem:window-rescaling}
Let $\kappa\ge1$ and $W\ge\max(1,\kappa^3,\kappa^{400})$.  Put $Z=\kappa W$.  For every $x\in[Z,Z^{3/2}]$, the number $x'=x/\kappa$ satisfies
\[
W\le x'\le W^2,
\qquad \lfloor x^{400}\rfloor\le\lfloor(x')^{401}\rfloor.
\]
For every $\eps>0$ one also has $\exp((x')^\eps/2)\le\exp(x^\eps/2)$.
\end{lemma}
\begin{proof}
The lower bound $x'\ge W$ follows from $x\ge Z$.  At the other endpoint,
\[
x'\le\frac{Z^{3/2}}{\kappa}
=\kappa^{1/2}W^{3/2}\le W^2
\]
by $W\ge\kappa$.  Since $x'\ge W\ge\kappa^{400}$,
\[
x^{400}=\kappa^{400}(x')^{400}\le(x')^{401}.
\]
Taking floors preserves the inequality.  The final assertion follows from $x'\le x$ and monotonicity of the exponential.
\end{proof}

This short calculation explains two conspicuous choices in the source.  A cohomological window of ratio $2$ is ample enough to contain a suitably rescaled geometric window of ratio $3/2$.  The cohomological exponent $401$ is ample enough to contain a fixed multiple of the geometric twist range $x^{400}$ once the scale is large.  The source uses the monotonicity of \texttt{VVle} to pass from the larger bound to the smaller one.  Neither exponent should be presented as intrinsic to the Riemann Hypothesis.

\section{Finiteness and uniform bounds for obstruction lengths}\label{sec:window-finiteness}

The source's \texttt{shaLen} takes values in the extended nonnegative integers before a finiteness theorem is applied.  That type is mathematically sensible: a global cohomology group invisible under the chosen local restrictions need not be known finite in advance.  The predicate \texttt{VVle} therefore asserts both the finiteness of each relevant length and the uniform inequality~\eqref{eq:vvle}.  These assertions should not be compressed into a vague statement that the obstruction groups are ``small.''

Once the lengths are finite, the finite-set formulation yields a useful consequence.  Fix the number field $E$ and a finite twist cutoff $t$.  If infinitely many coefficient primes $\ell$ had a positive inner sum $\sum_{e<t}a_{\ell,e}(E)$, the partial sum over any $N$ such primes would be at least $N\log2$, because the lengths are nonnegative integers.  It would eventually exceed the fixed bound $B$.  Thus for a fixed field and twist range, only finitely many coefficient primes contribute, and the infinite weighted total is bounded by $B$.

The same observation does not let us take a bound proved for one field and reuse it for every field in the tower.  Each $x$ may have a different host field, so the theorem must be uniform in the manner stated in~\eqref{eq:window-vv-explicit}.  This is precisely the kind of issue that the logical detail of the Lean formulation makes difficult to overlook.

\section{Iwasawa-theoretic input to the window theorem}\label{sec:iwasawa-role}

The chosen host fields contain a growing real cyclotomic component.  This brings the problem into the range of Iwasawa theory, which studies the growth of arithmetic invariants along towers of number fields. For the classical theory of such towers see \textcite{Iwasawa1973}.  It is worth distinguishing two appearances of the Greek letter $\mu$ in our story.  The mean slope $\mu(E)$ of an Hermitian lattice measures a normalized degree; Iwasawa's $\mu$-invariant measures a different kind of growth in a module over an Iwasawa algebra.  There is no implication between them simply because the notation is the same.

For abelian number fields, the vanishing of the relevant Iwasawa $\mu$-invariant is a classical theorem of Ferrero and Washington \parencite{FerreroWashington1979}. The named Lean literature input \texttt{Published.FerreroWashington} is formulated more concretely: under its specified abelian-base-field and root-of-unity hypotheses, the cardinalities of class-group quotients by the image of the $\ell$-power map are bounded uniformly along the indicated cyclotomic tower. We use the classical $\mu$-invariant language as motivation, not as a claim that the source declaration literally defines that invariant or that its equivalence has been independently established here.\footnote{See \sourcefile{src/Development/AutomorphicArithmetic/AutomorphicVarieties/PublishedTheorems/FerreroWashington.lean}{FerreroWashington.lean}, definition \texttt{UpperTheory.V.Published.FerreroWashington}.} This literature input is useful for controlling certain tower contributions, but it does not by itself imply the very strong window bound~\eqref{eq:window-vv-explicit}. The formal development combines it with local and global duality and separate estimates for character components, including cases where residual representations behave differently.  Some of those components are controlled by $p$-adic $L$-functions and Stickelberger-type information; others require automorphic congruence estimates.

The distinctions are not merely historical.  The final arithmetic bound sums lengths over \emph{all} the coefficient primes and over a rapidly growing range of twists.  An assertion of finiteness for each fixed $(\ell,e)$, even if true for all pairs separately, does not provide the uniform weighted sum needed by the lattice-index calculation.  Nor can one select separate good intervals for each character component and simply intersect them: arbitrarily late intervals need not overlap.  The source's joint window theorem exists to produce one interval that controls the entire family simultaneously.  Appendix~\ref{app:window-technical} summarizes the deferred arithmetic cases and identifies the broad source assembly; it does not claim a case-by-case Lean declaration concordance.

\chapter{Character averaging and the positive family}\label{ch:fourier}

\section{Orthogonality and character selection}

The theta-lift and slope calculations produce arithmetic inequalities for expressions involving the auxiliary quadratic field $K=\QQ(\sqrt{-d})$ and characters of the large host field $E$. Yet the analytic theorem ultimately needs a bound for a family containing a \emph{fixed} Dirichlet character $\psi$. The mismatch matters mathematically. If the inequality controls a combination of character sums with coefficients of mixed signs, it may say nothing useful about the particular sum one needs.

Finite Fourier orthogonality supplies the projection onto the desired quotient-character components. The arithmetic application requires more: the chosen coset weights must give the selected family positive coefficients, and the remaining quadratic-twist components must receive an independent estimate. The following calculation separates these algebraic and analytic requirements.

Let $A$ be a finite abelian group and let $\widehat A$ denote its character group. Every character $\chi\in\widehat A$ is a homomorphism $A\to\CC^\times$ with values on the unit circle. The elementary orthogonality identity is
\begin{equation}\label{eq:characterorthog}
\frac1{|A|}\sum_{a\in A}\chi(a)
=\begin{cases}1,&\chi=1,\\0,&\chi\ne1.\end{cases}
\end{equation}

\begin{proof}[Proof of the nontrivial case]
Choose $b\in A$ with $\chi(b)\ne1$. Translating the summation variable gives
\[
\sum_{a\in A}\chi(a)
=\sum_{a\in A}\chi(ba)
=\chi(b)\sum_{a\in A}\chi(a).
\]
Since $\chi(b)\ne1$, the sum must be zero.
\end{proof}

If $H\subset A$ is a subgroup, the same argument shows that averaging over $H$ annihilates exactly those characters that are nontrivial on $H$. This is a projection operator in the elementary Fourier sense.

\section{The quotient-character projection}

Suppose the character $\psi$ under investigation has order $o$. In the full host-character calculation the relevant finite abelian group decomposes as $G=E_0\times H$, with $E_0\cong(\ZZ/2\ZZ)^2$ carrying the two auxiliary quadratic factors. The field $K_\psi$ determines a subgroup $A_0\subset H$ fixing that field, with $H/A_0$ cyclic of order $o$. Averaging over $A_0$ gives the projection
\[
\frac1{|A_0|}\sum_{h\in A_0}F(gh),\qquad g\in G.
\]
A character survives precisely when it is trivial on $A_0$. Every character of $E_0$ is automatically trivial on this subgroup, so the surviving characters consist of all four quadratic factors multiplied by each of the $o$ quotient characters. They can be represented as
\[
\chi_{-q}^{i}\chi_{-d}^{j}\psi^k,
\qquad i,j\in\{0,1\},\quad 0\le k<o.
\]
Thus the exact projection retains $4o$ components, not merely the $o$ untwisted powers. The desired $o$ components are those with $i=j=0$; the other $3o$ are unwanted quadratic twists and must be controlled by separate estimates, including the auxiliary-prime selection discussed below. Orthogonality alone does not eliminate them. The quantitative arithmetic construction must both arrange the desired positive Fourier coefficients and bound the components not belonging to that positive family.

\section{Positive weights from character orthogonality}

The next task is to give the \emph{desired} coefficients positive weights, while treating the unwanted quadratic twists separately. This matters because the localized Landau theorem does not work for an arbitrary linear combination of oscillating character sums. The positive family defined in~\eqref{eq:Hpsi} uses weight $1$ for the trivial character and $1/(3o+1)$ for the primitive inducing character of each nontrivial power of $\psi$.

The denominator $3o+1$ looks arbitrary until one inspects the coset weights. Consider the quotient $H/A_0$, a cyclic group of order $o$. Put probability mass
\begin{equation}\label{eq:cosetweights}
\frac4{3o+1}\quad\text{on the identity coset},
\qquad \frac3{3o+1}\quad\text{on each other coset}.
\end{equation}
These masses add to $1$, because $4+3(o-1)=3o+1$. Their Fourier coefficients can be computed explicitly.

\begin{proposition}[The positive character weights]\label{prop:weights}
Let $\nu$ be the probability measure on the cyclic group of order $o$ defined by~\eqref{eq:cosetweights}. For the trivial character, $\widehat\nu(1)=1$. For every nontrivial character $\eta$,
\begin{equation}\label{eq:positive-fourier}
\widehat\nu(\eta)=\frac1{3o+1}.
\end{equation}
\end{proposition}
\begin{proof}
For the trivial character the Fourier coefficient is the total mass, namely $1$. For nontrivial $\eta$, orthogonality gives $\sum_{a\ne 1}\eta(a)=-1$. Thus
\[
\widehat\nu(\eta)
=\frac4{3o+1}+\frac3{3o+1}\sum_{a\ne1}\eta(a)
=\frac{4-3}{3o+1}=\frac1{3o+1}.
\]
\end{proof}

The difference between the weights $4$ and $3$ makes every nontrivial quotient Fourier coefficient positive! These coefficients define the target statistic $H_\psi$ in~\eqref{eq:Hpsi}; the three quadratic-twist families retained by the full host projection remain outside this positive family. The calculation is elementary, and that is useful: it isolates the finite Fourier identity from the harder question of arranging these weights by averaging over admissible arithmetic fields and controlling the unwanted twists. The latter question carries the genuine number-theoretic cost.

\begin{figure}[H]
\centering
% Extracted from DESIGN_PROTOTYPES/FIG03_fourier_character_selection_book.tex
% Requires TikZ and arrows.meta.
\begin{tikzpicture}[>=Latex,font=\small,
 bx/.style={draw=black!65,rounded corners=1.5pt,fill=black!2,minimum height=12mm,minimum width=42mm,align=center,inner sep=2.5pt},
 ar/.style={-{Latex[length=1.7mm]},line width=.7pt},
 tt/.style={font=\scriptsize,align=center}]
\node[font=\small\bfseries,anchor=west] at (-.9,1.75) {Ideal coset distribution on $H/A_0\cong C_o$};
\node[bx] (a) at (1.3,.65) {$\nu(1)=\dfrac4{3o+1}$\\$\nu(g\ne1)=\dfrac3{3o+1}$};
\node[bx] (b) at (7.0,.65) {$\widehat\nu(1)=1$\\$\widehat\nu(\eta\ne1)=\dfrac1{3o+1}>0$};
\draw[ar] (a.east)--node[above,tt] {Fourier}(b.west);
\node[tt,anchor=west] at (-.85,-.54) {$o=3$: $(4,3,3)/10\mapsto(1,1/10,1/10)$ (ideal calculation)};
\draw[black!25] (-.9,-1.07)--(9.3,-1.07);
\node[font=\small\bfseries,anchor=west] at (-.9,-1.47) {The actual host-character projection};
\node[bx] (c) at (1.2,-2.65) {$G=E_0\times H$\\average over $A_0$};
\node[bx] (d) at (7.0,-2.65) {$4o$ retained characters\\$\chi_{-q}^i\chi_{-d}^j\psi^k$};
\draw[ar] (c.east)--node[above,tt] {projection}(d.west);
\node[bx,minimum width=95mm] at (4.05,-4.3) {$o$ desired powers of $\psi$\qquad$3o$ unwanted quadratic twists};
\node[tt,text width=9.5cm,anchor=west] at (-.84,-5.27) {The ideal Fourier masses must be approximated by an arithmetic top-set distribution. The $3o$ twists need separate analytic bounds.};
\end{tikzpicture}

\caption{Two different Fourier operations in the character extraction. The upper calculation gives positive coefficients for an ideal coset distribution; its arithmetic approximation has a controlled discrepancy. The lower projection retains $4o$ host characters, of which $o$ belong to the desired family and $3o$ are auxiliary quadratic twists requiring separate bounds.}
\label{fig:v108-fourier-projection}
\end{figure}


For example, take $o=3$. The probability distribution on the three cosets is
\[
\nu=(4/10,\;3/10,\;3/10).
\]
The nontrivial characters of the cyclic group have values $1,\omega,\omega^2$, where $\omega^3=1$ and $1+\omega+\omega^2=0$. Their Fourier coefficient is
\[
\frac4{10}+\frac3{10}(\omega+\omega^2)=\frac1{10}.
\]
With a slightly larger weight on the distinguished coset, every nontrivial Fourier mode in the $o=3$ model acquires a small positive coefficient. The actual finite top-set weights need only approximate these ideal coefficients with a controlled discrepancy, while preserving the auxiliary-field and theta-lift conditions. The discrepancy belongs to the later error estimate; exact realization of the ideal coset masses is not assumed.

\section{Bounds for auxiliary quadratic twists}

The three residual quadratic-twist families have coefficients without the required positivity. Figure~\ref{fig:v108-fourier-projection} distinguishes these residual $3o$ twists from the positive ideal quotient coefficients. They therefore need a separate mean-square estimate over the admissible auxiliary prime pairs $q,d$, as described in the architecture guide. Since the pair may be chosen afresh for a selected window, such an average can support the choice of a pair for which the residual contribution is sufficiently small. This selection is a distinct input from the exact Fourier projection and must respect the source's other admissibility conditions.

Three families of quantities must be kept separate. The subgroup projection is an exact finite identity involving all $4o$ retained host characters. The ideal coset distribution assigns positive Fourier coefficients to the $o$ desired quotient characters. The arithmetic averaging must approximate that distribution and control the other $3o$ twists through a separate error estimate; positivity of the desired coefficients does not force the residual character sums to vanish. This last estimate is why the choice of admissible auxiliary primes remains part of the quantitative argument.

There is an elementary principle behind this selection. If $B_1,\ldots,B_N$ are nonnegative numbers with average at most $A$, then at least one $B_i\le A$. If there are several families of bad quantities, one may sum their squares and use the same argument to find an index at which the total is controlled. The real theorem must establish such an average in the specific ranges and congruence classes allowed for $q,d$; the elementary pigeonhole principle alone does not supply that arithmetic input.

The guide attributes a bound of the form
\begin{equation}\label{eq:twist-bound}
|g_\chi(x)|\le x^{1/2+\eps}
\end{equation}
for the selected twisted characters on the required window. In isolation,~\eqref{eq:twist-bound} is as difficult as a Riemann-Hypothesis-strength estimate. The strategy is not to obtain it uniformly for each predetermined twist. It is to choose auxiliary twists from a sufficiently large family so that the unwanted contributions can be averaged away. That is the difference between proving an estimate \emph{for all} characters and finding \emph{one suitable collection} of auxiliary characters.

\section{Progression-free index sets}

A final combinatorial device appears in the character-averaging stage: a relatively large set of indices with no short linear relations. Alp\"oge's timeline identifies Behrend's construction as part of the tower choice; the combinatorial construction is classical \parencite{Behrend1946}.\footnote{Alp\"oge, \emph{Remark}, p.~4, item~19; the specific arithmetic implementation requires the source hypotheses.} Why should a sphere enter a proof about primes? Frobenius iteration creates short arithmetic progressions among the indices. Too many selected points on one progression would count the same local contribution repeatedly.

Behrend's idea starts from the strict convexity of a sphere. Suppose $u,v,w$ are vectors of equal Euclidean norm and $u+w=2v$. Then the parallelogram identity gives
\[
\|u-w\|^2=2\|u\|^2+2\|w\|^2-4\|v\|^2=0,
\]
so $u=v=w$. Thus distinct points on a sphere cannot form a nontrivial three-term arithmetic progression. By placing integer lattice points on an appropriately chosen sphere and encoding their coordinates into integers without carries, one constructs a comparatively large progression-free set. The stronger source-specific avoidance conditions still have to be verified separately.

The full version used in the source excludes a broader family of short relations, not merely three-term progressions. In each case, geometric separation in a high-dimensional sphere turns into combinatorial sparsity along short arithmetic trajectories. A \emph{large} set is needed so averaging does not lose too much, while a \emph{sparse-incidence} property is needed so local terms are not repeatedly overcounted. The construction must satisfy both requirements simultaneously.

\section{The finite-place and Fourier identity}\label{sec:arithmetic-fourier}

Alp\"oge describes replacing Gauss-sum cancellation with ``iterated Parseval'' after simplifying a quadratic expression.\footnote{Alp\"oge, \emph{Remark}, p.~3, note~3. His account is exploratory and does not identify the exact finite identity below with that replacement.} The finite analogue needed here is character orthogonality: summing over characters first turns an oscillatory expression into an indicator for a Frobenius coset. This identity can be verified directly, without identifying it with the more elaborate argument in his remark.

Proposition~\ref{prop:weights} computes the Fourier coefficients of an ideal coset distribution. To connect these coefficients with the divisibility inequalities for integral theta-lift evaluations, one needs an exact finite-group identity expressing a weighted local sum in terms of centered character prime sums. That identity can be established for an abstract finite abelian group without reconstructing the auxiliary number field in which the arithmetic argument applies it.

Let $G$ be a finite abelian group, written additively, and let $\Phi\subset G$ be a subgroup of index two. The letter $\Phi$ now denotes a subgroup, unrelated to the meromorphic Mellin transform $\Phi(s)$ of Chapter~\ref{ch:statistic}.  Set $f=|\Phi|$, so $|G|=2f$, and let $\eta$ be the quadratic character with kernel $\Phi$.  At a rational prime $p$, let $I_p\subset G$ be its inertia subgroup and $F_p$ a representative of the Frobenius element in $G/I_p$.  For a character $\chi$ of $G$, define its prime-power coefficient by
\begin{equation}\label{eq:inertia-character}
v_\chi(p^r)=\frac1{|I_p|}\sum_{h\in I_p}\chi(rF_p+h).
\end{equation}
The coefficient is independent of the representative $F_p$, because changing that representative merely permutes the elements of $I_p$.  The average is $\chi(rF_p)$ when $\chi$ is trivial on inertia, and $0$ otherwise.  This is how ramification enters the primitive one-dimensional Artin $L$-function without separate casework for every character.

Finite character orthogonality now gives
\begin{equation}\label{eq:coset-delta}
\frac1f\sum_{\chi\in\widehat G}\overline{\chi(g)}\,v_\chi(p^r)
=\frac{2}{|I_p|}\mathbf1_{\{g\in rF_p+I_p\}}.
\end{equation}
To check it, substitute~\eqref{eq:inertia-character}, interchange the two finite sums, and use
\[
\sum_{\chi\in\widehat G}\chi(h)\overline{\chi(g)}
=|G|\mathbf1_{\{h=g\}}.
\]
Every expression in~\eqref{eq:coset-delta} is finite.  There are no convergence issues and no analytic assumptions hidden in the identity.

It is convenient to keep the weights variable.  Choose a positive parameter $a$, a positive function $w:\Phi\to\RR_{>0}$, and a bijection $\rho_t:G\setminus\Phi\to\Phi$ determined by an auxiliary translate $t$.  Define the nonnegative arithmetic sum
\begin{equation}\label{eq:finite-gain-A}
\mathcal A(a,t,w)=\sum_{p}\sum_{r\ge1}\frac{2\log p}{|I_p|}
\sum_{g\in(rF_p+I_p)\setminus\Phi}
\left(\frac{a}{p^r}-w(\rho_t(g))\right)_+.
\end{equation}
For fixed $a$ and positive weights, this is a finite sum: only prime powers smaller than $a/\min w$ can contribute.  The operation $(x)_+=\max(x,0)$ is the same ramp that entered Chapter~\ref{ch:statistic}.

For each character of $G$, define the corresponding primitive prime sum using the coefficients $v_\chi$, with the same centering as before; write it again as $G_\chi(y)$. Here $S_\chi$ denotes the uncentred ramp sum from Chapter~\ref{ch:statistic}.  For nontrivial $\chi$ set $\ell_\chi=L'(1,\chi)/L(1,\chi)$, with $G_\chi=S_\chi+y\ell_\chi$.  For the trivial character use $G_1=S_1-y\log y+(1+\gamma)y$.  Finally set
\begin{equation}\label{eq:arith-main-M}
\mathcal M(a,w)=a\left(\log a-1-\gamma+\ell_\eta
-\frac1f\sum_{\rho\in\Phi}\log w(\rho)\right).
\end{equation}
The complete identity is
\begin{equation}\label{eq:balance-exact}
\boxed{\displaystyle
\mathcal A(a,t,w)-\mathcal M(a,w)
=\frac1f\sum_{g\in G\setminus\Phi}\sum_{\chi\in\widehat G}
\overline{\chi(g)}\,w(\rho_t(g))
G_\chi\left(\frac{a}{w(\rho_t(g))}\right).}
\end{equation}
For this display the Fourier convention is fixed by~\eqref{eq:coset-delta}; replacing all characters by their inverses gives an equivalent form.  The corresponding source identity is \texttt{ch\_balance} in
\sourcefile{src/Development/AutomorphicArithmetic/DedekindLFunctions/LandauTheorems/ChBalance.lean}{ChBalance.lean}.  The exact host formalism carries additional cutoff parameters, but the elementary character calculation is the one displayed here.

Here is why the term $\mathcal M$ is exactly right.  Apply~\eqref{eq:coset-delta} to the uncentred ramp sums $S_\chi$; this gives $\mathcal A$ term by term.  Replacing $S_\chi$ by $G_\chi$ contributes correction terms.  For the trivial character the correction is
\[
a\log\frac{a}{w(\rho)}-(1+\gamma)a.
\]
For each nontrivial character it is $-a\ell_\chi$.  When summed over $g\notin\Phi$, character orthogonality removes all nontrivial terms except the index-two character $\eta$, for which $\eta(g)=-1$.  Thus the nontrivial correction contributes $a\ell_\eta$.  The trivial part contributes $a\log a-(1+\gamma)a-(a/f)\sum_\rho\log w(\rho)$.  Their sum is precisely~\eqref{eq:arith-main-M}.

\begin{figure}[ht]
\centering
\begin{tikzpicture}[>=Latex,node distance=10mm, every node/.style={font=\small,align=center},bx/.style={draw=black!35,rounded corners=2pt,minimum width=42mm,minimum height=15mm,text width=36mm,inner xsep=2pt}]
\node[bx,fill=green!5] (p) {Prime powers and\\inertia cosets\\\scriptsize $(rF_p+I_p)\setminus\Phi$};
\node[bx,right=38mm of p,fill=blue!4] (f) {All group characters\\\scriptsize finite orthogonality};
\node[bx,below=12mm of p] (a) {Finite-place gain\\\scriptsize $\mathcal A(a,t,w)$};
\node[bx,below=12mm of f] (g) {Centered prime\\statistics\\\scriptsize $G_\chi(a/w)$};
\draw[->,thick] (p)--node[above,font=\scriptsize,fill=white,inner sep=1pt] {Fourier\\expansion}(f);
\draw[->,thick] (p)--(a);
\draw[->,thick] (f)--(g);
\draw[<->,thick,accent] (a)--node[below,font=\scriptsize,fill=white,inner sep=1pt] {subtract exact\\main term $\mathcal M$}(g);
\end{tikzpicture}
\caption{The exact character identity converts a finite-place divisibility statistic into compensated prime sums. Ramification is included through the inertia average, rather than dropped as an error.}
\label{fig:arith-fourier}
\end{figure}

Two details of this identity matter in the arithmetic application.  First, if $w$ is constant, orthogonality over $G\setminus\Phi$ retains only the characters $1$ and $\eta$.  The desired nonquadratic character would disappear.  The weights must therefore \emph{vary} among embeddings.  Second, the factors $1/|I_p|$ in~\eqref{eq:finite-gain-A} cannot be omitted.  At a prime where $I_p$ has order two and meets both cosets of $\Phi$, each relevant Frobenius coset has one element outside $\Phi$; the coefficient of a corresponding $\log p$ term is then $\log p$, rather than $2\log p$.  These are exact local factors, not perturbations that become negligible at large scale.

\section{Variable weights and prime-sum arguments}\label{sec:variable-weights}

Each term $G_\chi$ on the right of~\eqref{eq:balance-exact} is evaluated at $a/w(\rho)$. For a fixed positive $a$, a larger weight $w(\rho)$ gives a smaller argument, and a smaller weight gives a larger argument. The embeddings used for the desired Fourier projection have weights near $b$, leading to $a/b\approx x$ at the target scale. Other embeddings have smaller weights and sample $G_\chi$ at larger arguments. In the source the numerator is a per-index quantity $A=a(t)$, and the stability estimate controls its deviation from the idealized relation $A=xb$.

In the source's scale model one has
\[
b=e^{x^\eps},\qquad a=xb,\qquad X'=e^{x^{\eps/3}},
\qquad b'=\frac{a}{X'}.
\]
An embedding of weight $w\approx b$ samples $G_\chi$ near $x$.  An embedding of weight $w\le2b'$ samples at least $a/(2b')=X'/2$ in this exact model. In the actual source, the per-index parameter $A=a(t)$ may satisfy $|A/b-x|\le1$ rather than $A=xb$; for $x\ge3$ the safe lower bound is $A/w\ge (1-1/x)X'/2\ge X'/3$. The estimate for the undesired characters is more manageable at these larger arguments, provided the auxiliary primes have been selected to give the required zero-free rectangles and bounds for the relevant twists.  This is why one cannot replace the variable metric by a single average metric without changing the proof.

The finer source construction distinguishes top, thick and thin portions of the embedding set.  The top portion engineers the distribution~\eqref{eq:cosetweights}; translating the thick set $T_2\subset A_0$ by all $u\in A_0$ and averaging over those subgroup elements projects the group characters onto those compatible with the chosen quotient; the thin and leftover portions are controlled by a separate large-argument estimate.  These operations are coordinated, because the same $w$ must be used on both sides of~\eqref{eq:balance-exact}.

The quantitative character-extraction bound controls the positive family $H_\psi(x)$ of Proposition~\ref{prop:weights} by several separately estimated terms. Write $\Sigma(x)$ for the sum of absolute values of the discarded host-character remainders, $E$ for the geometric error in the averaged arithmetic estimate, and $r,m,\delta_2,b$ for the positive normalization parameters of the selected top and thick sets. With this notation the source estimate has the shape
\begin{equation}\label{eq:fourier-error-budget}
H_\psi(x)\le\Sigma(x)+M R_x
+\frac{2E}{rm\delta_2b}+\frac{4}{\delta_2x},
\end{equation}
where $M$ bounds the number of annihilator characters in the relevant Fourier projection ($M=4o$ in the source); it need not equal the number retained. The quantity $R_x$ bounds the per-character change of argument and the top-set Fourier-coefficient discrepancy, which need not vanish even if the arguments equal $x$. The costs of replacing thick weights by $b$ and bounding complementary places are included in the separate final term $4/(\delta_2 x)$ of~\eqref{eq:fourier-error-budget}. The underlying arithmetic estimate supplies the required bound on $E$ uniformly over the translates used in the averaging.  This is the quantitative version of the elementary projection calculation earlier in the chapter; a single exact Fourier identity cannot replace the inequalities that make the projection effective.

The last two terms of~\eqref{eq:fourier-error-budget} are especially instructive.  A small absolute error $E$ may become large when divided by a sparse-set density $\delta_2$; the construction therefore has to control the normalized error, not merely $E$ itself.  Conversely, the auxiliary-prime selection is permitted to vary with the final window and can arrange bounds for the discarded twists that would be far too strong if demanded uniformly for every fixed auxiliary prime.  Appendix~\ref{app:fourier-technical} records the actual source interfaces used at this point.

\section{The final exponent budget}

The numerical constant attached to $\eps$ in Chapter~\ref{ch:statistic} is not a new critical exponent. It is the sum of losses accumulated in this selection process and in the earlier slope estimates. The architectural guide records the illustrative bookkeeping
\[
\frac{13}{8}+5+1+\frac38+1=9.
\]
These terms illustrate the lift, slope, density, constant-absorption, and discarded-twist losses; they are not a substitute for the source's quantitative hypotheses. Specifically, the formal theorem \texttt{core\_gm} in \sourcefile{src/Development/AutomorphicArithmetic/PrimitiveCosets/CoreGm.lean}{CoreGm.lean} gives an exponent $1/2+(e_D+e_L+7)\eps$ when its named inputs hold. The pinned producers \texttt{link\_gl} in \sourcefile{src/Development/AbelianVarietyTheory/Proofs/GlobalInterpolation/LinkGl.lean}{LinkGl.lean} and \texttt{shapeWin} in \sourcefile{src/Development/AutomorphicArithmetic/CharacterBoxes/ShapeWin.lean}{ShapeWin.lean} supply the specialization $e_D=2$, $e_L=0$ at the source-contract level, giving the displayed $9$. This does not independently reprove the many geometric premises of those producer theorems. The analytic criterion needs a \emph{fixed multiple of an arbitrarily small $\eps$}, not this particular integer. It could not tolerate a fixed exponent $2/3$ replacing a bound approaching $1/2$.

\enlargethispage{4\baselineskip}
The Fourier identity and error budget of~\eqref{eq:fourier-error-budget} now provide an estimate for the positive family defined in Proposition~\ref{prop:weights}. The theta-lift slope bounds supply its arithmetic terms, while the residual twists and stability errors are controlled separately. The resulting family is the one required by the localized window criterion; the common choice of parameters and scales is addressed in the final arithmetic balance.

\chapter{The arithmetic balance}\label{ch:balance}

The finite-place divisibility estimate and the archimedean slope estimate must be evaluated at the same fixed translate $\tau$ and over the same selected index set $I=W(u_0)$. Their principal terms contain finite-character, level and period normalizations in different forms. Those factors have to match exactly before the remaining terms can be estimated. The frame bounds ensure that the integral theta-lift evaluations can be compared over the common target despite losses of source rank.

Three normalization matches are tracked in the calculation below. The finite correction $F$ appears with opposite signs in the column and row estimates. The common level denominator contributes matching terms $mA\log g_0$. Finally, the period relation in~\eqref{eq:Omega-period} matches the factors involving $v/\pi$ and the logarithmic derivative $\ell_\eta$. These cancellations are identities among terms computed for the same translate. If a factor were estimated separately or evaluated for a different fiber, the displayed remainder would not follow.

Throughout this chapter, fix one admissible translate $\tau=\operatorname{tr}(u_0)$ and work with its column fiber $I=W(u_0)$. Write $G_\tau$ for the favorable finite-place gain attached to that same translate. The existence of uniform estimates as $u_0$ varies is a further source input; it is not obtained by allowing the translate to change independently inside a top-place sum.

\section{Degrees of fractional-ideal lattices}\label{sec:balance-degree}

A familiar one-dimensional example helps fix the sign.  Let $L=\ZZ$ inside $\RR$ with its usual norm.  Replacing $L$ by $p^{-a}\ZZ$ enlarges it: the fundamental interval has length $p^{-a}$ rather than $1$.  If degree is defined as the negative logarithm of covolume, the degree therefore \emph{increases} by $a\log p$.  In the conventional normalized Arakelov setting, the same calculation is performed at every finite place and at every embedding of a number field.

For a full Hermitian lattice $E$ of rank $d$ over a number field $K$, its normalized degree is determined by the norm of a rational generator of $\det E$ together with its finite valuations.  If a fractional ideal $J$ acts on every coordinate of $E$, then
\begin{equation}\label{eq:ideal-degree-shift}
\widehat\deg(JE)=\widehat\deg(E)+d\,\widehat\deg(J),
\qquad
\mu(JE)=\mu(E)+\widehat\deg(J),
\end{equation}
where $\mu=\widehat\deg/d$ is the mean slope.  For a fractional ideal with local exponent $-a_\mathfrak p$, the local contribution to its degree is $a_\mathfrak p\log N\mathfrak p$, with the usual degree normalization by $[K:\QQ]$.

The source's Hermitian--slope convention also includes a field-dependent normalization. In the notation of \texttt{StatementPl.lean}, its \texttt{muDef} is twice the geometric slope, shifted by $-\log 2+\log|D_K|/[K:\QQ]$. Thus $\mu_2=2\mu_{\mathrm{Ar}}$ is accurate only when $\mu_{\mathrm{Ar}}$ is itself defined with half that shift included. The shift cancels in differences over a fixed field, but it must not be dropped when comparing absolute bounds; map heights use averaged logarithms of squared operator norms. Appendix~\ref{app:slope-technical} records this precise convention.

The arithmetic column lattice is produced from a reference lattice by a controlled fractional-ideal adjustment.  Its mean slope consequently has a lower bound of the form
\begin{equation}\label{eq:balance-column}
\mu_i\ge -M_0+G_\tau+F-E_i-L_i.
\end{equation}
Here $G_\tau$ denotes the favorable finite-place gain for the fixed translate, $F$ the common finite-place correction, $E_i$ a loss associated with level and local comparison, $L_i$ a remaining nonnegative error, and $M_0$ the reference-lattice bound.  The source defines $F$ by a local sum over the supported places, including the distinguished ramified prime, the two selected level primes, and certain inert or signal primes below a cutoff; it is not supported only at primes ramified in the host field.  These symbols are bookkeeping for \emph{defined} arithmetic quantities in the construction, not estimates inferred from the elementary example.  Equation~\eqref{eq:balance-column} is one of the deep source inputs.

\section{Rank-weighted slope inequalities}\label{sec:balance-family}

Let the finite index set of column lattices be $I\ne\varnothing$, and let $d_i>0$ be their ranks.  Their contributions to a mean-slope inequality are weighted by rank:
\[
p_i=\frac{d_i}{\sum_{j\in I}d_j},\qquad
\langle Z_i\rangle=\sum_{i\in I}p_iZ_i.
\]
Because $I$ is nonempty and every $d_i$ is positive, the total rank is strictly positive and the weights are well defined.  An unweighted average over columns would generally give a different answer.  The frame theorem applies to the pooled list of source slopes, so the rank weighting is built into the argument.  The instantiated family in the source has nonempty fibers of positive rank.

For every $i\in I=W(u_0)$, the column inequality~\eqref{eq:balance-column} uses the same $u_0$, $G_\tau$, $F$, and reference bound $M_0$, although $E_i$, $L_i$ and $h_i$ may vary with the column. Taking its rank-weighted mean yields $\langle\mu_i\rangle\ge -M_0+G_\tau+F-\langle E_i+L_i\rangle$. Let $h_i$ bound the height of the corresponding evaluation map, and let $\mu_T$ be the slope of the dual row lattice.  The frame and Harder--Narasimhan calculation of Section~\ref{sec:slope-integration} produces an upper estimate
\begin{equation}\label{eq:balance-transfer}
\langle\mu_i-h_i\rangle\le\mu_T+E_{\mathrm{HN}},
\end{equation}
with $E_{\mathrm{HN}}\ge0$ controlled by the frame defects, the spread of source slopes, and the lower slope cutoff.  The direction of this inequality is easy to get wrong: it says that a sufficiently large mean slope among the sources forces a correspondingly large mean slope in the target, after allowances for map heights and defects.

Combining~\eqref{eq:balance-column} with~\eqref{eq:balance-transfer} gives
\begin{equation}\label{eq:balance-first-rearrange}
G_\tau\le\mu_T+\langle h_i\rangle-F
+M_0+E_{\mathrm{HN}}+\langle E_i+L_i\rangle.
\end{equation}
The common finite-place correction has appeared with a minus sign.  It will disappear only when the height of the row evaluation is inserted.  This matters: it is a match of two specifically normalized local contributions across the full support of $F$, not a claim that ramification is small.

\section{Period and level normalization identities}\label{sec:balance-period}

Retain the same fixed translate $\tau=\operatorname{tr}(u_0)$ and the same column fiber $I=W(u_0)$. Let $T$ denote its selected top places, and let $1\le j\le m$ enumerate the relevant theta stages. Assign each $s\in T$ a positive amplitude $a_s$, with $A=\sum_{s\in T}a_s$. For each $j$ write $w_j^\tau(\rho)=w(\tau,j,\rho)>0$ for $\rho\in\Phi$, where $f=|\Phi|$. The translation parameter $\tau$ is fixed while $s$ varies; these are different variables in the formal \texttt{ChTop} statement.

The period calculation introduces a positive scalar $v$ and a factor $\Omega$ satisfying
\begin{equation}\label{eq:Omega-period}
\log\Omega=\gamma+\log(\pi/v)-\Rea\ell_\eta,
\qquad \ell_\eta=\frac{L'(1,\eta)}{L(1,\eta)}.
\end{equation}
The period and character normalizations must match those of the height and section-count estimates. Once they do, the cancellation is exact—and elementary.  The archimedean comparison supplies an upper bound of the shape
\begin{equation}\label{eq:balance-height}
\mu_T+h_i-F\le
\sum_{s\in T}m a_s(\log a_s-1-\log\Omega)
-A\log T_i+E_D.
\end{equation}
Here $T_i$ is a positive level or section-count parameter, and $E_D$ the geometric error.  The local level estimate and the lower bound for the available sections have the respective forms
\begin{align}\label{eq:balance-level}
E_i&\le mA\log g_0+E_C,\\
\log T_i&\ge\frac1f\sum_{j=1}^m\sum_{\rho\in\Phi}
\log\left(\frac{w_j^\tau(\rho)v}{\pi}\right)+m\log g_0-E_T,
\end{align}
where $g_0\ge1$ is a common denominator and $E_C,E_T\ge0$.

The height, level and section-count inputs above are asserted \emph{for every $i\in I$ in this fixed fiber}, with the same $\tau$, $T$, $a_s$, $\Omega$, and weights $w_j^\tau$, and with the displayed error bounds valid uniformly in $i$. Since $\sum_{i\in I}p_i=1$, their rank-weighted mean gives
\[
\begin{aligned}
\mu_T+\langle h_i\rangle-F &\le \sum_{s\in T}m a_s(\log a_s-1-\log\Omega)-A\langle\log T_i\rangle+E_D,\\
\langle E_i\rangle &\le mA\log g_0+E_C,\\
\langle\log T_i\rangle &\ge \frac1f\sum_{j=1}^m\sum_{\rho\in\Phi}\log\left(\frac{w_j^\tau(\rho)v}{\pi}\right)+m\log g_0-E_T.
\end{aligned}
\]
The use of a single common main term and uniform error bounds is the substantive source hypothesis; averaging an estimate for just one selected column would not suffice. Substitute these rank-weighted inequalities into~\eqref{eq:balance-first-rearrange}. First, the $-F$ in~\eqref{eq:balance-first-rearrange} cancels the $+F$ obtained when~\eqref{eq:balance-height} is rearranged as an upper bound on $\mu_T+h_i$.  Second, the term $mA\log g_0$ from the level loss cancels the same term with the opposite sign from $-A\log T_i$.  After these two exact cancellations, the main expression remaining is
\begin{equation}\label{eq:balance-main-unexpanded}
\sum_{s\in T}\sum_{j=1}^m a_s(\log a_s-1-\log\Omega)
-\frac{A}{f}\sum_{j=1}^m\sum_{\rho\in\Phi}
\log\left(\frac{w_j^\tau(\rho)v}{\pi}\right).
\end{equation}
The terms still involving $v$ must not be estimated separately.  They are meant to cancel exactly.

\begin{proposition}[The period--Fourier identity]\label{prop:period-balance}
Under the normalization~\eqref{eq:Omega-period}, expression~\eqref{eq:balance-main-unexpanded} equals
\begin{equation}\label{eq:balance-main-exact}
\sum_{s\in T}\sum_{j=1}^m\Rea\mathcal M(a_s,w_j^\tau),
\end{equation}
where $\mathcal M$ is defined in~\eqref{eq:arith-main-M}.  The real part agrees directly with the source theorem \texttt{ch\_top\_main\_eq} in \sourcefile{src/Development/AutomorphicArithmetic/DedekindLFunctions/LandauTheorems/ChTop.lean}{ChTop.lean}.  In this quadratic case $\eta$ is real and self-inverse, and $L(1,\eta)\ne0$, so $\ell_\eta=L'(1,\eta)/L(1,\eta)$ is real; retaining $\Rea$ makes the relevant codomain explicit.
\end{proposition}
\begin{proof}
Expand the first occurrence of $-\log\Omega$ using~\eqref{eq:Omega-period}:
\[
-\log\Omega=-\gamma-\log(\pi/v)+\Rea\ell_\eta.
\]
In the second term of~\eqref{eq:balance-main-unexpanded}, expanding $\log(w_j^\tau(\rho)v/\pi)$ produces
\[
-\frac Af\sum_{j,\rho}\log(v/\pi)=-mA\log(v/\pi)=mA\log(\pi/v),
\]
because $|\Phi|=f$.  The latter cancels the $-mA\log(\pi/v)$ obtained from the first term.  The remaining terms are
\[
\sum_{s,j}a_s\left(\log a_s-1-\gamma+\Rea\ell_\eta
-\frac1f\sum_{\rho\in\Phi}\log w_j^\tau(\rho)\right),
\]
which equals the right side of~\eqref{eq:balance-main-exact} after taking real parts, by definition.
\end{proof}

\begin{figure}[ht]
\centering
\renewcommand{\arraystretch}{1.27}
\begin{tabularx}{\linewidth}{@{}>{\raggedright\arraybackslash}p{.32\linewidth}>{\raggedright\arraybackslash}p{.32\linewidth}X@{}}
\toprule
\textbf{Finite or geometric term}&\textbf{Matching term}&\textbf{Result}\\\midrule
$-F$ in the rearranged column bound&$+F$ from the row-height upper bound&The common finite-place correction cancels\\
$+mA\log g_0$ in level loss&$-mA\log g_0$ from section count&The level denominator cancels\\
$-mA\log(\pi/v)$ from period&$+mA\log(\pi/v)$ from embedding metrics&Volume normalization cancels\\\bottomrule
\end{tabularx}
\caption{Three exact matches required before errors are estimated.  The source must prove that both sides use compatible local and archimedean normalizations; this table displays the resulting algebra.}
\label{fig:three-cancellations}
\end{figure}

The cancellation takes only a few lines; arranging compatible periods, heights, and local levels takes much more. Those normalizations matter because estimating the large terms separately would leave losses the final prime-sum bound cannot absorb.

\section{The resulting bound for centered prime sums}\label{sec:balance-error}

For the fixed translate $\tau$, the geometric construction supplies one further pointwise inequality comparing the finite-place prime-power functional with the favorable gain:
\begin{equation}\label{eq:balance-finite-bound}
\sum_{s\in T}\sum_{j=1}^m\mathcal A(a_s,\tau,w_j^\tau)\le G_\tau+E_P.
\end{equation}
The untruncated prime-power expression is finite because the top and stage sets are finite and all the weights are positive. To invoke the formal \texttt{ch\_balance} identity, however, choose \emph{one finite cutoff $N$} with $a_s/w_j^\tau(\rho)\le N$ for every $s\in T$, stage $j$, and $\rho\in\Phi$. The formal \texttt{IdentU} compatibility supplies this common-cutoff premise for the chosen construction.
Combining this bound with~\eqref{eq:balance-first-rearrange}--\eqref{eq:balance-main-exact} yields the useful output
\begin{equation}\label{eq:balance-final}
\boxed{\displaystyle
\sum_{s\in T}\sum_{j=1}^m
\bigl(\mathcal A(a_s,\tau,w_j^\tau)-\Rea\mathcal M(a_s,w_j^\tau)\bigr)
\le E_{\mathrm{tot}},}
\end{equation}
where one permitted error budget is
\begin{equation}\label{eq:balance-errors}
E_{\mathrm{tot}}=M_0+E_{\mathrm{HN}}+\langle L_i\rangle
+E_D+E_C+AE_T+E_P.
\end{equation}
All main terms cancel, leaving controlled errors. Here $E_P$ and $E_{\mathrm{tot}}$ denote bounds for the chosen $\tau$; their use as \emph{common} bounds when Chapter~\ref{ch:fourier} averages over translates requires the source's uniform family hypotheses. For the fixed translate, the source uses pointwise \texttt{LossBounds} and \texttt{each\_unit\_bound\_statement} in \sourcefile{src/Formalism/SizeWindowBrackets/ChainSlopeWindows/LevelOfCongruence.lean}{LevelOfCongruence.lean}, through \texttt{PointwiseBoundVariant} in \sourcefile{src/Formalism/SizeWindowBrackets/QuotientSlopeBoundsWith.lean}{pointwise interface} and \texttt{ch\_top} in \sourcefile{src/Development/AutomorphicArithmetic/DedekindLFunctions/LandauTheorems/ChTop.lean}{ChTop.lean}.  By contrast, \texttt{MeanBoundSecondForm} bounds only an average over translates; it cannot alone justify the fixed-translate inequality~\eqref{eq:balance-final}.  Simultaneous geometric realization with the same translate, lattices, metrics and period remains an imported, unverified input (Appendix~\ref{app:theta-technical}).

Substituting the exact Fourier identity~\eqref{eq:balance-exact} into the summands of~\eqref{eq:balance-final} expresses the finite-place contribution as a linear combination of centered character prime sums at the arguments selected by the geometric weights. The preceding chapter's averaging retains the desired positive family and assigns the unwanted characters to a separate error estimate. This gives the inequality
\[
\begin{gathered}
\text{integral evaluation maps and finite-place gains}\\
\Downarrow\quad\text{frame, period and level comparisons}\\
\sum_{s,j}(\mathcal A-\Rea\mathcal M)\le E_{\mathrm{tot}}\\
\Downarrow\quad\text{finite Fourier identity and averaging}\\
H_\psi(x)\le \text{controlled analytic and geometric errors}.
\end{gathered}
\]
The auxiliary fields, theta-lift chains and cohomological windows must therefore satisfy their quantitative bounds at a common scale, with the target family and each window's permitted parameters chosen in the prescribed order. In particular, the geometric and discarded-character errors must fit the exponent budget used in the final zero-free deduction.

\part{Nonvanishing and the critical line}

\chapter{The final nonvanishing argument}\label{ch:assembly}

\section{Window bounds for the positive character family}
Recall the real weighted statistic $G_F(x)$ attached to a positive family of Dirichlet characters. The analytic theorem of Chapter~\ref{ch:statistic} allows zeros to be excluded when that statistic admits an upper bound $Cx^\theta$, with $\theta<1$, on arbitrarily late intervals of the form $[Y,Y^{1+\delta}]$. The arithmetic argument is designed to provide such intervals with exponents arbitrarily close to $1/2$.

The family must retain the desired primitive character. For a character $\psi$ of order $o$, the real statistic $H_\psi=G_{P(\psi)}$ is defined explicitly in~\eqref{eq:Hpsi}, using the \emph{distinct primitive inducing characters} of the powers of $\psi$ with coefficients $1$ and $1/(3o+1)$. In the pinned source this family is \texttt{Level.cyc}\,$\psi$; \texttt{Level.mem\_cyc} includes $\psi$, and \texttt{Level.Gm} assigns the positive character weights. Identifying the resulting source statistic with the finite-place Fourier sum additionally uses \texttt{GmAEq.lean} and its compatible host and normalization hypotheses. The finite Fourier calculation in Chapter~\ref{ch:fourier} explains these positive numbers, while the three unwanted quadratic-twist families are controlled separately. The much harder assertion is that the theta-lift and cohomological constructions control the resulting \emph{centered} prime statistic on the required windows.

A typical final bound has an exponent $1/2+C\eps$ for fixed $C$ and arbitrarily small $\eps>0$. Different intermediates in the source use different constants multiplying $\eps$; these must not be identified without checking the exact scale and family. For the final analytic implication, the important point is that $C$ remains fixed while $\eps$ can be made as small as needed. A bound at a fixed exponent such as $2/3$ would not suffice. This is why the period and integral-structure estimates cannot be treated as small technical improvements.

\section{Conditional analytic completion}

The final analytic deduction is short once the arithmetic hypothesis has been stated.
\begin{theorem}[Conditional analytic completion]\label{thm:conditional-completion}
Assume that for every primitive Dirichlet character $\psi$ and every $\theta>1/2$ sufficiently close to $1/2$, the associated positive character family $P(\psi)$, containing $\psi$ with positive weight, has a centered statistic $G_{P(\psi)}$ satisfying the hypotheses of the localized zero-free theorem, including an upper bound
\[
G_{P(\psi)}(x)\le C_\theta x^\theta
\]
on arbitrarily late intervals $[Y,Y^{1+\delta}]$ for some fixed $\delta>0$, with the family and analytic pole data fixed as required by that theorem. Then all primitive Dirichlet $L$-functions are zero-free in $\Rea(s)>1/2$, apart from the possible pole of the trivial character at $s=1$. Consequently RH and the generalized Riemann Hypothesis for Dirichlet $L$-functions follow.
\end{theorem}
\begin{proof}
Suppose a primitive Dirichlet $L$-function $L(s,\psi)$ had a zero $\rho$ with $\Rea(\rho)>1/2$. Choose $\theta$ with $1/2<\theta<\Rea(\rho)$ and use the assumed window bound. By the localized zero-free theorem, $L(s,\psi)$ has no zero in $\Rea(s)>\theta$, contradiction. Thus all primitive characters satisfy the required one-sided nonvanishing.

Every imprimitive Dirichlet character is induced by a primitive character $\psi^*$. Its $L$-function differs from $L(s,\psi^*)$ by finitely many Euler factors of the form $1-\psi^*(p)p^{-s}$, with $p$ dividing the larger modulus. Each factor is nonzero for $\Rea(s)>0$, because $|\psi^*(p)p^{-s}|\le p^{-\Rea(s)}<1$. Nonvanishing therefore extends to imprimitive characters in the right half-plane.

Finally, for a primitive character the functional equation of the completed $L$-function sends a zero in the critical strip at $s$ to a zero for the conjugate character at $1-s$. If a zero lay to the left of $\Rea(s)=1/2$, its reflection would lie to the right, where nonvanishing has just been established. Zeros in the critical strip can therefore lie only on the critical line. Taking the trivial character modulo $1$ gives the Riemann zeta function and RH.
\end{proof}

The conclusion is conditional on the arithmetic window hypothesis stated in the theorem. The deep construction theorems intended to furnish that hypothesis are identified separately by their source declarations; their proofs have not been independently reconstructed in this book. Thus the theorem specifies the analytic implication without enlarging the verification claim.

\section{Compatibility of estimates in one family}\label{sec:common-threshold}

A collection of individually correct estimates is not yet a usable collection of estimates.  Consider a simple problem.  Suppose that for every sufficiently large $x$ one can find a field with a good theta-norm bound, and that for every sufficiently large $x$ one can find a field with a good cohomological bound.  It does not follow that the \emph{same} field satisfies both.  The difficulty remains even if the fields are drawn from a very large supply.  The formal construction is careful about this logical point because it must coordinate local divisibility, geometry, character selection and the window theorem before taking $x$ large.

The source organizes its requirements through admissible families equipped with thresholds.  A threshold records lower bounds on parameters, sometimes as functions of data chosen earlier.  Replacing two thresholds by their pointwise maximum preserves any property that is monotone under making the parameters larger.  The following elementary statement captures the essential argument.

\begin{proposition}[A common threshold]\label{prop:common-threshold}
Fix $\eps>0$ in a permitted range, and suppose a class of admissible families has the following two properties.  First, for each of finitely many assertions $P_1,\ldots,P_r$, there is a threshold $K_i(\eps)$ such that every admissible family satisfying that threshold also satisfies $P_i$.  Second, for every threshold $K$ there exists an admissible family satisfying $K$.  Assume thresholds have finite pointwise maxima and satisfying a stronger threshold implies satisfying a weaker one.  Then some single family satisfies all $P_i$.
\end{proposition}
\begin{proof}
Take $K=\max_iK_i(\eps)$.  By the existence property choose a family satisfying $K$.  It satisfies every $K_i(\eps)$ and therefore every $P_i$.
\end{proof}

Two lines settle the deduction. The demanding part is the hypothesis: a family must be available \emph{after the threshold has been specified}. In the actual geometric problem the nine \texttt{Closed} requirements are combined by \texttt{topDesWin} in \sourcefile{src/Development/AutomorphicArithmetic/CharacterBoxes/DesWin.lean}{DesWin.lean}, with one family chosen \emph{after} the common threshold is assembled for each fixed sufficiently small $\eps$; this does not assert a uniform choice of family across all $\eps$. Its separate \texttt{DesignExists} and \texttt{HypV} inputs, along with the upstream geometric producers, are indispensable: taking a maximum of thresholds does not prove these premises. Appendix~\ref{app:assembly-technical} identifies the nine interfaces. The common-threshold proposition explains the last assembly step without silently changing the family halfway through the argument.

There is a second compatibility issue across arithmetic and analysis.  The analytic auxiliary-prime selection must choose one pair $q,d$ to suppress unwanted twists.  The cohomological window theorem then has to work for that same pair on the selected window. For fixed $\psi$ and $\eps$, the positive character family $P(\psi)$ remains fixed, while $q,d$ may be reselected at each sufficiently late window; within one window they are shared by the twist, cohomological, and geometric estimates.  The quantifier order in Theorem~\ref{thm:quantified-window} was designed for exactly this purpose.  With the scale rescaling of Lemma~\ref{lem:window-rescaling}, the resulting family has the bounds needed by the integral lattice construction on an interval $[Z,Z^{3/2}]$.  That ratio $3/2$ is inherited by the final analytic theorem; it is not produced anew at the last step.

\section{Specialization of the exponent to \texorpdfstring{$1/2+9\eps$}{1/2+9 epsilon}}\label{sec:final-numerical}

Four error terms separate the normalized Fourier inequality~\eqref{eq:fourier-error-budget} from the zero-free conclusion.  The discarded-character term $\Sigma(x)$ is controlled by the auxiliary-prime selection; the stability term $MR_x$ comes from argument shifts and the top-set Fourier discrepancy; and the remaining terms arise from geometric errors and the density of the chosen top set.  The choice of parameters is arranged so that, for a fixed primitive character $\psi$ of order $o$, all these contributions can be bounded on the selected window with an exponent that is a fixed multiple of $\eps$ above $1/2$.

The full source uses exponents $e_D,e_L\ge0$ for two portions of the geometric loss.  Its final scale comparison gives, in the notation of the centered positive character family, a bound of the form
\begin{equation}\label{eq:final-one-scale}
H_\psi(x)\le x^{1/2+(e_D+e_L+6)\eps}+4o\,x^{1/2+\eps}.
\end{equation}
Here the constants and order $o$ are fixed before $x$ increases. This intermediate inequality traces through \texttt{Deduction.ch\_core} in \texttt{ChCore.lean} and the bound \texttt{sigA\_le} in \texttt{SigALe.lean}; \texttt{GmAEq.lean} identifies the weighted statistic under the compatible host data. It is available only after their named geometric and character-selection hypotheses have been met. The final exponent in \texttt{CoreGm.lean} uses the same application-level nonnegative loss parameters. The significance of displaying the exact exponents is that their final absorption is simple and can be checked independently.

Put $a=e_D+e_L\ge0$.  For $x\ge1$,
\[
x^{1/2+\eps}\le x^{1/2+(a+6)\eps},
\]
so~\eqref{eq:final-one-scale} implies
\[
H_\psi(x)\le(1+4o)x^{1/2+(a+6)\eps}.
\]
Once $x^\eps\ge1+4o$, the fixed prefactor can be absorbed into one extra power $x^\eps$:
\begin{equation}\label{eq:final-general-bound}
H_\psi(x)\le x^{1/2+(a+7)\eps}.
\end{equation}
The pinned construction explicitly specializes $e_D=2$ and $e_L=0$ (hence $a=2$): \texttt{LinkGl.lean} supplies the conditional \texttt{LinkGLOBAL} input and \texttt{ShapeWin.lean} calls \texttt{topDesWin} with parameters $(C_{\mathrm{abs}},C_0,e_D,e_L)=(2,3,2,0)$. These source-level producers must still satisfy their numerous named geometric, family-existence and cohomological hypotheses; this exposition has not independently reconstructed those proofs. On that source-defined specialization,~\eqref{eq:final-general-bound} gives
\begin{equation}\label{eq:final-nine}
\boxed{\displaystyle H_\psi(x)\le x^{1/2+9\eps}
\qquad (Y\le x\le Y^{3/2})}
\end{equation}
for arbitrarily large selected $Y$.  The constants implicit in the preceding estimates depend on the fixed character and on $\eps$, but not on the running scale once its threshold has been chosen.  That last clause is part of the quantitative assertion.  If a supposedly fixed constant were allowed to grow arbitrarily with $x$, the exponent absorption would be invalid.

The number $9$ is particular to this parameter choice. What the analytic argument needs is a coefficient independent of $\eps$: as $\eps\to0$, the exponent approaches $1/2$ and falls below the real part of any hypothetical off-line zero. The pinned Lean interfaces select $e_D+e_L=2$, yielding $9$ after the checked prefactor absorption. This identifies the source of the numerical value; it does \emph{not} amount to an independent derivation of the upstream construction theorems furnishing the hypotheses of \texttt{LinkGLOBAL}, \texttt{ch\_core} or the window theorem. The calculation from~\eqref{eq:final-one-scale} to~\eqref{eq:final-nine} is elementary; establishing the intermediate bound from the arithmetic objects is the difficult part.

\section{The zero-free consequence of the window bound}\label{sec:nine-to-zero-free}

Take a primitive character $\psi$ and suppose it had a zero $\rho$ with $\Rea(\rho)>1/2$.  Choose a real number $\theta$ satisfying
\[
\frac12<\theta<\min\{\Rea(\rho),1\}.
\]
Because $9$ is fixed, we can choose a sufficiently small \emph{fixed} $\eps>0$ so that $1/2+9\eps\le\theta$, and then invoke the source-provided compatible family for that fixed $\eps$. For each selected window the auxiliary primes $q,d$ may be chosen anew, but the positive family $P(\psi)$ is unchanged and all estimates for that window use the same pair. Thus Equation~\eqref{eq:final-nine} gives
\[
H_\psi(x)\le x^\theta
\qquad (Y\le x\le Y^{3/2})
\]
on arbitrarily late intervals.

\begin{figure}[H]
\centering
% Central quantitative contradiction, not a generic proof flowchart.
% The right plot is the ratio of two scales, not the actual value of A_k.
\begin{tikzpicture}[x=.85cm,y=.82cm,>=Latex,every node/.style={font=\footnotesize}]
\node[anchor=west,font=\footnotesize\bfseries] at (.35,8.55) {Arithmetic input on arbitrarily late windows};
\node[anchor=west,text=black!75,font=\scriptsize] at (.35,8.02) {Theta evaluations, periods, lattice slopes, cohomological windows, character averaging};
\draw[accent!50,line width=.8pt] (.35,7.67)--(13.65,7.67);
\node[align=center,fill=blue!5,draw=accent!45,inner sep=6pt,rounded corners=2pt,text width=11.8cm] at (7.0,6.93)
 {$\displaystyle H_\psi(x)\le x^{1/2+9\varepsilon}$\quad for $Y\le x\le Y^{3/2}$, with one fixed positive family};
\node[anchor=west,font=\footnotesize\bfseries] at (.35,5.93) {Hypothetical pole to the right of $\theta$};
\node[anchor=west,font=\footnotesize\bfseries] at (7.35,5.93) {Comparison of derivative scales};
% geometric pole sketch
\draw[->,black!70] (.9,2.12)--(6.3,2.12);
\node[font=\scriptsize,anchor=west] at (6.33,2.23) {$\Re(s)$};
\draw[->,black!70] (.9,2.0)--(.9,5.25);
\node[font=\scriptsize,anchor=east] at (.78,5.12) {$\Im(s)$};
\draw[black!40,dashed] (1.94,2.12)--(1.94,5.36);
\draw[accent!70,dashed] (2.74,2.12)--(2.74,5.36);
\draw[accent!70] (5.65,2.12)--(5.65,5.33);
\filldraw[red!70!black] (3.86,3.78) circle (2.05pt);
\filldraw[accent] (5.65,3.96) circle (2.05pt);
\draw[<->,black!65,line width=.85pt] (3.92,4.01)--(5.59,4.18);
\node[font=\scriptsize,fill=white,inner sep=1.4pt] at (4.77,4.28) {$d$};
\node[font=\scriptsize,anchor=east] at (3.77,3.74) {$\rho$};
\node[font=\scriptsize,anchor=west] at (5.72,3.88) {$z$};
\node[below,font=\scriptsize] at (1.94,2.11) {$\frac12$};
\node[below,font=\scriptsize] at (2.74,2.11) {$\theta$};
\node[below,font=\scriptsize] at (3.86,2.11) {$\beta$};
\node[below,font=\scriptsize] at (5.65,2.11) {$\sigma$};
\draw[<->,accent!75] (2.76,1.39)--(5.64,1.39);
\node[fill=white,inner sep=1.8pt,font=\scriptsize] at (4.19,1.39) {$\sigma-\theta$};
\node[font=\scriptsize,anchor=west] at (1.25,.63) {$\theta<\beta$ and, after isolation, $d<\sigma-\theta$};
% ratio plot
\draw[->,black!75] (7.68,2.06)--(13.64,2.06) node[right,font=\scriptsize] {$k$};
\draw[->,black!75] (7.68,2.06)--(7.68,5.47);
\draw[densely dashed,gray!70] (7.7,2.82)--(13.53,2.82);
\node[font=\scriptsize,anchor=east] at (7.59,2.82) {$0$};
\node[font=\scriptsize,align=left,anchor=west] at (8.06,5.35) {$\log\!\left(\dfrac{\text{pole scale}}{\text{ceiling}}\right)$};
\draw[accent,line width=1.2pt] (8.05,2.24)--(13.11,4.76);
\foreach \x/\y in {8.85/2.64,9.85/3.14,11.1/3.77,12.25/4.34} {
  \filldraw[fill=accent,draw=white,line width=.4pt] (\x,\y) circle (2.0pt);
}
\node[font=\scriptsize,text=accent,align=left,anchor=west] at (10.43,3.08) {positive slope\\$\log((\sigma-\theta)/d)$};
\node[font=\scriptsize,anchor=west] at (8.0,1.51) {Dots: orders selected by successive windows};
\draw[black!35] (.35,.36)--(13.7,.36);
\end{tikzpicture}

\caption{How the arithmetic estimate is used to exclude a hypothetical off-line zero. The geometric sketch shows a pole $\rho$ to the right of $\theta$ and a comparison center $z$ at distance $d<\sigma-\theta$. The right-hand plot shows the \emph{scale ratio} of its isolated leading pole term to the permitted derivative ceiling; it grows along the selected derivative orders. This is not a plot of actual zeros or prime-sum values. The source's nearest-pole dichotomy and positive real residues are essential for converting this scale comparison into a contradiction (Chapter~\ref{ch:statistic} and Appendix~\ref{app:landau-derivatives}).}
\label{fig:final-chain}
\end{figure}

By Proposition~\ref{prop:positive-poles}, the positive weights of Chapter~\ref{ch:fourier} prevent a zero of $L(s,\psi)$ from disappearing in the combined pole residues. The localized argument of Chapter~\ref{ch:statistic} then excludes that pole, establishing $L(s,\psi)\ne0$ for $\Rea(s)>1/2$.

By Lemma~\ref{lem:euler}, the extra Euler factors of an imprimitive character are nonzero throughout $\Rea(s)>0$, so the primitive-character conclusion extends to its imprimitive forms. The functional equation then reflects a hypothetical zero with real part below $1/2$ to a zero above $1/2$, contrary to the established right-half-plane nonvanishing. For $\zeta$, the zeros outside the critical strip remain the familiar trivial zeros at the negative even integers.



The estimate in Figure~\ref{fig:final-chain} is obtained on selected late windows for one fixed centered positive character family. Each choice of window determines a suitable derivative order. The isolated pole term would exceed the resulting derivative ceiling if a zero lay to the right of $\theta$; the analytic theorem handles the pole-separation and positivity conditions needed for that comparison. The exact source hypotheses for constructing the arithmetic bound are recorded in the appendices.

\section{Formal statement of the RH conclusion}

The top-level source file \sourcefile{src/Main.lean}{Main.lean} contains theorem declarations for the Riemann Hypothesis and Dirichlet-character analogues. It imports the compiled formal argument from \texttt{Supplements.Kernels.MainTheoremsAll}. The relevant bridge, in
\sourcefile{src/Supplements/Kernels/MainTheorems.lean}{MainTheorems.lean}, has the mathematical structure
\[
\begin{gathered}
\texttt{allChi}:\ \texttt{UpperTheory.EveryCharacter},\\
\texttt{rhOfAll}:\ \texttt{RiemannHypothesis}.
\end{gathered}
\]
The second statement is obtained by applying the source's analytic reduction to the first. The definition of \texttt{UpperTheory.EveryCharacter} is exactly the one-sided nonvanishing statement
\[
\forall N,\chi,s,\qquad
\Rea(s)>\tfrac12,
\quad(\chi\ne1\ \text{or }s\ne1)
\quad\Longrightarrow\quad L(s,\chi)\ne0.
\]
It is important that the source quantifies over all Dirichlet characters, not only primitive ones. The finite Euler-factor calculation in the preceding proof explains why such a strengthening is natural.

A concise source-level dependency chain is:
\begin{enumerate}
\item The Galois and local-global development provides the hypotheses fed into \texttt{everyCharacter\_of\_row} in \texttt{MainConclusion.lean}.
\item A cohomological window theorem is converted into suitable host supply; the file \texttt{SupplyChi.lean} visibly takes \texttt{VWinAt}\,$\psi$\,$401$ as an input.
\item Theta-lift, period, slope, and arithmetic character-family results provide the quantitative input of the nonvanishing machinery, represented in the formalism by propositions including \texttt{TopG}.
\item The analytic window theorem converts bounded positive character statistics into the appropriate zero-free assertion.
\item \texttt{UpperTheory.EveryCharacter} feeds the reflection and finite Euler-factor reductions, yielding the final target.
\end{enumerate}
A full source audit must verify every arrow, with the definitions at the two ends of each arrow fixed at the same source revision. The filenames alone are not a dependency proof; the statement is an editorial description of the mathematical architecture and of the source-level theorems we have directly inspected.

\section{Scope and remaining verification}
Under their stated hypotheses, this exposition proves the Euler-factor reduction, ramp transform, pole-residue positivity, Fourier orthogonality, frame inequalities, integrated slope calculation, displayed cancellations, window rescaling, and exponent absorption. The theta and integral models, uniform covolume and height bounds, auxiliary-prime selection, and Galois-cohomological window theorem rely on deeper cited Lean declarations. Appendix~\ref{app:verification} distinguishes checks of the formal archive from review of this exposition.

Andrew Sutherland's public report, updated 21 November 2026 at 12:55 UTC, records that all 175,187 Lean modules compile, that the proof replays through Lean's kernel and independent kernels, and that five public statements match a separately prepared challenge. Comparator and export-replay checks were still pending at that time. The report and this manuscript's pinned public mirror are separate cited artifacts; no archive-wide byte comparison between them has been established. This exposition did not repeat those checks, which do not certify that it correctly explains every deep dependency. Appendix~\ref{app:verification} records the reported checks and their limits.

Further review should examine the intermediate theorems, their normalizations, and their compatibility with the source definitions. Any simplification must preserve their mathematical content. The present chapters identify where the deepest inputs enter the prime-sum estimate, explain the hypotheses they must supply, and prove the shorter deductions connecting them to the zero-free theorem. This distinguishes what the formal source asserts from what the present exposition has established independently, and gives subsequent reviewers precise points at which the reconstruction can be checked or improved.

\appendix
\part{Technical notes and source concordance}

\chapter{Analytic normalization and localization}\label{app:analytic}

The main argument uses the ramp transform and the concentration of derivatives of a Laplace transform.  The ordinary calculations are collected here for readers who want to check the normalizations and the one-sided estimate without interrupting the exposition of the number-theoretic construction.

\section{The compensation at \texorpdfstring{$s=1$}{s=1}}\label{app:mellin-normalization}

This section supplies the principal-character companion to Proposition~\ref{prop:transform} in Chapter~\ref{ch:statistic}. Recall $S_\chi(y)=\sum_{n\le y}\Lambda(n)\chi(n)(y/n-1)$.  For a nontrivial primitive character, put $G_\chi(y)=S_\chi(y)+y\ell_\chi$, where $\ell_\chi=L'(1,\chi)/L(1,\chi)$.  For the primitive trivial character use $G_1(y)=S_1(y)-y\log y+(1+\gamma)y$.  Direct integration gives, initially for $\Rea(s)>1$,
\begin{align*}
F_\chi(s)&:=\int_1^\infty G_\chi(y)y^{-s-1}\,dy
=-\frac{L'(s,\chi)/L(s,\chi)-s\ell_\chi}{s(s-1)},\qquad\chi\ne1,\\
F_1(s)&=-\frac{\zeta'(s)/\zeta(s)+1/(s-1)+1-(1+\gamma)s}{s(s-1)}.
\end{align*}
The second formula follows from
\[
\int_1^\infty y\log y\,y^{-s-1}\,dy=\frac1{(s-1)^2},
\quad\int_1^\infty y\,y^{-s-1}\,dy=\frac1{s-1}.
\]
At $s=1$, the first numerator vanishes, since $L'/L$ is regular there.  For the trivial character, the local expansion $\zeta'/\zeta=-1/(s-1)+\gamma+O(s-1)$ shows that its numerator also vanishes.  The compensation therefore removes the apparent singularity at $s=1$.  Away from $s=0,1$, a zero of order $m$ of an included $L$-function still supplies residue $-m/[\rho(\rho-1)]$.  A different centering would give a different formula at $s=1$ and must be followed consistently through the final arithmetic identity.

\section{An exponential tail estimate for the kernel}\label{app:kernel}

Fix $\delta>0$ and put
\[
\ell=2+\frac6\delta,\quad M=\ell(1+\delta),\quad
\sigma=M(\ell+2),\quad b=1+\frac\delta2,\quad
k=\lceil b\sigma U\rceil.
\]
The values $\ell,M,\sigma,b$ depend only on the fixed window-width parameter $\delta$; $k$ is the integer chosen anew for each $U$. The chosen constants yield a uniform estimate; minimizing the derivative order is unnecessary for this sufficient tail bound.  Set $V=(1+\delta)U$ and $a_\pm=\sigma-2\pm\ell$, both positive.  The elementary bounds
\[
\mathbf1_{u>V}\le e^{\ell(u-V)},\qquad
\mathbf1_{u<U}\le e^{\ell(U-u)}
\]
and the gamma integral $\int_0^\infty u^ke^{-au}\,du/k!=a^{-k-1}$ imply
\begin{equation}\label{eq:app-tail}
\int_{[0,\infty)\setminus[U,V]}\frac{u^k}{k!}e^{-(\sigma-2)u}\,du
\le \frac{e^{-\ell V}}{a_-^{k+1}}+\frac{e^{\ell U}}{a_+^{k+1}}.
\end{equation}
Now $b(\ell+2)=M-1$ and $b(\ell-2)=\ell+1$.  Using $\log(\sigma/a)\le(\sigma-a)/a$ for $a>0$, together with $b\sigma U\le k\le b\sigma U+1$, shows
\[
e^{-\ell V}a_-^{-k-1}\le\frac e{a_-}\sigma^{-k},
\qquad e^{\ell U}a_+^{-k-1}\le\frac1{a_+}\sigma^{-k}.
\]
The first follows because $(\ell+2)/a_-=1/(M-1)$; the second uses $b\sigma(\ell-2)\ge\ell a_+$.  Thus, for every $U\ge1$,
\begin{equation}\label{eq:tail-strong}
\int_{[0,\infty)\setminus[U,(1+\delta)U]}
\frac{u^k}{k!}e^{-(\sigma-2)u}\,du
\le C_\delta\sigma^{-k},\qquad
C_\delta=\frac e{a_-}+\frac1{a_+}.
\end{equation}
For $0<\theta<1$ as required by the localized zero-free theorem, the tail is exponentially small relative to the main scale $(\sigma-\theta)^{-k}$: the ratio is $(1-\theta/\sigma)^k$. Merely applying Chebyshev's inequality to the normalized gamma distribution would not be enough once the coarse growth bound $|h(u)|\le C_0e^{2u}$ has been restored. The exponential tail estimate is the needed quantitative fact.

In the simpler sufficient \emph{pointwise-window} case, for a real $h$ bounded in absolute value by $C_0e^{2u}$ and with $h(u)\le e^{\theta u}$ on arbitrarily distant windows, the inequality
\[
|A_k(\sigma+it)|+\Rea A_k(\sigma)
\le 2\int_0^\infty\max(h(u),0)\frac{u^k}{k!}e^{-\sigma u}\,du
\]
then gives an upper bound by a fixed multiple of $(\sigma-\theta)^{-k}$.  To rule out a pole, compare that growth with $|c_\rho|d^{-k-1}$ for a uniquely nearest pole at distance $d<\sigma-\theta$.  A small vertical perturbation of the center separates competing nonreal poles: within a fixed compact region, equality of distances to two distinct poles occurs at only finitely many heights, except for points with equal imaginary part, whose unequal real parts already separate them because the center $\sigma>2$ lies to the right of the relevant poles. A pole on the real axis is handled by its positive residue.  These are the two geometric cases in the localization argument of
\sourcefile{src/Development/AutomorphicArithmetic/DedekindLFunctions/LandauTheorems/landauWindow.lean}{landauWindow.lean}.  The argument needs the meromorphic continuation, local finiteness of the pole set, nonzero residues, and positivity at real poles, exactly as described in the main chapter.  This pointwise-window argument is a sufficient special case; the source's \texttt{WindowHypU} permits a more flexible integrated upper envelope, and its \texttt{PoleData} carries additional analytic hypotheses. The comparison is established in this sufficient pointwise-window setting; the remaining parameterized source hypotheses are not independently audited here.



\section{Derivative bounds and the pole comparison}\label{app:landau-derivatives}

The tail estimate also clarifies the analytic mechanism behind the localized zero-free criterion. Let $h:[0,\infty)\to\RR$ be measurable, with $|h(u)|\le C_0e^{2u}$, and initially set
\[
\Phi(s)=\int_0^\infty h(u)e^{-su}\,du,\qquad \Rea(s)>2.
\]
At a real $\sigma>2$, its normalized derivatives are
\begin{equation}\label{eq:app-landau-deriv}
 A_k(\sigma+i\gamma):=\frac{(-1)^k}{k!}\Phi^{(k)}(\sigma+i\gamma)
 =\int_0^\infty h(u)\frac{u^k}{k!}e^{-(\sigma+i\gamma)u}\,du.
\end{equation}
All these integrals converge absolutely in the stated half-plane. Since $h$ is real, the triangle inequality gives the useful one-sided inequality
\begin{equation}\label{eq:app-landau-positive}
\begin{aligned}
 |A_k(\sigma+i\gamma)|+\Rea A_k(\sigma)
 &\le\int_0^\infty\bigl(|h(u)|+h(u)\bigr)\frac{u^k}{k!}e^{-\sigma u}\,du\\
 &=2\int_0^\infty h(u)_+\frac{u^k}{k!}e^{-\sigma u}\,du.
\end{aligned}
\end{equation}
There is no lower bound on $h$ here. The nonnegative quantity $|h|+h=2h_+$ makes an \emph{upper} estimate over selected intervals sufficient.

For clarity, suppose first that on arbitrarily late intervals $[U,(1+\delta)U]$ the stronger pointwise bound $h(u)\le e^{\theta u}$ holds, with $0<\theta<1$ and fixed $\delta>0$. Choose $\sigma$ and $k=k(U)\ge U$ as in the preceding section. On the interval, the right side of~\eqref{eq:app-landau-positive} is bounded by the complete gamma integral with exponent $\sigma-\theta$; outside it, use the coarse bound $h_+(u)\le C_0e^{2u}$ and~\eqref{eq:tail-strong}. Thus, uniformly in $\gamma$,
\begin{equation}\label{eq:app-landau-combined}
 |A_k(\sigma+i\gamma)|+\Rea A_k(\sigma)
 \le 2(\sigma-\theta)^{-k-1}+2C_0C_\delta\sigma^{-k}
 \le C(\sigma-\theta)^{-k},
\end{equation}
where $C$ is independent of $U,k,\gamma$. The source's \texttt{WindowHypU} is more general: it assumes an integrated upper envelope rather than this pointwise hypothesis. The same positive-part calculation is used there after the envelope has been inserted.

Suppose now that the meromorphic continuation of $\Phi$ has an isolated simple pole at $\rho$ with nonzero residue $c_\rho$. If $z$ is a point at distance $d=|z-\rho|$ from that pole, and the analytic remainder extends to a disk of radius $r>d$ around $z$, Cauchy's estimate gives
\begin{equation}\label{eq:app-pole-asym}
 A_k(z)=\frac{c_\rho}{(z-\rho)^{k+1}}+O(r^{-k}).
\end{equation}
Its pole term has magnitude $|c_\rho|d^{-k-1}$. When $d<\sigma-\theta$, this grows faster with $k$ than the permitted upper scale $(\sigma-\theta)^{-k}$. The remaining issue is to make that comparison without losing the lower bound through other poles or through the real term in~\eqref{eq:app-landau-combined}. The formal argument uses a nearest-pole dichotomy: a suitable shifted center isolates a nonreal pole while controlling the real-center contribution, or a real pole has the positive residue needed to prevent cancellation. Local finiteness of the pole set, analyticity away from its zeros and the sign condition on real residues are therefore genuine hypotheses, not consequences of the kernel computation. The source implements these cases in \sourcefile{src/Development/AutomorphicArithmetic/DedekindLFunctions/LandauTheorems/landauWindow.lean}{landauWindow.lean}. This appendix derives the quantitative bounds; it does not claim to reprove that entire source-level pole-separation argument.

\chapter{Integral theta maps and slope normalizations}\label{app:geometry}

\section{Integral theta comparisons}\label{app:theta-technical}

The source's geometric input contains more than the standard existence of a theta correspondence.  The selected spaces depend on parameters $k,m,t$ with, among other rank constraints, $k,m$ even, $t$ odd, $m\ge2$, $m+3\le k$, $t\ge k+m+4$, and $t>3m$. Increasing-rank spaces $X_j$ and $Y_j$ and fixed-rank changing-signature spaces $W_j$ are arranged so that successive theta operators admit the required pairings.  For $1\le j\le m$, the fixed-rank spaces $W_j$ have schematic signature $(k+j-1,m-j)$ and become definite at $j=m$; this index range is part of the description, not an extrapolation to larger $j$. The pinned geometric data remain the authority for the exact stage indexing.

The column operators are weighted sums over finite lattice-class sets, with stabilizer denominators.  If $\mathcal C_j$ is such a set and $a_j(c)$ the order of the stabilizer of $c$, a typical stage has the form
\[
\Phi_{j+1}(c')=\sum_{c\in\mathcal C_j}a_j(c)^{-1}K_j(c,c')\Phi_j(c).
\]
The kernel $K_j$ includes the theta-unfolding operator and determinant-character factors at the archimedean places. This is a schematic kernel, not the exact exponent-bearing stage operator: the prescribed archimedean powers are not displayed here and must be checked in the pinned construction. A common integer denominator $D_\Gamma$ is divisible by all stabilizer orders occurring along the chain.  It is indispensable to the integral comparison: the rational operator cannot simply be treated as integral by changing notation.

The final pairing is first a rational comparison $N_u\otimes K\to R^\vee\otimes K$ over a common coefficient field. It is not thereby an integral map $N_u\to R^\vee$: compatible containment of the chosen lattices must be established after clearing a common denominator. The Bessel comparison identifies the rational pairing with the analytic one up to unitary changes of coordinates and a nonzero column-dependent scalar. Neither this identification nor rationality alone bounds that scalar, the integral denominators, or the operator height; those estimates are separate linked divisibility and row/column inputs.  The source also estimates the covolumes of the integral automorphic lattices through covering levels, arithmetic section counts, and comparisons of Hermitian norms.  The results are assembled in source interfaces including
\sourcefile{src/Development/AutomorphicArithmetic/CharacterBoxes/ShapeWin.lean}{ShapeWin.lean} and the geometric covolume modules named \texttt{TopCovolSlim2Use} and \texttt{CovolumeFromGramAndBounds}.  These are cited construction results.  Their proofs are not reproduced in this book, and their uniformity in varying fields and weights must not be attributed to theta correspondence alone.

\section{Frame and Hermitian-slope conventions}\label{app:slope-technical}

For a Hermitian lattice $E$ of rank $d$, write the geometric mean slope as $\mu_{\mathrm{geom}}(E)=\widehat\deg(E)/d$ before field normalization. The source's \texttt{muDef} in \texttt{StatementPl.lean} is $2\mu_{\mathrm{geom}}(E)-\log2+\log|D_K|/[K:\QQ]$; the analogous maximum/minimum slopes have the same shift. Thus a notation $\mu_2=2\mu_{\mathrm{Ar}}$ is valid only if $\mu_{\mathrm{Ar}}$ has already absorbed half this field shift. The corresponding source operator height averages logarithms of squared norms. All slopes, heights and normalization shifts must be reconciled before importing numerical bounds.

Let $E_i\to F$ be the integral maps from the columns to the row target.  Their high-slope filtrations obey a morphism estimate of the form
\[
B_i(E_i^{\ge t})\subset F^{\ge t-h_i},
\]
where $h_i$ is a logarithmic operator-height bound. The vector-space frame at a chosen embedding is required to be almost isometric on the source side and almost flat on the target side; the source compares the arithmetic target rank $\dim\kappa$ with the lifted frame rank $\dim(\kappa\oplus\kappa_1)$ through $(1-\delta_F)(\dim\kappa+\dim\kappa_1)\le\dim\kappa$. Thus $\delta_F$ measures this rank-comparison allowance, not an independently measured rate at which original target directions were discarded. If the shifted source slopes have mean $s$ and maximum at most $s+u$, and the \emph{minimum} target slope is at least $s-H_0$, the conclusion takes the form
\[
\mu_{\mathrm{src}}(F)\ge s-\frac{4u\log(1+H_0/u)+(\eta_F+4e_F)H_0}{1-\delta_F}.
\]
Here $\mu_{\mathrm{src}}(F)$ denotes the target's normalized \texttt{muDef} value in \texttt{StatementPl.lean} (the source theorem writes \texttt{muDef b' X' cc}), \emph{not} its unshifted geometric mean slope. The source parameters $s,u,H_0$ and all source slope comparisons use the same normalization.
Here $e_F=\sqrt{e_{\mathrm{Lean}}}$: the source \texttt{TheoremFSharp} has a loss $4\sqrt{e_{\mathrm{Lean}}}$, while Proposition~\ref{prop:frame-cover} writes its Hilbert--Schmidt defect as $e_F^2R$ and correspondingly obtains $4e_F$. The displayed logarithmic term is interpreted by continuity at $u=0$. The finite-dimensional estimate is proved in Proposition~\ref{prop:frame-cover}; the integrated rank-loss calculation is in Section~\ref{sec:slope-integration}.  The full arithmetic theorem, including the field normalization, integral maps and quotient realization, is packaged in
\sourcefile{src/Formalism/HarderNarasimhan/HarderNarasimhanProofs/TheoremFSharp.lean}{TheoremFSharp.lean} and \sourcefile{src/Formalism/HarderNarasimhan/ShiftInvariance.lean}{ShiftInvariance.lean}.  The main text uses this theorem precisely where it converts the source means into the target slope inequality~\eqref{eq:balance-transfer}.

\chapter{Arithmetic selection and character errors}\label{app:selection}

\section{Inertia, weights and discarded characters}\label{app:fourier-technical}

The exact identity~\eqref{eq:balance-exact} explains how an arithmetic gain becomes a Fourier series of centered prime sums. In the error bound~\eqref{eq:fourier-error-budget}, the multiplicity $M=4o$ counts the annihilator-character projection, not merely the final retained family. Its two delicate companions are the distribution of embedding weights and the estimates for characters not retained by the desired cyclic family.  The relevant source modules are
\sourcefile{src/Development/AutomorphicArithmetic/DedekindLFunctions/LandauTheorems/ChBalance.lean}{ChBalance.lean}, \sourcefile{src/Development/AutomorphicArithmetic/DedekindLFunctions/LandauTheorems/ChCore.lean}{ChCore.lean}, and the final normalization in \texttt{CoreGm.lean}.

A fixed primitive character $\psi$ of order $o$ contributes the character family consisting of the primitive inducing characters for $1,\psi,\ldots,\psi^{o-1}$.  The construction places normalized mass $4/(3o+1)$ on the identity coset and $3/(3o+1)$ on each other coset; the elementary Fourier transform gives weight $1/(3o+1)$ for every nontrivial character in the quotient.  The host can carry additional quadratic twists introduced by the auxiliary primes.  Their prime sums are not allowed to join the retained positive family with arbitrary signs.  The auxiliary-prime selection instead supplies bounds for these discarded characters over the same range of scales, including separate attention to real characters and possible exceptional zeros.

For an arithmetic metric that varies among embeddings, the expression $wG_\chi(a/w)$ has a useful local continuity bound, provided $a>0$ and $w,w'>0$:
\[
\bigl|wG_\chi(a/w)-w'G_\chi(a/w')\bigr|
\le\frac{3a}{\min(w,w')}|w-w'|.
\]
It follows from the 1-Lipschitz property of the positive-part ramp (including its breakpoints) and a coarse bound on the sum of von Mangoldt weights. For nontrivial characters the correction $a\ell_\chi$ is constant as $w$ varies; the trivial character's $-a\log(a/w)+(1+\gamma)a$ contribution varies with $w$ and is handled separately in the source's character-by-character continuity theorem.  This estimate lets the construction replace approximately equal thick weights by a common central value, at an explicit cost.  The fine definitions of top, thick and thin embedding sets and the necessary sparse-incidence conditions belong to the arithmetic character-selection modules; the simplified explanation in the main text does not prove their existence.

\section{Cohomological window inputs}\label{app:window-technical}

The local-global kernel in~\eqref{eq:sha} uses the coefficient module $\QQ_\ell/\ZZ_\ell(-e)$ with a negative cyclotomic twist.  Its length need not be finite before the relevant theorem is applied.  The source organizes the estimate by character components, coefficient prime $\ell$ and twist $e$.  Among the cases treated separately are generic odd-prime components, resonant primes where the residual character changes, the trivial-character component, and boundary twists.  A bound for each case must be summed before the window is selected; choosing a new window for each case would not supply a common interval. The producer modules cited below identify the aggregate construction, not a one-to-one reference for each of the listed cases.

The tower argument uses the vanishing of Iwasawa's $\mu$-invariant for abelian number fields, together with local/global duality and cohomological estimates.  The fundamental classical result is due to Ferrero and Washington.  The conclusion required here is considerably more quantitative than that classical theorem alone.  The formal assembly passes through \texttt{VwinOfRow.lean}, \texttt{ExistsPageN.lean}, and the host-supply theorem in
\sourcefile{src/Development/AutomorphicArithmetic/CharacterBoxes/SupplyChi.lean}{SupplyChi.lean}.  The last converts the cohomological window hypothesis for a primitive character with twist cutoff $\lfloor x^{401}\rfloor$ into the host existence needed by the geometric construction: the cohomological statement ranges over positive twists below this cutoff, rather than singling out twist $401$.  The main exposition gives the quantifier order in Theorem~\ref{thm:quantified-window} and proves the $400$-to-$401$ rescaling in Lemma~\ref{lem:window-rescaling}.  This appendix identifies, but does not independently prove, the deep arithmetic estimates used to establish the window theorem itself.

\chapter{The final source interfaces and verification status}\label{app:final}

\section{Nine compatible requirements}\label{app:assembly-technical}

The finite-threshold assembly described in Section~\ref{sec:common-threshold} combines nine named \texttt{Closed} requirements on one admissible family. In the source, these are \texttt{LinkN} (local arithmetic/cohomological loss), \texttt{LinkEta} (local covering/twist compatibility), \texttt{LinkIRR} (exceptional row measure), \texttt{SmallLevelRankEq} (rank identification), \texttt{LinkDIVt} (conditional divisibility), \texttt{NuFullF} together with \texttt{J.MuRank} (full column lattices and row ranks), \texttt{LinkG} (finite-place gain and index control), \texttt{LinkS} (Gram, frame and rank comparisons), and \texttt{LinkGLOBAL} (the quantitative global comparison). The requirements are assembled in \texttt{DesWin.lean} and supplied by the larger \texttt{ShapeWin.lean} construction. Design existence, \texttt{HypV} and the enclosing \texttt{TopShape} hypotheses are \emph{additional} prerequisites, not secretly among the nine. The formal requirements are conditional where their stated premises demand it, and for each fixed sufficiently small $\eps$, one family is chosen after the corresponding $\eps$-dependent common threshold has been assembled; no single family uniform across all $\eps$ is asserted.

After these inputs are selected together, the arithmetic balance~\eqref{eq:balance-final} and finite Fourier extraction~\eqref{eq:fourier-error-budget} leave a normalized geometric error, a discarded-character bound, and stability errors.  In the source the final specialization uses $e_D=2$, $e_L=0$, and window ratio $3/2$, through \sourcefile{src/Development/AbelianVarietyTheory/Proofs/GlobalInterpolation/LinkGl.lean}{LinkGl.lean}, \sourcefile{src/Development/AutomorphicArithmetic/CharacterBoxes/ShapeWin.lean}{ShapeWin.lean} and \sourcefile{src/Development/AutomorphicArithmetic/CharacterBoxes/DesWin.lean}{DesWin.lean}; this is a checked source-level parameter choice, not an independent reproof of the underlying geometric lemmas. The theorem \texttt{zeroFreeAll\_of\_topG} packages the analytic conclusion. The short route via \texttt{EveryCharacter.plugAll} is one convenient interface, whereas the actual public-theorem route in \texttt{MainConclusion.lean} also passes through the literature-input and \sourcefile{src/Development/AutomorphicArithmetic/ThetaCorrespondence/ThetaPairLifts/ThirdPlaceSeries/TopPlug.lean}{TopPlug.lean} bridges before \texttt{MainTheorems.lean}. These are related source routes, not a claim that \texttt{MainConclusion.lean} literally imports \texttt{plugAll}.  The text has developed the available elementary algebra and analytic implications while citing, rather than independently reconstructing, the massive geometric and cohomological existence results.

\section{Independent verification of the formal archive}\label{app:verification}
\enlargethispage{2\baselineskip}

Three verification tasks are distinct: checking that the formal archive compiles and its proof objects are accepted by kernels; checking that its public theorem statements have the intended Riemann Hypothesis and Dirichlet-character meanings; and checking that the mathematical explanations in this exposition accurately represent the deep declarations used by the formal proof. Public third-party reports provide substantial evidence for the first two tasks. The third requires source-level mathematical review of the reconstruction.

Andrew Sutherland's report at \url{https://math.mit.edu/~drew/grh/}, updated 21 November 2026 at 12:55 UTC, records successful compilation of all 175,187 modules with Lean and Mathlib 4.33.1, module-by-module Lean kernel replay, acceptance of the complete exported proof by the independent \texttt{nanoda} and \texttt{con-leche} kernels, and acceptance of every compiled module by \texttt{lean4lean}. It also reports an accepted completed \texttt{con-leche --compare} check against the separately prepared five-theorem challenge. The distinct end-to-end \texttt{lake comparator} job, and full-export runs through the default and paranoid Lean kernels, were still in progress at that report's timestamp; the per-module \texttt{leanchecker} and \texttt{lean4lean} checks had completed. The public statements use Mathlib's definitions of zeta and Dirichlet $L$-functions and depend only on Lean's standard axioms. These are external formal checks of the published proof archive identified in that report, rather than checks performed by this exposition. The public \texttt{CocoSgt} source mirror used here is pinned at a separate revision; the correspondence of that snapshot with the tested archive has not been established by byte comparison.

Neither an accepted Lean proof object nor a readable informal derivation automatically proves that every intermediate explanation offered here is the optimal or even the most natural mathematical description of the source.  In particular, the construction-specific normalizations of periods and integral lattices merit continued human review.  The unattributed 23-page paper \emph{Dirichlet $L$-functions and the Critical Line}, shared in November 2026 at \url{https://ulam.ai/research/rh-short.pdf}, was useful for identifying the missing mathematical interfaces in an earlier edition of this book.  Authorship has not been established from the copy examined.  It is a secondary exposition, not a substitute for the pinned Lean declarations.  Wherever this edition records a deep result as a cited input, the distinction remains explicit.

\backmatter
\printbibliography[heading=bibintoc,title={References}]
\chapter*{Sources and verification notes}
\addcontentsline{toc}{chapter}{Sources and verification notes}
\begin{enumerate}\raggedright
\item L. Alp\"oge, \emph{Remark on the announced Riemann Hypothesis argument}, November 2026, \url{https://alpo.ge/remark.pdf}. Pages 1, 3 and 4 are cited at the mathematical decisions they illuminate: Stepanov multiplicity, the ramp sum, seesaw periods, the exponent loss, and the cohomological and Iwasawa components. This personal remark is a preliminary account, not a written proof.
\item A. Sutherland, \emph{Verification files for the Lean proof of the Riemann Hypothesis}, MIT, updated 21 November 2026, \url{https://math.mit.edu/~drew/grh/}. External report of compilation, kernel replay, independent kernel acceptance and theorem-statement comparison; individual jobs still running at the reported timestamp.
\item \emph{Dirichlet $L$-functions and the Critical Line}, unattributed 23-page preliminary mathematical exposition circulated via \url{https://x.com/prz_chojecki/status/2124066186522431812} and available at \url{https://ulam.ai/research/rh-short.pdf}. Used as a technical comparison reference for identifying missing analytic, Fourier, frame and cancellation interfaces; authorship was not established from the inspected copy.
\item \emph{Riemann Hypothesis in Lean --- public mirror}, \newline source revision \texttt{7cbf7a0c32fc}\ldots, \url{https://github.com/CocoSgt/riemann-hypothesis-lean}. The mirror makes the Lean source available at a fixed revision. Sutherland's verification applies to the published archive identified in his report, and a byte-for-byte match with this mirror revision has not been established here; see Appendix~\ref{app:verification}.
\item \emph{The proof of the Riemann hypothesis announced by L. Alp\"oge: a guide to its architecture}, provisional independent explanatory guide, \url{https://nasqret.github.io/riemann-rev/guide/proof-guide.pdf}. Consulted for the announced proof's analytic statistic and the architecture of its arithmetic estimates. It is not itself the formal proof and has not been independently audited here.
\item \sourcefile{src/Statements.lean}{Statements.lean}, \sourcefile{src/Main.lean}{Main.lean}, \sourcefile{src/Supplements/Kernels/MainTheoremsAll.lean}{MainTheoremsAll.lean}, \sourcefile{src/Supplements/Kernels/MainTheoremsCritical.lean}{MainTheoremsCritical.lean}, and \sourcefile{src/Supplements/Kernels/MainTheorems.lean}{MainTheorems.lean}: the target statements, intermediate imports, and final theorem assembly. The conceptual zero-free route and the actual \texttt{MainConclusion} producer chain should not be conflated.
\item \sourcefile{src/Development/AutomorphicArithmetic/CharacterBoxes/EveryCharacter.lean}{EveryCharacter.lean}, \sourcefile{src/Development/AutomorphicArithmetic/ThetaCorrespondence/ThetaPairLifts/ThirdPlaceSeries/TopPlug.lean}{TopPlug.lean}, \sourcefile{src/Development/AutomorphicArithmetic/ThetaCorrespondence/ThetaPairLifts/ThirdPlaceSeries/EveryCharacterOfLiteratureOfVWin.lean}{EveryCharacterOfLiteratureOfVWin.lean}, \sourcefile{src/Development/AutomorphicArithmetic/CharacterBoxes/ClaimStrip.lean}{ClaimStrip.lean}, \sourcefile{src/Development/AutomorphicArithmetic/CharacterBoxes/Imprimitive.lean}{Imprimitive.lean}, and \sourcefile{src/Formalism/GammaFactorNonvanishing/CritLine.lean}{CritLine.lean}: alternate nonvanishing bridges, reflection, finite Euler factors, and the critical-line reduction.
\item \sourcefile{src/Development/AutomorphicArithmetic/Foundations/GaloisRepresentations/ThetaParseval/SplitLineSubtraction/RowDefOfOdd.lean}{RowDefOfOdd.lean} and \sourcefile{src/Development/AutomorphicArithmetic/Foundations/GaloisRepresentations/ThetaParseval/SplitLineSubtraction/SubtractOf.lean}{SubtractOf.lean}: the two-enlargement and subtraction propositions.
\item \sourcefile{src/Development/AutomorphicArithmetic/DedekindLFunctions/LandauTheorems/ZeroFreeOfRpowWindow.lean}{ZeroFreeOfRpowWindow.lean}: the sparse-window zero-free criterion described in Chapter~\ref{ch:statistic}.
\item \sourcefile{src/Formalism/LatticeSlopes/LatticeSlopesProofs/SlopeInequality.lean}{SlopeInequality.lean}: a formalized slope-comparison principle.
\item \sourcefile{src/Development/AutomorphicArithmetic/CharacterBoxes/ExistsPageN.lean}{ExistsPageN.lean} and \sourcefile{src/Formalism/GaloisCohomology/Statement.lean}{Statement.lean}: the definitions of the cohomological window and the underlying lengths, reserved for later chapters.
\end{enumerate}
\noindent\textit{Verification boundary.} This manuscript uses pinned source declarations and reconstructs selected elementary and conditional arguments. Independent formal checks reported by Sutherland support kernel acceptance of the original proof objects and comparison of the final statements. This exposition has not repeated those checks or independently reconstructed all deep arithmetic inputs. Its scope remains a readable mathematical reconstruction rather than a replacement for the formal archive.
\end{document}
